\chapter{Finite-energy functions and measures}
\index{finite-energy potential!BJ}\index{finite-energy measure!BJ}\index{Green function!BJ}
\label{chap:bj-finite-energy}

Energy turns the Monge--Amp\`ere equation into a variational problem.  The
potential side is the class \(\mathcal E^1\) introduced in the preceding
chapter.  The measure side is its Legendre dual.  Uniform differentiability
shows that an almost-maximizer has the desired Monge--Amp\`ere measure;
the envelope property supplies compactness and hence an actual maximizer.

That sentence is the whole chapter in four lines, and every one of the four
hides a hypothesis.  ``Variational problem'' presupposes a functional that
is finite somewhere, which is why \Cref{def:bj-e1} had to be made before
anything could be maximized.  ``Legendre dual'' presupposes a pairing
between potentials and measures that does not produce
\((+\infty)-(+\infty)\), which is why \Cref{prop:ch46-no-mass-outside-lin}
had to be proved first.  ``Uniform differentiability'' presupposes that the
obstacle envelope \(P_\omega\) is an operator with values in
\(\PSH(\omega)\), which is \Cref{imp:ch45-obstacle-operator} and is
available only after \Cref{chap:bj-envelope-consequences}.  And
``compactness'' is \Cref{cor:ch45-compactness}, which before that chapter
was a conjecture.  The purpose of the orientation below is to make the four
hypotheses visible before they are used, and the purpose of
\Cref{sec:ch47-imported} is to list, once, the four statements that are
imported rather than proved.

\section{Orientation}
\label{sec:ch47-orientation}

\subsection{What is proved, and out of what}
\label{sec:ch47-what}

Two objects are constructed here and one theorem is proved about them.

The first object is a \emph{geometry} on the class \(\cE^1(\omega)\).  It
is not a norm, because \(\cE^1(\omega)\) is convex but is not a vector
space (\Cref{rem:ch46-vector-space}); it is not a metric either, because
the natural quantity \(I(\phi,\psi)\) satisfies the triangle inequality
only up to a dimensional constant.  What it is, is a translation-invariant
quasi-metric whose square-root behaves like a norm on differences and whose
sublevel sets are exactly the sets on which energy is bounded.  All of its
properties are consequences of the cocycle formula
\eqref{eq:bj-energy-cocycle} and of the positivity of the Dirichlet form
\Cref{cor:ch46-dirichlet-e1}; none of them requires any new analysis.  This
is \Cref{sec:ch47-geometry}.

The second object is the \emph{Legendre dual} \(E^\vee_\omega\) of the
energy, a convex, weakly lower semicontinuous functional on Radon
probability measures whose finiteness locus \(\cM^1(\omega)\) is the
measure-side counterpart of \(\cE^1(\omega)\).  Duality is elementary once
the tangent inequalities \eqref{eq:bj-energy-tangent-inequalities} are
available; what is not elementary is that the two sides correspond, and
that correspondence is the theorem.  This is \Cref{sec:ch47-dual}.

The theorem is the \emph{non-Archimedean Calabi--Yau theorem}: under the
envelope property, \(\MA_\omega:\cE^1(\omega)\to\cM^1(\omega)\) is
surjective, and injective modulo constants.  Its proof is variational and
has exactly four ingredients, which it is worth naming in advance because
each is supplied by a different earlier chapter:

\begin{enumerate}
\item\label{itm:ch47-ingredient:usc}
      \emph{Upper semicontinuity and concavity of \(E_\omega\)}, which make
      the functional \(F_\mu=E_\omega-\int\cdot\,\rd\mu\) upper
      semicontinuous and concave.  This is
      \Cref{prop:ch46-energy-usc}, and it is intersection theory plus a
      monotone limit.
\item\label{itm:ch47-ingredient:compact}
      \emph{Compactness of \(\PSH^{\sup}(\omega)\)}, which produces a
      maximizer.  This is \Cref{cor:ch45-compactness}, and it is a
      consequence of the envelope theorem.
\item\label{itm:ch47-ingredient:diff}
      \emph{Uniform differentiability of the equilibrium energy}, which
      identifies the Monge--Amp\`ere measure of an almost-maximizer.  This
      is \Cref{thm:bj-energy-differentiability}, and it is the one place
      where a quantitative estimate from \cite{BJ22} is imported rather
      than proved.
\item\label{itm:ch47-ingredient:concave}
      \emph{Uniform concavity of \(E_\omega\)}, which makes every
      maximizing net Cauchy for \(I\) and hence yields uniqueness and, in
      combination with completeness, existence for an arbitrary
      \(\mu\in\cM^1(\omega)\).  This is
      \Cref{thm:bj-uniform-concavity}, and it is proved here in full, with
      an explicit constant.
\end{enumerate}

Ingredient \cref{itm:ch47-ingredient:compact} is the reason the chapter
could not have been written before
\Cref{chap:bj-envelope-consequences}.  Ingredient
\cref{itm:ch47-ingredient:diff} is the reason it cannot be written
entirely from the material of \Cref{chap:bj-energy} either.

\subsection{The five tasks}
\label{sec:ch47-tasks}

\begin{enumerate}
\item\label{itm:ch47-task:geometry}
      \emph{\(\cE^1\) and its metric geometry.}  Define \(J_\psi\), \(I\)
      and \(J\); prove the comparison
      \eqref{eq:bj-I-J-comparison} with the sharp binomial constants,
      the symmetry and vanishing properties, uniform concavity
      \eqref{eq:bj-uniform-concavity} with an explicit \(c_n\), and the
      behaviour under the scaling action.  Record the quasi-metric
      \(\mathbb I\) of \eqref{eq:bj-I-normalized}, say what belongs here
      and what is deferred to \Cref{chap:bj-strong-topologies}.
      \Cref{sec:ch47-geometry}.
\item\label{itm:ch47-task:dual}
      \emph{Finite-energy measures and dual energy.}  Define
      \(E^\vee_\omega\) by Legendre duality; prove that the test class may
      be taken to be any of four natural families; prove convexity, lower
      semicontinuity, compactness of sublevels, the Fenchel identities, the
      characterization of \(E^\vee=0\), homogeneity for the scaling action,
      and the vanishing of \(\mu\) on the polar set of a finite-energy
      potential.  \Cref{sec:ch47-dual}.
\item\label{itm:ch47-task:diff}
      \emph{Differentiability of energy.}  Prove the exact one-sided
      tangent inequality for the equilibrium energy, the polynomial
      expansion along segments in \(\cE^1\), and deduce
      \eqref{eq:bj-uniform-differentiability} and
      \eqref{eq:bj-first-variation} from the one imported quadratic
      estimate.  \Cref{sec:ch47-differentiability}.
\item\label{itm:ch47-task:orthogonality}
      \emph{Orthogonality and the variational equation.}  Prove that
      envelopes are in contact with their obstacles almost everywhere for
      the Monge--Amp\`ere measure, that maximizing nets solve the equation
      in the weak limit, and that a maximizer is a solution; record the
      domination principle and uniqueness.
      \Cref{sec:ch47-orthogonality}.
\item\label{itm:ch47-task:cy}
      \emph{The non-Archimedean Calabi--Yau theorem.}  Prove coercivity and
      existence by the direct method for an explicit class of measures;
      prove that maximizing nets are Cauchy; assemble the equivalence
      \Cref{thm:bj-cy-envelope-equivalence} and the unconditional
      \Cref{cor:bj-cy-unibranch}.  \Cref{sec:ch47-cy}.
\end{enumerate}

\subsection{Standing hypotheses and notation}
\label{sec:ch47-hypotheses}

We retain the hypotheses and probability normalization of
\Cref{chap:bj-energy}: \(X\) is integral and projective of dimension
\(n\), \(\omega\in\Amp(X)\), \(V=(\omega^n)\), and
\(\MA(\phi)=V^{-1}(\omega+\ddc\phi)^n\).  Write

\[
 \mathcal E^1:=\mathcal E^1(\omega),\qquad
 \mathcal M:=\{\hbox{Radon probability measures on }X^{\An}\}.
\tag{47.1}\label{eq:bj-e1-m-probability}
\]

All integrals of an \(\mathcal E^1\)-function are interpreted only after
integrability has been established.  Since these functions are real-valued
on \(X^{\lin}\) but may be \(-\infty\) at proper-support points, this is a
genuine safeguard.

Beyond this we assume, as in \Cref{chap:bj-energy}, that \(X\) is
\emph{unibranch}, so that \Cref{thm:bj-final-envelope} applies and the
obstacle envelope is an operator.  The ground field \(k\) is algebraically
closed and trivially valued.  As in \Cref{sec:ch46-hypotheses} we abbreviate
\(\PSH(\theta):=\PSH_X(\theta)\), \(\PSH^{\sup}(\theta):=\PSH^{\sup}_X(\theta)\)
for the normalized space \eqref{eq:normalized-psh-space}, and
\(\PL^\theta(X):=\PL_{\mathbf R}(X)\cap\PSH(\theta)\).  We write

\[
 \begin{aligned}
 \cE^1_0(\omega)
 &:=\{\phi\in\cE^1(\omega)
      \mid\textstyle\sup_{X^{\An}}\phi=0\},\\
 \PSH^{\sup}(\omega)
 &:=\{\phi\in\PSH(\omega)
      \mid\textstyle\sup_{X^{\An}}\phi=0\}.
 \end{aligned}
\tag{47.2}\label{eq:ch47-normalized-classes}
\]

so that \(\cE^1_0(\omega)=\cE^1(\omega)\cap\PSH^{\sup}(\omega)\), and every
\(\phi\in\cE^1(\omega)\) is uniquely \(\phi_0+c\) with
\(\phi_0\in\cE^1_0(\omega)\) and \(c=\sup\phi\in\mathbf R\).

Three facts about the trivially valued setting are used so often that they
are recorded here rather than rediscovered each time.

\begin{lemma}
\label{lem:ch47-three-normalizations}
Let \(\omega\in\Amp(X)\).  Then:
\begin{enumerate}
\item\label{itm:ch47-norm:trivial-measure}
      \(\MA_\omega(0)=\delta_{v_{\triv}}\), where \(v_{\triv}\) is the
      trivial valuation of \(k(X)\);
\item\label{itm:ch47-norm:max}
      \(\sup_{X^{\An}}\phi=\phi(v_{\triv})\) for every
      \(\phi\in\PSH(\omega)\);
\item\label{itm:ch47-norm:J}
      consequently, for \(\phi\in\cE^1(\omega)\) the two natural
      normalizing functionals coincide:
      \[
       J(\phi):=\sup_{X^{\An}}\phi-E_\omega(\phi)
       =E_\omega(0)-E_\omega(\phi)+\int(\phi-0)\,\MA_\omega(0)
       =J_0(\phi),
      \]
      with \(J_0\) as in \eqref{eq:bj-I-J-potentials}.  In particular
      \(J\ge0\), \(J(\phi+c)=J(\phi)\), and \(J(\phi)=0\) if and only if
      \(\phi\) is constant.
\end{enumerate}
\end{lemma}

\begin{proof}
\cref{itm:ch47-norm:trivial-measure}.  Apply
\Cref{prop:ch46-pl-measure} on the trivial test configuration
\(\cX=X\times\mathbf A^1\) with \(D=0\).  Its central fibre is
\(\cX_0=X\times\{0\}\), irreducible and reduced, so \(b_{\cX_0}=1\), and
the associated valuation is \(v_{\cX_0}=v_{X,\triv}\) by
\Cref{def:ch33-component-valuation}.  The single coefficient is
\((\omega_{\cX}^{\,n}|_{\cX_0})=(\omega^n)=V\), so
\(\MA^{\mathrm u}_\omega(0)=V\delta_{v_{\triv}}\) and, dividing by \(V\),
\(\MA_\omega(0)=\delta_{v_{\triv}}\).  This computation was already used in
the proof of \Cref{prop:ch46-energy-usc}\,\cref{itm:ch46-Eusc:strict}.

\cref{itm:ch47-norm:max}.  For a Fubini--Study function this is the maximum
principle \Cref{prop:ch34-generic-finiteness}\,\cref{itm:ch34-gf:decreasing}:
in a presentation \(\phi=m^{-1}\max_j(\log|s_j|+\lambda_j)\) each term is
decreasing along every ray emanating from \(v_{\triv}\), because
\(\log|s_j|(tv)=-t\,v(s_j)\le0\) for \(t\ge0\), so the maximum is attained
at \(t=0\).  For general \(\phi\in\PSH(\omega)\), take a decreasing net
\((\phi_i)\) of Fubini--Study functions with \(\phi_i\downarrow\phi\)
(\Cref{def:bj-psh}\,(3)).  Then
\[
 \sup_{X^{\An}}\phi=\sup_{x}\inf_i\phi_i(x)
 \ \le\ \inf_i\sup_x\phi_i(x)
 =\inf_i\phi_i(v_{\triv})=\phi(v_{\triv})
 \ \le\ \sup_{X^{\An}}\phi ,
\]
so all the terms are equal.  Only the elementary inequality
\(\sup\inf\le\inf\sup\) and the already-proved Fubini--Study case were
used.

\cref{itm:ch47-norm:J}.  Combine the two previous items: by
\cref{itm:ch47-norm:trivial-measure},
\(\int\phi\,\MA_\omega(0)=\phi(v_{\triv})\), and by
\cref{itm:ch47-norm:max} this equals \(\sup\phi\); and
\(E_\omega(0)=0\) by \eqref{eq:bj-trivial-energy-zero}.  Hence
\(J_0(\phi)=-E_\omega(\phi)+\sup\phi=J(\phi)\).
Nonnegativity is \Cref{prop:ch46-energy-usc}\,\cref{itm:ch46-Eusc:identities}
(\(E_\omega(\phi)\le\sup\phi\)); translation invariance follows from
\(E_\omega(\phi+c)=E_\omega(\phi)+c\) and
\(\sup(\phi+c)=\sup\phi+c\).  If \(J(\phi)=0\) then
\(E_\omega(\phi)=\sup\phi=E_\omega(\sup\phi)\) with
\(\phi\le\sup\phi\), so \(\phi=\sup\phi\) almost everywhere for
\(\MA_\omega(\sup\phi)=\delta_{v_{\triv}}\) by
\Cref{prop:ch46-cocycle}\,\cref{itm:ch46-coc:mono}; this gives only
\(\phi(v_{\triv})=\sup\phi\), which is automatic.  The genuine conclusion
that \(\phi\) is constant is \Cref{cor:ch47-J-vanishes}, which uses the
\(I\)-\(J\) comparison and is proved in \Cref{sec:ch47-geometry}.
\end{proof}

\begin{warning}{Item \cref{itm:ch47-norm:J} has no counterpart in \Cref{part:bfj}}
\label{warn:ch47-J-identity}
Over a discretely valued base the reference measure \(\MA_\omega(0)\) is
the Chambert--Loir measure of the trivial model metric, which is a sum of
several atoms with weights \(b_E(\cL^{\,n}|_E)\), and there is no point of
\(X^{\An}\) at which every psh function attains its maximum.  Accordingly
\Cref{def:ch26-ij} keeps \(I\) and \(J\) apart from the sup-normalization,
and \Cref{prop:ch26-energy-lipschitz} carries the correction term
explicitly.  The identity
\(J(\phi)=\sup\phi-E_\omega(\phi)=J_0(\phi)\) of
\Cref{lem:ch47-three-normalizations}\,\cref{itm:ch47-norm:J} is a feature of
the trivially valued theory, produced by the cone structure of
\(X^{\An}\): the cone point \(v_{\triv}\) carries all of
\(\MA_\omega(0)\) and is where every psh function is largest.  It is used
below whenever a normalized supremum is traded for an energy, and any
attempt to transplant such a step to \Cref{part:bfj} must reinstate the
correction term.
\end{warning}

\subsection{The normalization, inherited}
\label{sec:ch47-normalization}

\Cref{chap:bj-energy} declared a departure from the global convention of
the front matter, in \Cref{warn:ch46-normalization}: the unadorned symbol
\(\MA_\omega\) denotes the \emph{probability} measure
\(V^{-1}(\omega+\ddc\phi)^n\), the unnormalized measure of mass \(V\)
carries a superscript \(\mathrm u\), and consequently
\(E_\omega(\phi+c)=E_\omega(\phi)+c\).  That declaration is in force for
the whole of \Crefrange{chap:bj-energy}{chap:bj-stability} and is
\emph{not} repeated or amended here.  Two consequences should be kept in
view, because both are used constantly below.

First, \(\mathcal M\) in \eqref{eq:bj-e1-m-probability} consists of
probability measures, and the Legendre duality of
\Cref{sec:ch47-dual} is a duality between \(\cE^1(\omega)\), on which the
additive group \(\mathbf R\) acts by translation, and \(\mathcal M\), on
which that action is dual to nothing at all --- which is exactly why
\(E^\vee_\omega\) is finite only on probability measures
(\Cref{lem:ch47-mass-forced}).

Second, every numerical comparison with \Cref{part:bfj} carries a factor
\(V\).  \Cref{def:dual-energy-measure} there is stated for measures of
total mass \(V\) and its Fenchel inequality
\eqref{eq:measure-fenchel-energy} is the unnormalized one; the dictionary
\eqref{eq:ch28-probability-dictionary} records the conversion.  Every
displayed comparison below applies it explicitly.

\subsection{Extended-real ledger}
\label{sec:ch47-ledger}

\Cref{sec:ch46-ledger} listed three places where a value \(-\infty\) is
manipulated in the construction of the energy.  Two more appear here, and
they are the only new ones.

\begin{enumerate}
\item\label{itm:ch47-ledger:integral}
      For \(\phi\in\PSH(\omega)\) and \(\mu\in\mathcal M\), the integral
      \(\int\phi\,\rd\mu\) is defined in \([-\infty,\sup\phi]\) by
      monotone truncation: \(\int\phi\,\rd\mu:=\lim_j\int\max\{\phi,-j\}\rd\mu\),
      the limit of a decreasing sequence of real numbers.  No convention is
      needed for the value \(+\infty\), which does not occur, because psh
      functions are bounded above.  The Fenchel expression
      \(E_\omega(\phi)-\int\phi\,\rd\mu\) is therefore well defined in
      \([-\infty,+\infty]\) except when \(E_\omega(\phi)=-\infty\) and
      \(\int\phi\,\rd\mu=-\infty\) simultaneously; this is why
      \eqref{eq:bj-dual-energy} takes its supremum over \(\cE^1\), where
      the first term is finite, and why
      \Cref{prop:ch47-test-class} has to be proved rather than assumed.
\item\label{itm:ch47-ledger:sup}
      A supremum of the form \(\sup_i(a_i-b_i)\) with \(a_i\) finite and
      \(b_i\in[-\infty,+\infty)\) is unambiguous, and equals \(+\infty\)
      as soon as some \(b_i=-\infty\).  This is the mechanism by which
      \(E^\vee_\omega(\mu)=+\infty\) detects that \(\mu\) charges the polar
      set of a finite-energy potential
      (\Cref{prop:ch47-no-polar-mass}).
\end{enumerate}

The prohibition on \((-\infty)-(-\infty)\) remains binding.  In particular
the difference \(\phi-\psi\) of two members of \(\cE^1(\omega)\) is formed
on \(X^{\lin}\), where both are finite by
\eqref{eq:bj-e1-sublinear-growth}, and every mixed measure occurring in
this chapter is carried there by
\Cref{prop:ch46-no-mass-outside-lin}.

\subsection{What is imported, and what is proved}
\label{sec:ch47-imported}

\begin{important}{The four imported statements}
\label{imp:ch47-imported}
Everything in this chapter is proved from
\Crefrange{chap:bj-fs-pl-psh}{chap:bj-energy} except for the following four
statements, each of which is quoted from \cite{BJ22} with an exact
reference and used only where indicated.
\begin{enumerate}
\item\label{itm:ch47-import:diff}
      \emph{The uniform quadratic lower bound for the equilibrium energy},
      \cite[Theorem~8.5]{BJ22}: for \(f\in\PL_{\mathbf R}(X)\) there is
      \(C=C(f)\), independent of \(\phi\), with
      \(E^\downarrow(\phi+\varepsilon f)\ge
      E_\omega(\phi)+\varepsilon\int f\,\MA_\omega(\phi)-C\varepsilon^2\)
      for all \(\phi\in\cE^1(\omega)\) and \(\varepsilon\in[0,1]\).
      The matching upper bound is proved here
      (\Cref{lem:ch47-one-sided-tangent}) and needs no hypothesis at all.
      Used in \Cref{thm:bj-energy-differentiability} and hence in
      \Cref{thm:bj-maximizers-converge} and everything after it.
\item\label{itm:ch47-import:uniqueness}
      \emph{Uniqueness for the Monge--Amp\`ere equation in \(\cE^1\)},
      \cite[Corollary~10.4]{BJ22}: if \(\phi,\psi\in\cE^1(\omega)\) and
      \(\MA_\omega(\phi)=\MA_\omega(\psi)\), then \(\phi-\psi\) is
      constant.  Proved here for \(n=1\)
      (\Cref{prop:ch47-curve-uniqueness}); used for the last clause of
      \Cref{prop:bj-energy-I-identities} and for the uniqueness half of
      \Cref{cor:bj-cy-unibranch}.
\item\label{itm:ch47-import:quasitriangle}
      \emph{The quasi-triangle inequality} \eqref{eq:bj-I-quasitriangle},
      \cite[Theorem~3.31]{BJ22}.  Proved here for \(n=1\) with the sharp
      constant \(C_1=2\) (\Cref{prop:ch47-curve-I-norm}); used for the
      statement that \(\mathbb I\) is a quasi-metric, and once inside a
      proof, namely for the energy bound along an \(I\)-Cauchy net in the
      proof of \Cref{thm:bj-cy-envelope-equivalence}.
\item\label{itm:ch47-import:completeness}
      \emph{Continuity of energy along increasing nets and completeness of
      the strong space}, \cite[Theorems~12.4 and~12.5]{BJ22}.  Used only in
      the proof of \Cref{thm:bj-cy-envelope-equivalence}, for the passage
      from an \(I\)-Cauchy maximizing net to an actual maximizer, and
      developed in \Cref{chap:bj-strong-topologies}.  For the class of
      measures treated in \Cref{thm:ch47-existence-admissible} it is not
      needed: there the direct method applies and the proof is complete.
\end{enumerate}
No other external input is used.  In particular the capacity theory of
\Cref{chap:bj-strong-topologies} is \emph{not} used here; the logical
direction is the reverse, as the source note of that chapter records.
\end{important}

The reason for listing the imports so prominently is that the four
statements are of very different depths.
Item~\cref{itm:ch47-import:quasitriangle} is a Hodge-type estimate that
belongs to the finitary theory of \Cref{sec:ch46-model}; withholding it
would cost \(\mathbb I\) its quasi-metric status --- it would be merely a
symmetric separating gauge --- and would remove the uniform energy bound
along an \(I\)-Cauchy net in the proof of
\Cref{thm:bj-cy-envelope-equivalence}, which is where the compactness
argument there starts.
Item~\cref{itm:ch47-import:uniqueness} is a genuine theorem whose failure
would leave the Calabi--Yau theorem an existence statement without
uniqueness.  Item~\cref{itm:ch47-import:diff} is the analytic heart:
without it, \Cref{thm:bj-maximizers-converge} is false as stated, and the
variational method produces maximizers that need not solve any equation.
Item~\cref{itm:ch47-import:completeness} is a topological statement about
\(\cE^1\) that is developed in the next chapter and is used here only to
remove the admissibility hypothesis from the existence theorem.

\begin{important}{The envelope theorem is used without further comment}
\label{imp:ch47-envelope}
As in \Cref{imp:ch46-envelope-imported}, \(X\) is unibranch and
\Cref{thm:bj-final-envelope} is available.  The three consequences used
here are:
\begin{enumerate}
\item\label{itm:ch47-env:operator}
      \(P_\omega\) is an operator from admissible obstacles
      (\Cref{def:ch45-admissible-obstacle}) to \(\PSH(\omega)\), and is the
      identity on \(\PSH(\omega)\)
      (\Cref{imp:ch45-obstacle-operator},
      \Cref{prop:ch45-obstacle-calculus-2}\,\cref{itm:ch45-P:projection});
\item\label{itm:ch47-env:compact}
      \(\PSH^{\sup}(\omega)\) is compact for the topology of pointwise
      convergence on \(X^{\divis}\) (\Cref{cor:ch45-compactness});
\item\label{itm:ch47-env:calculus}
      the calculus of \(P_\omega\)
      (\Cref{prop:bj-final-obstacle-calculus},
      \Cref{prop:ch45-obstacle-calculus-2}) and the monotone statements
      \Cref{thm:ch45-decreasing-obstacle-nets,thm:ch45-increasing-obstacles}.
\end{enumerate}
\end{important}

\subsection{Relation to \texorpdfstring{\Cref{part:bfj}}{Part III}}
\label{sec:ch47-part3}

This chapter is the trivially valued counterpart of
\Cref{chap:bfj-energy,chap:bfj-differentiability,chap:bfj-finite-energy-measures},
and the parallel is close enough that the differences must be tabulated
before they are met.  The rule of thumb is: \emph{every argument that uses
only multilinearity, the tangent inequalities and the Dirichlet form has
the same algebraic scheme; every argument that uses a reference measure of full
support, a model, or a dual complex does not; and every statement involving
the scaling action is new.}  Here ``transfers'' concerns the algebraic
proof pattern, not an identification of ambient pairings: in the
discretely valued column every energy term retains an actual vertical
factor, whereas in the trivially valued column pair powers are computed on
\(\mathbf P^1\)-compactified test configurations.

\begin{center}
\begin{tabular}{p{.24\textwidth}p{.35\textwidth}p{.32\textwidth}}
\toprule
& \Cref{part:bfj} (discretely valued) & this chapter (trivially valued)\\
\midrule
\(I\) and \(J\) &
\Cref{def:ch26-ij}, \eqref{eq:ch26-I-definition}--\eqref{eq:ch26-J-definition};
unnormalized &
\eqref{eq:bj-I-J-potentials}; probability normalized
(\Cref{warn:ch46-normalization})\\
comparison \(I\asymp J\) &
\Cref{thm:ch26-mixed-energy-identities}\,\cref{itm:ch26-ij:comparison} &
\Cref{prop:bj-energy-I-identities}; \emph{same proof}, the binomial
coefficients of \eqref{eq:ch47-bernstein-derivative}\\
sup-normalization &
\(J\) and \(\sup\varphi-E\) differ
(\Cref{prop:ch26-energy-lipschitz}) &
they coincide
(\Cref{lem:ch47-three-normalizations}\,\cref{itm:ch47-norm:J},
\Cref{warn:ch47-J-identity})\\
dual energy &
\Cref{def:dual-energy-measure}, mass \(V\) &
\Cref{def:bj-finite-energy-measure}, mass \(1\)\\
existence of a maximizer &
\Cref{thm:ch28-existence-admissible}, for \emph{admissible} \(\mu\) &
\Cref{thm:ch47-existence-admissible} for admissible \(\mu\);
\Cref{thm:bj-cy-envelope-equivalence} for all
\(\mu\in\cM^1\)\\
what makes coercivity work &
\(M(\mu)<\infty\), \Cref{lem:ch28-coercivity} &
\(M(\mu)\le\int T_\omega\,\rd\mu\),
\Cref{lem:ch47-coercivity}; \emph{same proof}\\
orthogonality &
a hypothesis on \(\omega\) (\Cref{def:orthogonality-property}), proved in
special cases &
a theorem (\Cref{thm:bj-obstacle-orthogonality-energy}), because the
envelope property holds\\
surjectivity of \(\MA\) &
\Cref{prop:ch28-image-is-solvable}: the image is a proper subset in
general &
\Cref{thm:bj-cy-envelope-equivalence}: onto \(\cM^1\)\\
extra structure &
none &
the scaling action: \(I\), \(J\) and \(E^\vee\) are homogeneous of degree
one (\Cref{prop:ch47-scaling})\\
\bottomrule
\end{tabular}
\end{center}

Three rows deserve a comment.

\emph{Where the algebraic steps agree.}  The comparison
\eqref{eq:bj-I-J-comparison}, the coercivity estimate of
\Cref{lem:ch47-coercivity}, the identification of the value of the
variational problem with the dual energy, and the deduction of the equation
from maximality follow the arguments of
\Cref{thm:ch26-mixed-energy-identities},
\Cref{lem:ch28-coercivity},
\Cref{prop:ch28-sup-is-dual-energy} and
\Cref{prop:variational-maximizer-solves}.  They use only the cocycle
formula, the tangent inequalities and the positivity of mixed measures, all
of which hold in both regimes.  Where a proof below says ``as in
\Cref{part:bfj}'', this is what is meant, and the corresponding statement is
cited so that the reader can check the transfer.

\emph{Where they are not.}  Three mechanisms available over a discretely
valued base are unavailable here.  There is no reference measure of full
support: \(\MA_\omega(0)=\delta_{v_{\triv}}\) by
\Cref{lem:ch47-three-normalizations}\,\cref{itm:ch47-norm:trivial-measure},
so the classical proof of the domination principle --- comparison against a
nondegenerate measure --- collapses, which is why
\Cref{thm:ch47-domination} is imported and not proved
(\Cref{rem:ch47-no-nondegenerate-measure}).  There is no cofinal system of
models with dual complexes carrying the whole space, so the reductions of
\Cref{prop:fixed-skeleton-stability} and
\Cref{lem:ch28-skeleton-admissible} have only the partial analogue
\Cref{prop:ch47-skeletal-admissible}.  And sequences are never substituted
for nets: \Cref{warn:ch45-nets} is in force throughout, whereas
\Cref{rem:ch22-legitimate-sequences} permits some sequential reductions in
\Cref{part:bfj}.

\emph{What is genuinely new.}  The scaling action makes \(\cE^1(\omega)\)
and \(\cM^1(\omega)\) cones and makes \(I\), \(J\) and \(E^\vee_\omega\)
homogeneous of degree one (\Cref{prop:ch47-scaling}).  The Calabi--Yau
equation is equivariant: \(\MA_\omega(\phi)=\mu\) implies
\(\MA_\omega(t\cdot\phi)=t_\star\mu\).  Nothing of this kind exists in
\Cref{chap:bfj-finite-energy-measures}, and
\Cref{cor:ch47-scaling-cy} is the first place in the book where a
Monge--Amp\`ere equation is solved on a whole orbit at once.

\subsection{Roadmap}
\label{sec:ch47-roadmap}

\Cref{sec:ch47-geometry} builds the geometry of \(\cE^1\): the functionals
\(J_\psi\), \(I\), \(J\), the polynomial expansion along segments, the
comparison \eqref{eq:bj-I-J-comparison}, uniform concavity, and the
behaviour under scaling.  It uses nothing beyond
\Cref{chap:bj-energy}.

\Cref{sec:ch47-dual} constructs the dual energy and the class
\(\cM^1(\omega)\), proves the four descriptions of the test class, the
Fenchel identities, and the elementary structure theory of
\(\cM^1(\omega)\).

\Cref{sec:ch47-differentiability} treats differentiability.  The exact
one-sided inequality \Cref{lem:ch47-one-sided-tangent} is proved first,
because it is used four times afterwards and requires no hypotheses; the
matching lower bound is the single imported estimate.

\Cref{sec:ch47-orthogonality} proves orthogonality, deduces that
maximizing nets solve the equation, and records the domination principle
and uniqueness.

\Cref{sec:ch47-cy} proves the Calabi--Yau theorem: first for an explicit
admissible class of measures by the direct method, with every compactness
input attributed; then in general, by showing that maximizing nets are
\(I\)-Cauchy and appealing to completeness.

\Cref{sec:ch47-examples} carries four examples to explicit numbers, and
\Crefrange{sec:ch47-scope}{sec:ch47-source} record scope, the chapter
dictionary, and provenance.

\section{The \texorpdfstring{\(I\)}{I}-geometry of potentials}
\label{sec:ch47-geometry}

The class \(\cE^1(\omega)\) carries no norm, because it is not a vector
space, and no obviously canonical metric, because the energy is not
quadratic.  What it does carry is a family of nonnegative functionals
manufactured out of the cocycle formula, and these turn out to control each
other to within dimensional constants.  This section constructs them and
proves the four estimates that are used for the rest of the chapter: the
comparison \eqref{eq:bj-I-J-comparison}, uniform concavity
\eqref{eq:bj-uniform-concavity}, the monotone estimate
\Cref{cor:ch47-monotone-I}, and homogeneity
\Cref{prop:ch47-scaling}.  Everything rests on one computation, the
expansion of \(t\mapsto E_\omega((1-t)\phi+t\psi)\) in the Bernstein basis,
and it is worth doing that computation once and carefully.

\subsection{The functionals \texorpdfstring{\(J_\psi\)}{J psi},
\texorpdfstring{\(I\)}{I} and \texorpdfstring{\(J\)}{J}}
\label{sec:ch47-three-functionals}

For \(\phi,\psi\in\mathcal E^1\), put

\[
 \begin{split}
 J_\psi(\phi)&:=E(\psi)-E(\phi)
                  +\int(\phi-\psi)\MA(\psi),\\
 I(\phi,\psi)&:=\int(\phi-\psi)
                    \bigl(\MA(\psi)-\MA(\phi)\bigr)
               =J_\psi(\phi)+J_\phi(\psi).
 \end{split}
\tag{47.3}\label{eq:bj-I-J-potentials}
\]

The tangent inequalities \eqref{eq:bj-energy-tangent-inequalities} imply
that both quantities are nonnegative.  They do not see simultaneous
translation of the two potentials.

Alongside them we use the \emph{normalized energy functional}

\[
 J(\phi):=\sup_{X^{\An}}\phi-E_\omega(\phi)
 \ \in[0,+\infty),
 \qquad\phi\in\cE^1(\omega),
\tag{47.4}\label{eq:ch47-J-functional}
\]

which by \Cref{lem:ch47-three-normalizations}\,\cref{itm:ch47-norm:J}
coincides with \(J_0(\phi)\).  Thus \(J\) measures how far \(\phi\) is from
being constant, and it is the functional whose sublevel sets are the
compact sets of the theory: by \eqref{eq:ch47-normalized-classes},
\[
 \{\phi\in\cE^1_0(\omega)\mid J(\phi)\le C\}
 =\{\phi\in\PSH^{\sup}(\omega)\mid E_\omega(\phi)\ge-C\},
\]
which is a closed subset of the compact space \(\PSH^{\sup}(\omega)\) by
\Cref{prop:ch46-energy-usc}\,\cref{itm:ch46-Eusc:props} and hence compact.

\begin{lemma}
\label{lem:ch47-IJ-elementary}
Let \(\phi,\psi\in\cE^1(\omega)\) and \(a,b\in\mathbf R\).  Then all the
integrals in \eqref{eq:bj-I-J-potentials} and
\eqref{eq:ch47-J-functional} are finite, and:
\begin{enumerate}
\item\label{itm:ch47-IJ:nonneg}
      \(J_\psi(\phi)\ge0\), \(I(\phi,\psi)\ge0\), \(J(\phi)\ge0\);
\item\label{itm:ch47-IJ:translation}
      \(J_{\psi+b}(\phi+a)=J_\psi(\phi)\) and
      \(I(\phi+a,\psi+a)=I(\phi,\psi)\), while \(J(\phi+a)=J(\phi)\);
\item\label{itm:ch47-IJ:symmetry}
      \(I(\phi,\psi)=I(\psi,\phi)\), and \(I(\phi,\phi)=0\);
\item\label{itm:ch47-IJ:sup}
      \(I(\phi,\psi)\le2\sup_{X^{\divis}}|\phi-\psi|\) and
      \(|E(\phi)-E(\psi)|\le\sup_{X^{\divis}}|\phi-\psi|\);
\item\label{itm:ch47-IJ:monotone}
      if \(\phi\le\psi\), then
      \(0\le J_\psi(\phi)\le E(\psi)-E(\phi)\le J_\phi(\psi)+E(\psi)-E(\phi)\),
      and \(I(\phi,\psi)\le\int(\psi-\phi)\MA(\phi)\).
\end{enumerate}
\end{lemma}

\begin{proof}
Finiteness of every term is the last clause of \Cref{thm:bj-e1-basic}
together with \Cref{prop:ch46-no-mass-outside-lin}: the difference
\(\phi-\psi\) is formed on \(X^{\lin}\), where both functions are finite,
and it is integrable for every mixed measure of finite-energy pairs.

\cref{itm:ch47-IJ:nonneg}.  The right-hand inequality of
\eqref{eq:bj-energy-tangent-inequalities}, applied with the roles of
\(\phi\) and \(\psi\) as displayed, reads
\(E(\phi)-E(\psi)\le\int(\phi-\psi)\MA(\psi)\), which is exactly
\(J_\psi(\phi)\ge0\).  Summing this with the same inequality for the
reversed pair gives \(I\ge0\).  For \(J\) see
\Cref{lem:ch47-three-normalizations}\,\cref{itm:ch47-norm:J}.

\cref{itm:ch47-IJ:translation}.  Immediate from
\(E(\phi+a)=E(\phi)+a\) (\eqref{eq:bj-energy-basic-identities}),
\(\MA(\psi+b)=\MA(\psi)\), and the fact that \(\MA(\psi)\) is a probability
measure.  Thus the energy translation and the constant part of the
integral cancel even when only one argument is translated.

\cref{itm:ch47-IJ:symmetry} is visible in the displayed formula for \(I\),
which is unchanged by interchanging \(\phi\) and \(\psi\).

\cref{itm:ch47-IJ:sup}.  Put
$C=\sup_{X^{\divis}}|\phi-\psi|$.  If $C=\infty$, the inequalities are
automatic.  Otherwise, the inequalities $\phi\leq\psi+C$ and
$\psi\leq\phi+C$ on $X^{\divis}$ extend to $X^{\An}$ by
\Cref{cor:ch35-divisorial-determination}.  Since $\phi$ and $\psi$ are
integrable against both probability measures, this bounds their
difference almost everywhere by $C$.  The integral defining $I$ is
therefore bounded by $2C$; the energy bound is
\Cref{prop:ch46-cocycle}\,\cref{itm:ch46-coc:lipschitz}.

\cref{itm:ch47-IJ:monotone}.  If \(\phi\le\psi\), then
\(J_\psi(\phi)=E(\psi)-E(\phi)-\int(\psi-\phi)\MA(\psi)\le E(\psi)-E(\phi)\)
because the subtracted integral is nonnegative; and
\(I=J_\psi(\phi)+J_\phi(\psi)\) with
\(J_\phi(\psi)=E(\phi)-E(\psi)+\int(\psi-\phi)\MA(\phi)\), so
\(I=\int(\psi-\phi)\MA(\phi)-\int(\psi-\phi)\MA(\psi)
\le\int(\psi-\phi)\MA(\phi)\).
\end{proof}

\subsection{The segment expansion}
\label{sec:ch47-segment}

The following theorem is the computational core of the section.  It is the
trivially valued analogue of \Cref{prop:ch27-energy-polynomial}: it uses
multilinearity of the extended pairing on \(\vec\cE^{\,1}(X)\) and the
Riesz identification of mixed measures.  The powers below belong to the
pair pairing of \Cref{chap:bj-energy}; at the PL level they are absolute
intersection powers on a compactified test configuration over
\(\mathbf P^1\).  They are not absolute \((n+1)\)-fold powers of closed
model forms over a discretely valued field, where only the vertical-factor
formalism of \Cref{chap:bfj-energy} is defined.

\begin{theorem}
\label{thm:ch47-segment-expansion}
Let \(\phi,\psi\in\cE^1(\omega)\) and put
\(\phi_t:=(1-t)\phi+t\psi\) for \(t\in[0,1]\).  Then
\(\phi_t\in\cE^1(\omega)\), and with

\[
 \Lambda_k:=\frac{(\omega,\phi)^{\,n+1-k}(\omega,\psi)^{\,k}}{(n+1)V},
 \qquad
 \lambda_k:=\int(\psi-\phi)\,
   V^{-1}(\omega+\ddc\phi)^{n-k}\wedge(\omega+\ddc\psi)^{k},
\tag{47.5}\label{eq:ch47-bernstein}
\]

for \(0\le k\le n+1\) and \(0\le k\le n\) respectively, one has
\(\Lambda_0=E(\phi)\), \(\Lambda_{n+1}=E(\psi)\),
\(\lambda_k=(n+1)(\Lambda_{k+1}-\Lambda_k)\), and

\[
 g(t):=E_\omega(\phi_t)
 =\sum_{k=0}^{n+1}\binom{n+1}{k}(1-t)^{n+1-k}t^k\Lambda_k,
 \qquad
 g'(t)=\sum_{k=0}^{n}\binom{n}{k}(1-t)^{n-k}t^k\lambda_k .
\tag{47.6}\label{eq:ch47-bernstein-derivative}
\]

In particular \(g\) is a real polynomial of degree at most \(n+1\),

\[
 \begin{aligned}
 g'(t)&=\int(\psi-\phi)\,\MA_\omega(\phi_t),\\
 g'(0)&=\int(\psi-\phi)\,\MA(\phi),&
 g'(1)&=\int(\psi-\phi)\,\MA(\psi).
 \end{aligned}
\]

and, writing \(f:=\psi-\phi\in\vec\cE^{\,1}(X)\),

\[
 \begin{aligned}
  e_k&:=\lambda_k-\lambda_{k+1}
    =V^{-1}\langle f,f\rangle_{\mathrm D,k}\ge0,\\
  g''(t)&=-n\sum_{k=0}^{n-1}\binom{n-1}{k}
    (1-t)^{n-1-k}t^ke_k\le0,
 \end{aligned}
\]

where \(\langle\,\cdot\,,\,\cdot\,\rangle_{\mathrm D,k}\) is the Dirichlet
form of \Cref{cor:ch46-dirichlet-e1} with the \(n-1\) auxiliary pairs
consisting of \(n-1-k\) copies of \((\omega,\phi)\) and \(k\) copies of
\((\omega,\psi)\).  Finally

\[
 \sum_{k=0}^{n-1}e_k=\lambda_0-\lambda_n=I(\phi,\psi),
 \qquad
 E(\psi)-E(\phi)=g(1)-g(0)=\frac1{n+1}\sum_{k=0}^{n}\lambda_k .
\]
\end{theorem}

\begin{proof}
\emph{Step 1: the segment lies in \(\cE^1\).}  \(\PSH(\omega)\) is convex
by \Cref{thm:bj-basic-closure}\,(3) and \(E_\omega\) is concave by
\Cref{prop:ch46-energy-usc}\,\cref{itm:ch46-Eusc:props}, so
\(E_\omega(\phi_t)\ge(1-t)E_\omega(\phi)+tE_\omega(\psi)>-\infty\); this is
also the convexity clause of \Cref{thm:bj-e1-basic}.

\emph{Step 2: the pair \((\omega,\phi_t)\) is a convex combination.}  In
the vector space of pairs \((\theta,u)\in N^1(X)_{\mathbf R}\times
\vec\cE^{\,1}(X)\) on which the extended pairing
\eqref{eq:bj-e1-polarized-pairing} is multilinear, one has the identity
\[
 (\omega,\phi_t)=(1-t)(\omega,\phi)+t(\omega,\psi),
\]
because \((1-t)\omega+t\omega=\omega\) and
\((1-t)\phi+t\psi=\phi_t\).  All three pairs are finite-energy pairs, so
every monomial in the expansion below is a finite real number, again by the
last clause of \Cref{thm:bj-e1-basic}; this is the only place where
finiteness of \emph{mixed} energies is needed, and it is what
\Cref{warn:ch46-no-subtraction} demands.

\emph{Step 3: the expansion.}  Expanding \((\omega,\phi_t)^{n+1}\) by the
binomial theorem for the symmetric multilinear pairing gives
\[
 (\omega,\phi_t)^{n+1}
 =\sum_{k=0}^{n+1}\binom{n+1}{k}(1-t)^{n+1-k}t^k
   (\omega,\phi)^{n+1-k}(\omega,\psi)^k,
\]
and dividing by \((n+1)V\) gives the first formula of
\eqref{eq:ch47-bernstein-derivative} by
\eqref{eq:bj-energy-definition}.  The identification
\(\Lambda_0=E(\phi)\), \(\Lambda_{n+1}=E(\psi)\) is the same definition.

\emph{Step 4: the derivative.}  Differentiating the Bernstein form,
\[
 \frac{\rd}{\rd t}\sum_{k=0}^{m}\binom{m}{k}(1-t)^{m-k}t^ka_k
 =m\sum_{k=0}^{m-1}\binom{m-1}{k}(1-t)^{m-1-k}t^k(a_{k+1}-a_k),
\]
an identity of polynomials verified by comparing coefficients.  With
\(m=n+1\) and \(a_k=\Lambda_k\) this gives
\(g'(t)=(n+1)\sum_{k}\binom{n}{k}(1-t)^{n-k}t^k(\Lambda_{k+1}-\Lambda_k)\).
It remains to identify \((n+1)(\Lambda_{k+1}-\Lambda_k)\) with
\(\lambda_k\).  By multilinearity,
\[
 \begin{aligned}
  (n+1)V(\Lambda_{k+1}-\Lambda_k)
   &=\bigl[(\omega,\psi)-(\omega,\phi)\bigr]
     (\omega,\phi)^{n-k}(\omega,\psi)^k\\
   &=(0,\psi-\phi)(\omega,\phi)^{n-k}(\omega,\psi)^k,
 \end{aligned}
\]
because the two classes cancel and only the potentials differ; and by the
duality \eqref{eq:bj-mixed-ma-duality} of \Cref{thm:bj-mixed-ma} the
right-hand side is
\(V\int(\psi-\phi)\,V^{-1}(\omega+\ddc\phi)^{n-k}\wedge(\omega+\ddc\psi)^k\),
which is \(V\lambda_k\).  Expanding
\(\MA_\omega(\phi_t)=V^{-1}\bigl[(1-t)(\omega+\ddc\phi)+t(\omega+\ddc\psi)\bigr]^n\)
gives \(\int(\psi-\phi)\MA_\omega(\phi_t)
=\sum_k\binom nk(1-t)^{n-k}t^k\lambda_k=g'(t)\).

\emph{Step 5: the second derivative and the sign.}  Applying the same
differentiation identity to \(g'\) with \(m=n\) and \(a_k=\lambda_k\) gives
the displayed formula for \(g''\) with \(e_k=\lambda_k-\lambda_{k+1}\).
By the computation of Step~4 applied twice,
\[
 \begin{aligned}
 V(\lambda_k-\lambda_{k+1})
 &=(0,\psi-\phi)\bigl[(\omega,\phi)-(\omega,\psi)\bigr]
    (\omega,\phi)^{n-1-k}(\omega,\psi)^k\\
 &=-(0,f)(0,f)(\omega,\phi)^{n-1-k}(\omega,\psi)^k.
 \end{aligned}
\]
with \(f=\psi-\phi\).  By the identification of the Dirichlet form recorded
in \Cref{prop:ch46-dirichlet-vertical} --- the form is \emph{minus} the
intersection pairing of the two vertical slots --- the right-hand side is
\(\langle f,f\rangle_{\mathrm D,k}\), which is \(\ge0\) by
\Cref{cor:ch46-dirichlet-e1}.  Hence \(e_k\ge0\) and \(g''\le0\).

\emph{Step 6: the two summation identities.}  The first is a telescoping
sum, \(\sum_{k=0}^{n-1}(\lambda_k-\lambda_{k+1})=\lambda_0-\lambda_n\),
together with \(\lambda_0=\int(\psi-\phi)\MA(\phi)\),
\(\lambda_n=\int(\psi-\phi)\MA(\psi)\) and the definition of \(I\).  The
second is
\(g(1)-g(0)=\int_0^1g'(t)\,\rd t\) and
\(\int_0^1\binom nk(1-t)^{n-k}t^k\,\rd t=\binom nk\frac{k!(n-k)!}{(n+1)!}
=\frac1{n+1}\); it is the cocycle formula
\eqref{eq:bj-energy-cocycle} in the present notation.
\end{proof}

\begin{remark}[What the expansion says, in words]
\label{rem:ch47-expansion-meaning}
The sequence \(\lambda_0\ge\lambda_1\ge\dots\ge\lambda_n\) is the sequence
of ``slopes'' of the energy along the segment from \(\phi\) to \(\psi\),
measured against the \(n+1\) mixed Monge--Amp\`ere measures interpolating
between \(\MA(\phi)\) and \(\MA(\psi)\).  Concavity of the energy is
precisely the monotonicity of this sequence; the total drop
\(\lambda_0-\lambda_n\) is \(I(\phi,\psi)\); and the cocycle formula says
that the energy increment is the \emph{average} of the slopes.  From this
picture the comparison \eqref{eq:bj-I-J-comparison} is immediate --- an
average of a monotone sequence lies between the extreme terms in a
quantitatively controlled way --- and so is the uniform concavity
\eqref{eq:bj-uniform-concavity}, which says that a polynomial with a
prescribed total drop of slope cannot be too close to affine.  The rest of
this section is these two observations made precise.
\end{remark}

\subsection{The comparison between \texorpdfstring{\(I\)}{I} and
\texorpdfstring{\(J\)}{J}}
\label{sec:ch47-comparison}

\begin{proposition}\label{prop:bj-energy-I-identities}
For \(\phi,\psi\in\mathcal E^1\),

\[
 \frac1{n+1}I(\phi,\psi)\leq J_\psi(\phi)
       \leq\frac n{n+1}I(\phi,\psi),
\tag{47.7}\label{eq:bj-I-J-comparison}
\]

and there is a dimensional constant \(C_n\) such that

\[
 I(\phi_1,\phi_2)
 \leq C_n\bigl(I(\phi_1,\phi_3)+I(\phi_3,\phi_2)\bigr).
\tag{47.8}\label{eq:bj-I-quasitriangle}
\]

Moreover \(I(\phi,\psi)=0\) if and only if \(\phi-\psi\) is constant.
\end{proposition}

\begin{proof}
\emph{The comparison \eqref{eq:bj-I-J-comparison}.}  Adopt the notation of
\Cref{thm:ch47-segment-expansion}.  By its last two identities,
\[
 J_\psi(\phi)=\bigl(E(\psi)-E(\phi)\bigr)-\lambda_n
 =\frac1{n+1}\sum_{k=0}^{n}\lambda_k-\lambda_n
 =\frac1{n+1}\sum_{k=0}^{n}(\lambda_k-\lambda_n),
\]
using \(\lambda_n=\int(\psi-\phi)\MA(\psi)=-\int(\phi-\psi)\MA(\psi)\).
The sequence \((\lambda_k)_{k=0}^n\) is nonincreasing, because
\(\lambda_k-\lambda_{k+1}=e_k\ge0\); hence every term of the last sum lies
in \([0,\lambda_0-\lambda_n]=[0,I(\phi,\psi)]\), and the term with
\(k=n\) vanishes.  Therefore
\[
 \frac{1}{n+1}\,I(\phi,\psi)
 =\frac{1}{n+1}(\lambda_0-\lambda_n)
 \ \le\ J_\psi(\phi)\ \le\ \frac{n}{n+1}(\lambda_0-\lambda_n)
 =\frac{n}{n+1}I(\phi,\psi),
\]
the lower bound retaining only the term \(k=0\) and the upper bound
estimating each of the \(n\) nonzero terms by the largest.  Interchanging
\(\phi\) and \(\psi\) --- which reverses the segment and replaces
\((\lambda_k)\) by \((-\lambda_{n-k})\) --- gives the same two-sided bound
for \(J_\phi(\psi)\); this is consistent with
\(J_\psi(\phi)+J_\phi(\psi)=I(\phi,\psi)\), and shows that both bounds are
attained simultaneously only when \(n=1\), where they coincide.

\emph{Sharpness.}  The constants are optimal.  If \(\lambda_0=\dots
=\lambda_{n-1}\) and \(\lambda_n\) is strictly smaller --- which happens
when all the Dirichlet forms \(e_0,\dots,e_{n-2}\) vanish and \(e_{n-1}\)
does not --- then \(J_\psi(\phi)=\frac n{n+1}I\); the reversed
configuration gives the lower bound.  For \(n=1\) both constants equal
\(\tfrac12\) and \eqref{eq:bj-I-J-comparison} is the identity
\(J_\psi(\phi)=J_\phi(\psi)=\tfrac12I(\phi,\psi)\), which is visible in
\Cref{ex:ch47-p1-numbers}.

\emph{The quasi-triangle inequality \eqref{eq:bj-I-quasitriangle}.}  This
is the Hodge-type estimate \cite[Theorem~3.31]{BJ22}, applied to decreasing
finite-energy regularizations and passed to the limit; it is
item~\cref{itm:ch47-import:quasitriangle} of \Cref{imp:ch47-imported} and is
the only assertion of this section that is not proved here.  What makes it
a theorem rather than a computation is that the quadratic forms
\(\langle\,\cdot\,,\,\cdot\,\rangle_{\mathrm D,k}\) of
\Cref{thm:ch47-segment-expansion} depend on the two potentials being
compared, so that \(I^{1/2}\) is not the norm of a \emph{fixed} inner
product and Cauchy--Schwarz does not immediately apply.  For \(n=1\) the
forms do not depend on the potentials, \(I^{1/2}\) is a genuine seminorm,
and the sharp constant is \(C_1=2\); this is
\Cref{prop:ch47-curve-I-norm}.

\emph{The vanishing clause.}  If \(\phi-\psi\) is constant then
\(\MA(\phi)=\MA(\psi)\) and the integrand of \(I\) is constant, so
\(I(\phi,\psi)=0\).  Conversely suppose \(I(\phi,\psi)=0\).  By
\Cref{thm:ch47-segment-expansion} all the numbers
\(e_k=\langle f,f\rangle_{\mathrm D,k}/V\) vanish, with \(f=\psi-\phi\).
By the Cauchy--Schwarz inequality of \Cref{cor:ch46-dirichlet-e1}, applied
to the \(k\)-th form, \(\langle f,g\rangle_{\mathrm D,k}=0\) for every
\(g\in\vec\cE^{\,1}(X)\); taking \(g\in\PL_{\mathbf R}(X)\) and using
integration by parts \eqref{eq:bj-integration-by-parts} we get
\[
 \int g\,\ddc f\wedge(\omega+\ddc\phi)^{n-1-k}\wedge(\omega+\ddc\psi)^{k}=0
 \qquad\text{for all }g\in\PL_{\mathbf R}(X).
\]
Since \(\PL_{\mathbf R}(X)\) is uniformly dense in \(C^0(X^{\An})\)
(\Cref{thm:ch34-pl-density}), the signed Radon measure
\(\ddc f\wedge(\omega+\ddc\phi)^{n-1-k}\wedge(\omega+\ddc\psi)^{k}\)
vanishes.  Taking \(k=0,\dots,n-1\) in turn and telescoping, all \(n+1\)
mixed measures interpolating between \(\MA(\phi)\) and \(\MA(\psi)\)
coincide; in particular
\[
 \MA_\omega(\phi)=\MA_\omega(\psi).
\]
That \(\phi-\psi\) is then constant is the finite-energy uniqueness theorem
\cite[Corollary~10.4]{BJ22}, item~\cref{itm:ch47-import:uniqueness} of
\Cref{imp:ch47-imported}; it is proved here for \(n=1\) in
\Cref{prop:ch47-curve-uniqueness}.  Integrality of \(X\) is used at this
last step; for an equidimensional reducible variety the conclusion is
``locally constant''.
\end{proof}

\begin{corollary}
\label{cor:ch47-J-vanishes}
For \(\phi\in\cE^1(\omega)\) one has
\(\tfrac1{n+1}I(\phi,\sup\phi)\le J(\phi)\le\tfrac n{n+1}I(\phi,\sup\phi)\),
and \(J(\phi)=0\) if and only if \(\phi\) is constant.  Consequently, if
\(\phi\in\PSH(\omega)\) satisfies \(E_\omega(\phi)=\sup_{X^{\An}}\phi\),
then \(\phi\) is constant.
\end{corollary}

\begin{proof}
Apply \eqref{eq:bj-I-J-comparison} with \(\psi\) the constant
\(c:=\sup_{X^{\An}}\phi\); then
\(J_c(\phi)=E(c)-E(\phi)+\int(\phi-c)\MA(c)=c-E(\phi)+(\phi(v_{\triv})-c)
=\sup\phi-E(\phi)=J(\phi)\),
using \Cref{lem:ch47-three-normalizations}\,\cref{itm:ch47-norm:trivial-measure}
and \cref{itm:ch47-norm:max}.  Hence \(J(\phi)=0\) forces
\(I(\phi,c)=0\), so \(\phi-c\) is constant by
\Cref{prop:bj-energy-I-identities}, that is, \(\phi\) is constant.  For the
last assertion, \(E_\omega(\phi)=\sup\phi\) is finite, so
\(\phi\in\cE^1(\omega)\) and \(J(\phi)=0\).
\end{proof}

\begin{remark}[This settles a statement deferred from \Cref{chap:bj-energy}]
\label{rem:ch47-settles-ch46}
\Cref{rem:ch46-equality-scope} recorded that the equality clause of
\Cref{prop:ch46-energy-usc}\,\cref{itm:ch46-Eusc:strict} --- proved there
only for \(\phi\in\cH_{\mathrm{dom}}(\omega)\), by the equality case of the
Hodge index inequality on one model --- extends to arbitrary
\(\phi\in\PSH(\omega)\), and deferred the extension to this chapter.
\Cref{cor:ch47-J-vanishes} is that extension.  Note the logical shape: the
model-theoretic argument of \Cref{chap:bj-energy} is quantitative and
local to one test configuration, while the argument here is a formal
consequence of the comparison \eqref{eq:bj-I-J-comparison} and of the
uniqueness theorem, and is therefore available exactly when the latter is.
\end{remark}

\subsection{The curve case, in full}
\label{sec:ch47-curve-case}

Two of the three assertions of \Cref{prop:bj-energy-I-identities} that were
imported become elementary when \(n=1\), and it is worth seeing why: the
Dirichlet form ceases to depend on the potentials.

\begin{proposition}
\label{prop:ch47-curve-I-norm}
Let \(n=\dim X=1\).  For \(f,g\in\vec\cE^{\,1}(X)\) set
\(\langle f,g\rangle_{\mathrm D}:=-\int_{X^{\An}}f\,\ddc g\), a symmetric,
positive semidefinite bilinear form independent of any potential.  Then for
\(\phi,\psi\in\cE^1(\omega)\),
\[
 I(\phi,\psi)=V^{-1}\langle\phi-\psi,\phi-\psi\rangle_{\mathrm D},
 \qquad
 J_\psi(\phi)=J_\phi(\psi)=\tfrac12 I(\phi,\psi),
\]
the map \(\phi\mapsto\phi\) induces an injection of
\(\cE^1(\omega)/\mathbf R\) into the pre-Hilbert space
\((\vec\cE^{\,1}(X)/\mathbf R,\langle\,\cdot\,,\,\cdot\,\rangle_{\mathrm D})\),
and \eqref{eq:bj-I-quasitriangle} holds with \(C_1=2\), this being optimal.
\end{proposition}

\begin{proof}
For \(n=1\) there are no auxiliary pairs, so the Dirichlet form of
\Cref{cor:ch46-dirichlet-e1} is
\(\langle f,g\rangle_{\mathrm D}=-\int f\,\ddc g\); it is symmetric by
\Cref{prop:bj-integration-by-parts} and positive semidefinite with the
Cauchy--Schwarz inequality by \Cref{cor:ch46-dirichlet-e1}.  With
\(f=\psi-\phi\), \Cref{thm:ch47-segment-expansion} has a single number
\(e_0=V^{-1}\langle f,f\rangle_{\mathrm D}\) and
\(I(\phi,\psi)=e_0\).  Since \(\ddc\) annihilates constants,
\(\langle\,\cdot\,,\,\cdot\,\rangle_{\mathrm D}\) descends to
\(\vec\cE^{\,1}(X)/\mathbf R\).  The identity
\(J_\psi(\phi)=J_\phi(\psi)=\tfrac12I\) is
\eqref{eq:bj-I-J-comparison} with \(n=1\), where the two constants agree.

For the quasi-triangle inequality, write
\(\lVert f\rVert_\omega:=
 (V^{-1}\langle f,f\rangle_{\mathrm D})^{1/2}\), a seminorm by
Cauchy--Schwarz.  Then
\[
 \begin{aligned}
 I(\phi_1,\phi_2)&=\lVert\phi_1-\phi_2\rVert_\omega^2\\
 &\le\bigl(\lVert\phi_1-\phi_3\rVert_\omega
              +\lVert\phi_3-\phi_2\rVert_\omega\bigr)^2\\
 &\le2\bigl(\lVert\phi_1-\phi_3\rVert_\omega^2
              +\lVert\phi_3-\phi_2\rVert_\omega^2\bigr)\\
 &=2\bigl(I(\phi_1,\phi_3)+I(\phi_3,\phi_2)\bigr).
 \end{aligned}
\]
and the constant \(2\) is attained in the limit
\(\lVert\phi_1-\phi_3\rVert_\omega
=\lVert\phi_3-\phi_2\rVert_\omega\) with the two
differences proportional.
\end{proof}

\begin{proposition}
\label{prop:ch47-curve-uniqueness}
Let \(\dim X=1\) and \(\phi,\psi\in\cE^1(\omega)\) with
\(\MA_\omega(\phi)=\MA_\omega(\psi)\).  Then \(\phi-\psi\) is constant.  In
particular \(I(\phi,\psi)=0\) forces \(\phi-\psi\) constant, and
\Cref{prop:bj-energy-I-identities} is proved without import when
\(n=1\).
\end{proposition}

\begin{proof}
Put \(f:=\phi-\psi\in\vec\cE^{\,1}(X)\).  From
\(\MA(\phi)=\MA(\psi)\) and \eqref{eq:bj-curve-laplacian} we get
\(\ddc f=0\) as a signed Radon measure on \(X^{\An}\); equivalently, by
\Cref{ex:ch46-curve-laplacian}, the graph Laplacian of \(f\) vanishes.
Hence \(\langle f,f\rangle_{\mathrm D}=-\int f\,\ddc f=0\).

We show directly that \(f\) is constant.  Fix a closed point \(p\in X(k)\)
and consider the restriction of \(f\) to the ray
\(\{t\,\ord_p\mid t\ge0\}\), a difference of two convex functions of
\(t\) by \Cref{ex:ch34-psh-check}.  Vanishing of the graph Laplacian at
every interior point of the ray says that the two outgoing slopes at that
point sum to zero, that is, \(t\mapsto f(t\,\ord_p)\) is affine on
\((0,\infty)\); vanishing at the endpoint \(v_{\triv}\) says that the sum
over \emph{all} rays of the outgoing slopes at \(v_{\triv}\) vanishes.  But
both \(\phi\) and \(\psi\) have finite energy, so
\eqref{eq:bj-e1-sublinear-growth} applies to each of them on
\(X^{\lin}\ni t\,\ord_p\): with \(n=1\) the exponent there is
\(1-2^{-1}=\tfrac12\), and \(T_\omega(t\,\ord_p)=t\,T_\omega(\ord_p)\) by
\eqref{eq:bj-dinfty-homogeneity}, so
\(|f(t\,\ord_p)|=O(\sqrt t)\) as \(t\to\infty\).  An affine function of
\(t\) that is \(O(\sqrt t)\) has slope \(0\).  Hence \(f\) is constant on every ray, and the
common value at \(v_{\triv}\) makes it constant on
\(\bigcup_p\{t\ord_p\}\), which is dense in \(X^{\An}\); as \(f\) is a
difference of upper semicontinuous functions, finite on \(X^{\lin}\),
comparison at divisorial points (\Cref{thm:bj-final-envelope}, or directly
\cite[Theorem~4.22]{BJ22}) gives \(f\) constant on \(X^{\An}\).

Finally, if \(I(\phi,\psi)=0\) then \(\MA(\phi)=\MA(\psi)\) by the proof of
\Cref{prop:bj-energy-I-identities}, and the previous paragraph applies.
\end{proof}

\begin{remark}[Why the curve argument does not generalize]
\label{rem:ch47-curve-not-general}
Two features were used, and both fail for \(n\ge2\).  First, the Dirichlet
form was independent of the potentials, so that \(I^{1/2}\) was a seminorm
on a fixed vector space; for \(n\ge2\) the weights
\((\omega+\ddc\phi)^{n-1-k}\wedge(\omega+\ddc\psi)^k\) vary with the
argument, and comparing two such forms is exactly the content of the
Hodge-type estimates behind \eqref{eq:bj-I-quasitriangle}.  Second, a
function with vanishing Laplacian on a one-dimensional graph is affine on
each edge, and the growth condition then forces it to be constant; in
higher dimension ``\(\ddc f\wedge(\omega+\ddc\phi)^{n-1}=0\)'' is a
degenerate elliptic equation whose solutions are not classified by an
elementary argument.  This is the precise sense in which
item~\cref{itm:ch47-import:uniqueness} of \Cref{imp:ch47-imported} is a
genuine theorem.
\end{remark}

\subsection{Uniform concavity}
\label{sec:ch47-uniform-concavity}

\begin{theorem}\label{thm:bj-uniform-concavity}
For \(t\in[0,1]\) and \(\phi,\psi\in\mathcal E^1\),

\[
 E((1-t)\phi+t\psi)
 -(1-t)E(\phi)-tE(\psi)
 \geq c_n t(1-t)I(\phi,\psi)
\tag{47.9}\label{eq:bj-uniform-concavity}
\]

for a constant \(c_n>0\) depending only on \(n\).
\end{theorem}

\begin{proof}
Keep the notation of \Cref{thm:ch47-segment-expansion} and set
\[
 A(t):=g(t)-(1-t)g(0)-tg(1),
 \qquad t\in[0,1],
\]
so that \eqref{eq:bj-uniform-concavity} asserts
\(A(t)\ge c_nt(1-t)I(\phi,\psi)\).

\emph{Step 1: a Green's function representation.}  The function \(A\) is a
polynomial with \(A(0)=A(1)=0\) and \(A''=g''\).  Put \(F:=-g''\ge0\).  For
\(t\in[0,1]\),
\[
 A(t)=\int_0^1G(t,s)F(s)\,\rd s,
 \qquad
 G(t,s):=\begin{cases}s(1-t),& 0\le s\le t,\\ t(1-s),& t\le s\le1,\end{cases}
\]
as one checks by differentiating the right-hand side twice: its first
derivative is \(-\int_0^tsF+\int_t^1(1-s)F\) and its second is
\(-tF(t)-(1-t)F(t)=-F(t)=A''(t)\), while both sides vanish at \(t=0\) and
\(t=1\).  This is the classical Green's function of \(-\rd^2/\rd t^2\) on
\([0,1]\) with Dirichlet boundary conditions.

\emph{Step 2: a pointwise bound for the kernel.}  For all
\(s,t\in[0,1]\),
\[
 G(t,s)\ \ge\ t(1-t)\,s(1-s).
\]
Indeed for \(s\le t\) one has
\(t(1-t)s(1-s)=s(1-t)\cdot t(1-s)\le s(1-t)=G(t,s)\), since
\(t(1-s)\le1\); the case \(s\ge t\) is symmetric.

\emph{Step 3: integrating the second derivative.}  By
\Cref{thm:ch47-segment-expansion},
\[
 F(s)=n\sum_{k=0}^{n-1}\binom{n-1}{k}(1-s)^{n-1-k}s^ke_k,
\]
with
\(e_k\ge0\).  Using the Beta integral
\[
 \int_0^1s^{k+1}(1-s)^{n-k}\rd s
  =\frac{(k+1)!\,(n-k)!}{(n+2)!},
\]
\[
 \begin{aligned}
  \int_0^1s(1-s)F(s)\,\rd s
   &=n\sum_{k=0}^{n-1}\binom{n-1}{k}
     \frac{(k+1)!\,(n-k)!}{(n+2)!}\,e_k\\
   &=\frac{n!}{(n+2)!}
     \sum_{k=0}^{n-1}(k+1)(n-k)\,e_k,
 \end{aligned}
\]
where the middle expression was simplified by
\(\binom{n-1}{k}(k+1)!(n-k)!=(n-1)!\,(k+1)(n-k)\).  For
\(0\le k\le n-1\) the integer \((k+1)(n-k)\) is at least \(n\), with
equality exactly at the two ends \(k=0\) and \(k=n-1\).  Hence, using
\(\sum_ke_k=I(\phi,\psi)\),
\[
 \int_0^1s(1-s)F(s)\,\rd s
 \ \ge\ \frac{n\cdot n!}{(n+2)!}\,I(\phi,\psi)
 =\frac{n}{(n+1)(n+2)}\,I(\phi,\psi).
\]

\emph{Step 4: conclusion.}  Combining Steps 1--3,
\[
 A(t)\ \ge\ t(1-t)\int_0^1s(1-s)F(s)\,\rd s
 \ \ge\ \frac{n}{(n+1)(n+2)}\,t(1-t)\,I(\phi,\psi),
\]
which is \eqref{eq:bj-uniform-concavity} with
\(c_n=n/((n+1)(n+2))\).
\end{proof}

\begin{corollary}
\label{cor:ch47-concavity-constants}
\Cref{thm:bj-uniform-concavity} holds with
\[
 c_n=\frac{n}{(n+1)(n+2)} ,
\]
and for all \(t\in[0,1]\) and \(\phi,\psi\in\cE^1(\omega)\) one also has
the reverse estimate
\[
 E((1-t)\phi+t\psi)-(1-t)E(\phi)-tE(\psi)\ \le\ t(1-t)\,I(\phi,\psi).
\tag{47.10}\label{eq:ch47-concavity-upper}
\]
For \(n=1\) the two bounds are
\(\tfrac16t(1-t)I\le A(t)\le t(1-t)I\), while the exact value is
\(A(t)=\tfrac12t(1-t)I\); so neither constant is sharp in general, and
\(c_n\) is not claimed to be optimal.
\end{corollary}

\begin{proof}
The value of \(c_n\) is Step~4 above.  For \eqref{eq:ch47-concavity-upper}
apply the right-hand tangent inequality of
\eqref{eq:bj-energy-tangent-inequalities} twice, at the endpoints:
\[
 \begin{aligned}
 E(\phi_t)-E(\phi)
   &\le\int(\phi_t-\phi)\MA(\phi)
     =t\!\int(\psi-\phi)\MA(\phi),\\
 E(\phi_t)-E(\psi)
   &\le(1-t)\!\int(\phi-\psi)\MA(\psi).
 \end{aligned}
\]
Multiplying the first by \(1-t\), the second by \(t\), and adding gives
\[
 A(t)\le t(1-t)\left[\int(\psi-\phi)\MA(\phi)+\int(\phi-\psi)\MA(\psi)\right]
 =t(1-t)I(\phi,\psi).
\]
For \(n=1\) the polynomial \(g\) is quadratic with \(g''=-I\) by
\Cref{thm:ch47-segment-expansion}, so
\(A(t)=-\tfrac12g''t(1-t)=\tfrac12It(1-t)\).
\end{proof}

\begin{remark}[The shape of the estimate]
\label{rem:ch47-concavity-shape}
Both bounds of \Cref{cor:ch47-concavity-constants} say the same thing about
the segment: the energy along it is a concave polynomial whose total drop
of slope is exactly \(I(\phi,\psi)\), and a concave polynomial of degree
\(\le n+1\) with prescribed total drop of slope has a controlled gap above
its chord.  The degree bound is essential.  For a general concave function
with total slope drop \(\Delta\) --- say one whose derivative is constant
except for a jump of size \(\Delta\) concentrated near \(t=1\) --- the gap
above the chord can be made arbitrarily small compared with
\(t(1-t)\Delta\).  So \eqref{eq:bj-uniform-concavity} is not a fact about
concave functions; it is a fact about the energy, and the input is the
polynomiality supplied by multilinearity of the pairing on
\(\vec\cE^{\,1}(X)\).  This is worth recording because
\Cref{prop:ch47-maximizing-cauchy} --- the step that makes every maximizing
net Cauchy --- uses \eqref{eq:bj-uniform-concavity} at the single value
\(t=\tfrac12\), and nothing weaker would do.
\end{remark}

\subsection{The quasi-metric \texorpdfstring{\(\mathbb I\)}{I}}
\label{sec:ch47-quasimetric}

Thus \(I\) is a quasi-metric on \(\mathcal E^1/\mathbf R\), not a metric
on \(\mathcal E^1\).  A convenient non-quotient version is

\[
 \mathbb I(\phi,\psi):=I(\phi,\psi)
                 +|\sup\phi-\sup\psi|.
\tag{47.11}\label{eq:bj-I-normalized}
\]

\begin{proposition}
\label{prop:ch47-quasimetric}
The function \(\mathbb I\) on \(\cE^1(\omega)\times\cE^1(\omega)\) is
symmetric and nonnegative, vanishes exactly on the diagonal, and satisfies
\[
 \mathbb I(\phi_1,\phi_2)\le
 \max\{C_n,1\}\bigl(\mathbb I(\phi_1,\phi_3)+\mathbb I(\phi_3,\phi_2)\bigr)
\]
with \(C_n\) as in \eqref{eq:bj-I-quasitriangle}.  Moreover
\(\mathbb I(\phi,\psi)\le3\sup_{X^{\divis}}|\phi-\psi|\), and \(I\)
descends to a quasi-metric on \(\cE^1(\omega)/\mathbf R\) for which the
quotient map is \(1\)-Lipschitz.
\end{proposition}

\begin{proof}
Symmetry and nonnegativity are
\Cref{lem:ch47-IJ-elementary}\,\cref{itm:ch47-IJ:nonneg},
\cref{itm:ch47-IJ:symmetry}.  If \(\mathbb I(\phi,\psi)=0\) then
\(I(\phi,\psi)=0\), so \(\phi-\psi\) is a constant \(c\) by
\Cref{prop:bj-energy-I-identities}, and
\(|\sup\phi-\sup\psi|=|c|=0\); hence \(\phi=\psi\).  The quasi-triangle
inequality follows by adding \eqref{eq:bj-I-quasitriangle} to the genuine
triangle inequality for \(|\sup\phi-\sup\psi|\).  The sup-norm bound
combines \Cref{lem:ch47-IJ-elementary}\,\cref{itm:ch47-IJ:sup} with
\(|\sup\phi-\sup\psi|\le\sup|\phi-\psi|\), which uses
\Cref{lem:ch47-three-normalizations}\,\cref{itm:ch47-norm:max} to evaluate
both suprema at \(v_{\triv}\in X^{\divis}\).  That \(I\) descends is
\Cref{lem:ch47-IJ-elementary}\,\cref{itm:ch47-IJ:translation}.
\end{proof}

\begin{corollary}
\label{cor:ch47-monotone-I}
Let \(\phi,\psi\in\cE^1(\omega)\) with \(\phi\le\psi\).  Then
\[
 \begin{aligned}
  \frac{1}{n+1}I(\phi,\psi)
   &\le J_\psi(\phi)\le E(\psi)-E(\phi),\\
  I(\phi,\psi)
   &\le(n+1)\bigl(E(\psi)-E(\phi)\bigr).
 \end{aligned}
\]
Consequently, if \((\phi_i)\) is a net in \(\cE^1(\omega)\) with
\(\phi_i\le\phi\) for some \(\phi\in\cE^1(\omega)\) and
\(E(\phi_i)\to E(\phi)\), then \(I(\phi_i,\phi)\to0\); if moreover
\(\phi_i(v_{\triv})\to\phi(v_{\triv})\), then
\(\mathbb I(\phi_i,\phi)\to0\).
\end{corollary}

\begin{proof}
The first chain is \eqref{eq:bj-I-J-comparison} together with
\Cref{lem:ch47-IJ-elementary}\,\cref{itm:ch47-IJ:monotone}.  For the
consequence, apply it with \(\psi=\phi\):
\(I(\phi_i,\phi)\le(n+1)(E(\phi)-E(\phi_i))\to0\).  Under the additional
hypothesis, \(\sup\phi_i=\phi_i(v_{\triv})\to\phi(v_{\triv})=\sup\phi\) by
\Cref{lem:ch47-three-normalizations}\,\cref{itm:ch47-norm:max}, so the
second term of \eqref{eq:bj-I-normalized} also tends to \(0\).
\end{proof}

\begin{remark}[What belongs here and what belongs to \Cref{chap:bj-strong-topologies}]
\label{rem:ch47-strong-deferred}
\Cref{prop:ch47-quasimetric} and \Cref{cor:ch47-monotone-I} are the whole
of the metric geometry that this chapter needs, and they are proved from
\Cref{chap:bj-energy} alone.  Three further statements about the same
gauge belong to the next chapter and are used here only where explicitly
flagged.
\begin{enumerate}
\item\label{itm:ch47-strong:topology}
      That \(\mathbb I\) induces the \emph{strong topology}, that is, that
      \(\mathbb I(\phi_i,\phi)\to0\) if and only if \(\phi_i\to\phi\)
      pointwise on \(X^{\divis}\) with \(E(\phi_i)\to E(\phi)\), is
      \Cref{thm:bj-strong-e1-quasimetric}.  Half of it ---
      that convergence of energies along a monotone net implies
      \(\mathbb I\to0\) --- is \Cref{cor:ch47-monotone-I} above.
\item\label{itm:ch47-strong:complete}
      That \((\cE^1(\omega),\mathbb I)\) is \emph{complete} is
      item~(4) of \Cref{thm:bj-cy-envelope-equivalence}; it is stated here
      because it is one of the four equivalent forms of the envelope
      property, but its proof rests on
      item~\cref{itm:ch47-import:completeness} of
      \Cref{imp:ch47-imported} and is carried out in
      \Cref{chap:bj-strong-topologies}.
\item\label{itm:ch47-strong:d1}
      That the strong topology is induced by a genuine \emph{metric}
      \(d_1\), and not merely by the quasi-metric \(\mathbb I\), is
      \Cref{thm:stability-norm-completion} of
      \Cref{chap:bj-stability}, where \(d_1\) is produced from the norm
      geometry of filtrations; the relation to \(\mathbb I\) is
      bi-Lipschitz comparability, not equality, and no isometry statement
      is made anywhere in this book.
\end{enumerate}
The division is deliberate.  Everything that is used to \emph{solve} the
Monge--Amp\`ere equation is here; everything that is used to \emph{topologize}
the solution space is there.
\end{remark}

\subsection{Homogeneity under the scaling action}
\label{sec:ch47-scaling}

Nothing in \Cref{part:bfj} corresponds to the following statement, and it
is the first of the three places in this chapter where the trivially
valued theory has strictly more structure than the discretely valued one.

\begin{proposition}
\label{prop:ch47-scaling}
Let \(t>0\) and let \(t\cdot\phi\) denote the scaling action
\eqref{eq:function-scaling-action}, with \(t_\star\) the induced action
\eqref{eq:valuation-scaling} on \(X^{\An}\).  For
\(\phi,\psi\in\cE^1(\omega)\):
\begin{enumerate}
\item\label{itm:ch47-scal:e1}
      \(t\cdot\phi\in\cE^1(\omega)\), and
      \(\sup(t\cdot\phi)=t\sup\phi\);
\item\label{itm:ch47-scal:JI}
      \(J_{t\cdot\psi}(t\cdot\phi)=t\,J_\psi(\phi)\),
      \(I(t\cdot\phi,t\cdot\psi)=t\,I(\phi,\psi)\), and
      \(J(t\cdot\phi)=t\,J(\phi)\);
\item\label{itm:ch47-scal:quasi}
      \(\mathbb I(t\cdot\phi,t\cdot\psi)=t\,\mathbb I(\phi,\psi)\);
\item\label{itm:ch47-scal:segment}
      \(t\cdot\bigl((1-s)\phi+s\psi\bigr)=(1-s)(t\cdot\phi)+s(t\cdot\psi)\)
      for \(s\in[0,1]\), so the scaling action carries segments to segments
      and multiplies the polynomial \(g\) of
      \Cref{thm:ch47-segment-expansion} by \(t\).
\end{enumerate}
\end{proposition}

\begin{proof}
\cref{itm:ch47-scal:e1}.  Stability of \(\cE^1(\omega)\) is the last clause
of \Cref{thm:ch46-energy-homogeneous}.  For the supremum, use
\((t\cdot\phi)(v)=t\,\phi(t^{-1}v)\), so that
\(\sup_{X^{\An}}(t\cdot\phi)=t\sup_{X^{\An}}\phi\); alternatively evaluate
at \(v_{\triv}\), which is the fixed point of the action, and apply
\Cref{lem:ch47-three-normalizations}\,\cref{itm:ch47-norm:max}.

\cref{itm:ch47-scal:JI}.  By \eqref{eq:bj-homogeneous-ma-scaling},
\(\MA_\omega(t\cdot\psi)=t_\star\MA_\omega(\psi)\).  Hence, for any
\(u,w\in\cE^1(\omega)\),
\[
 \begin{aligned}
 \int(t\cdot u-t\cdot w)\,\MA_\omega(t\cdot\psi)
 &=\int\bigl[(t\cdot u)(t v)-(t\cdot w)(tv)\bigr]\,
      \rd\MA_\omega(\psi)(v)\\
 &=t\int(u-w)\,\MA_\omega(\psi).
 \end{aligned}
\]
because \((t\cdot u)(tv)=t\,u(v)\).  Combining this with
\(E_\omega(t\cdot u)=t\,E_\omega(u)\)
(\Cref{thm:ch46-energy-homogeneous}) and reading off
\eqref{eq:bj-I-J-potentials} gives the three identities; the one for \(J\)
also uses \cref{itm:ch47-scal:e1}.

\cref{itm:ch47-scal:quasi} is \cref{itm:ch47-scal:JI} plus
\cref{itm:ch47-scal:e1}.

\cref{itm:ch47-scal:segment}.  The action is defined pointwise by
\((t\cdot u)(v)=t\,u(t^{-1}v)\), which is linear in \(u\); apply it to the
convex combination.  The identity for \(g\) follows from
\cref{itm:ch47-scal:e1} and \Cref{thm:ch46-energy-homogeneous}.
\end{proof}

\begin{remark}[Reading \Cref{prop:ch47-scaling}]
\label{rem:ch47-cone-reading}
The proposition says that \(\bigl(\cE^1(\omega)/\mathbf R,I\bigr)\) is a
cone: it carries an action of \(\mathbf R_{>0}\) with a unique fixed point
(the class of the constants, by
\Cref{cor:ch47-J-vanishes}), and the gauge \(I\) is homogeneous of degree
one.  Two consequences are used later.  First, the sublevel sets
\(\{J\le C\}\) are the dilates of \(\{J\le1\}\), so a single compactness
statement at one level gives all of them; this is used implicitly in
\Cref{sec:ch47-cy}.  Second, the Calabi--Yau correspondence will be
equivariant (\Cref{cor:ch47-scaling-cy}), so that solving one equation
solves a whole orbit.  Neither statement has an analogue in
\Cref{chap:bfj-finite-energy-measures}: over a discretely valued base there
is no \(\mathbf G_m\)-action on the model system, and
\Cref{thm:finite-energy-measure-completion} produces a metric space with no
distinguished dilation.
\end{remark}

\subsection{A worked example on \texorpdfstring{\(\mathbf P^1\)}{P1}}
\label{sec:ch47-p1-numbers}

\begin{example}[All the functionals of this section, computed]
\label{ex:ch47-p1-numbers}
Keep the data of \Cref{ex:ch46-p1-energy}: \(X=\mathbf P^1_k\),
\(\omega=c_1(\cO(d))\), \(V=d\), \(p\in X(k)\), and for \(0\le c\le d\)
\[
 \begin{aligned}
  \phi_c&:=c\max\{\log|\mathfrak m_p|,-1\},\\
  \phi_c(s\,\ord_p)&=-c\min\{s,1\},
  &\phi_c&=0\quad\text{off the \(p\)-ray},
 \end{aligned}
\]
with
\[
 E_\omega(\phi_c)=-\frac{c^2}{2d},
 \qquad
 \MA_\omega(\phi_c)=\frac{d-c}{d}\,\delta_{v_{\triv}}
                    +\frac{c}{d}\,\delta_{\ord_p}.
\]
Recall \(\MA_\omega(0)=\delta_{v_{\triv}}\), which is the case \(c=0\), in
agreement with
\Cref{lem:ch47-three-normalizations}\,\cref{itm:ch47-norm:trivial-measure}.

\emph{The functionals.}  For \(0\le c,c'\le d\) the difference
\(\phi_c-\phi_{c'}\) equals \(0\) at \(v_{\triv}\) and \(-(c-c')\) at
\(\ord_p\), while
\(\MA(\phi_{c'})-\MA(\phi_c)=\frac{c-c'}{d}(\delta_{v_{\triv}}
-\delta_{\ord_p})\).  Hence
\[
 I(\phi_c,\phi_{c'})=\frac{(c-c')^2}{d},
 \qquad
 J_{\phi_{c'}}(\phi_c)=J_{\phi_c}(\phi_{c'})=\frac{(c-c')^2}{2d},
 \qquad
 J(\phi_c)=\frac{c^2}{2d}.
\]
The middle identity is \eqref{eq:bj-I-J-comparison} for \(n=1\), where the
two constants both equal \(\tfrac12\); the last one is
\(\sup\phi_c-E(\phi_c)=0+c^2/(2d)\) and confirms
\Cref{lem:ch47-three-normalizations}\,\cref{itm:ch47-norm:J}, since
\(J_0(\phi_c)=E(0)-E(\phi_c)+\int\phi_c\,\MA(0)=c^2/(2d)+\phi_c(v_{\triv})
=c^2/(2d)\).

\emph{The seminorm.}  The map \(c\mapsto[\phi_c]\) is, by the display
above, an isometry of \(([0,d],d^{-1/2}|\cdot|)\) onto its image in
\((\cE^1(\omega)/\mathbf R,I^{1/2})\), as
\Cref{prop:ch47-curve-I-norm} predicts.  In particular the triangle
inequality holds here with constant \(1\), not merely with \(C_1=2\); the
constant \(2\) in \Cref{prop:ch47-curve-I-norm} is needed only because
\(I\) is the \emph{square} of the seminorm.

\emph{Uniform concavity.}  Taking \(\psi=0\) and using
\(t\phi_c=\phi_{tc}\),
\[
 A(t)=E_\omega(t\phi_c)-tE_\omega(\phi_c)
 =-\frac{t^2c^2}{2d}+\frac{tc^2}{2d}
 =\frac{c^2}{2d}\,t(1-t)
 =\tfrac12\,t(1-t)\,I(\phi_c,0),
\]
which lies between the bound \(c_1t(1-t)I=\tfrac16t(1-t)I\) of
\Cref{thm:bj-uniform-concavity} and the bound \(t(1-t)I\) of
\eqref{eq:ch47-concavity-upper}, and equals the exact value
\(\tfrac12t(1-t)I\) recorded in
\Cref{cor:ch47-concavity-constants}.

\emph{Scaling.}  By \Cref{ex:ch46-homogeneity-p1},
\(t\cdot\phi_c=\max\{c\log|\mathfrak m_p|,-ct\}\) and
\(E_\omega(t\cdot\phi_c)=-tc^2/(2d)\).  Therefore
\(J(t\cdot\phi_c)=tc^2/(2d)=t\,J(\phi_c)\) and
\(I(t\cdot\phi_c,0)=tc^2/d=t\,I(\phi_c,0)\), confirming
\Cref{prop:ch47-scaling}\,\cref{itm:ch47-scal:JI}.  Notice that
\(t\cdot\phi_c\ne\phi_{tc}\): the first moves the corner of the graph from
\(\ord_p\) to \(t\,\ord_p\) keeping the slope \(-c\), while the second
keeps the corner and changes the slope to \(-tc\).  Their energies differ
too: \(E_\omega(t\cdot\phi_c)=-tc^2/(2d)\) is homogeneous of degree one in
\(t\), as \Cref{thm:ch46-energy-homogeneous} requires, whereas
\(E_\omega(\phi_{tc})=-t^2c^2/(2d)\) is homogeneous of degree \(n+1=2\).  In
\Cref{ex:ch46-toric-p2} the analogous two operations on \(\mathbf P^2\)
give \(-tc^3/3\) and \(-t^3c^3/3\), with exponents \(1\) and \(n+1=3\); the
pattern is the same in both dimensions, and only the first operation is the
scaling action.
\end{example}

\section{Dual energy of a measure}
\label{sec:ch47-dual}

The Legendre transform of a concave functional on a convex set is a convex
functional on the dual; here the convex set is \(\cE^1(\omega)\), the
functional is \(E_\omega\), and the dual pairing is integration against a
measure.  Because \(E_\omega(\phi+c)=E_\omega(\phi)+c\), the transform is
finite only on measures of total mass one
(\Cref{lem:ch47-mass-forced}), which is why the probability normalization
of \Cref{warn:ch46-normalization} is the natural one for this chapter and
not merely a convenience.

\subsection{The definition and the test class}
\label{sec:ch47-dual-def}

\begin{definition}[Finite-energy measure in the BJ theory]
\label{def:bj-finite-energy-measure}
For \(\mu\in\mathcal M\), define

\[
 E_\omega^\vee(\mu)
 :=\sup_{\phi\in\mathcal E^1}
       \left(E_\omega(\phi)-\int\phi\,\rd \mu\right)
 \in[0,+\infty].
\tag{47.12}\label{eq:bj-dual-energy}
\]

The extended integral in \eqref{eq:bj-dual-energy} is understood in the
usual truncation sense.  Equivalently, the supremum may first be taken over
the bounded class \(\mathcal H^{\mathrm{dom}}(\omega)\); finiteness then
implies that every \(\mathcal E^1\)-function is \(\mu\)-integrable, so the
displayed formula is unambiguous.
We write
\(\mathcal M^1(\omega):=\{\mu\in\mathcal M\mid E_\omega^\vee(\mu)<\infty\}\).
\end{definition}

The supremum can equivalently be taken over
\(\mathcal H^{\mathrm{dom}}(\omega)\), by
decreasing approximation and monotone convergence
\cite[Proposition~9.4]{BJ22}.  It follows at once that \(E^\vee\) is convex
and weakly lower semicontinuous.  Its sublevel sets

\[
             \mathcal M^1_C:=\{\mu\mid E^\vee(\mu)\leq C\}
\tag{47.13}\label{eq:bj-measure-energy-sublevel}
\]

are weakly compact.

The next three statements make all of this precise and prove it.

\begin{lemma}
\label{lem:ch47-mass-forced}
Let \(\mu\) be a positive Radon measure on \(X^{\An}\) of total mass
\(m=\mu(X^{\An})\).  If \(m\ne1\), then
\(\sup_{\phi\in\cE^1(\omega)}\bigl(E_\omega(\phi)-\int\phi\,\rd\mu\bigr)
=+\infty\).  Hence \(E^\vee_\omega\) is finite only on probability
measures, and the restriction to \(\mathcal M\) in
\eqref{eq:bj-e1-m-probability} loses nothing.
\end{lemma}

\begin{proof}
Test on the constants, which lie in \(\cE^1(\omega)\) because \(\omega\) is
nef (\Cref{lem:ch34-psh-nonempty}).  By
\eqref{eq:bj-energy-basic-identities}, \(E_\omega(c)=c\), so
\(E_\omega(c)-\int c\,\rd\mu=c(1-m)\), which tends to \(+\infty\) as
\(c\to+\infty\) if \(m<1\) and as \(c\to-\infty\) if \(m>1\).  This is the
trivially valued form of \Cref{lem:ch28-mass-forced}, where the same test
forces the mass to be \(V\) because the energy there translates by \(cV\).
\end{proof}

\begin{proposition}
\label{prop:ch47-test-class}
Let \(\mu\in\mathcal M\).  Then the supremum
\eqref{eq:bj-dual-energy} is unchanged if \(\cE^1(\omega)\) is replaced by
any of
\[
 \begin{aligned}
 \cH_{\mathrm{dom}}(\omega)
 &\subset \PL^\omega(X)
 \subset \{\phi\in\PSH(\omega)\mid \phi\ \text{continuous}\}\\
 &\subset \{\phi\in\PSH(\omega)\mid \phi\ \text{bounded}\}
 \subset \cE^1(\omega).
 \end{aligned}
\]
Moreover, if \(E^\vee_\omega(\mu)<\infty\), then every
\(\phi\in\cE^1(\omega)\) is \(\mu\)-integrable, and
\[
 \int(-\phi)\,\rd\mu\ \le\ E^\vee_\omega(\mu)+J(\phi)
 \qquad\text{whenever }\sup\nolimits_{X^{\An}}\phi=0 .
\]
\end{proposition}

\begin{proof}
\emph{The inclusions.}  \(\cH_{\mathrm{dom}}(\omega)\subset\PL^\omega(X)\)
is \Cref{def:ch34-hdom}; PL functions are continuous; continuous functions
on the compact space \(X^{\An}\) are bounded; and bounded \(\omega\)-psh
functions have finite energy by \Cref{thm:bj-e1-basic}.  Hence the
corresponding suprema increase from left to right, and it suffices to show
that the supremum over \(\cH_{\mathrm{dom}}(\omega)\) dominates the one
over \(\cE^1(\omega)\).

\emph{The reverse inequality.}  Let \(\phi\in\cE^1(\omega)\).  By
\Cref{lem:ch46-continuous-competitors} --- \Cref{lem:ch46-pl-domination}
being available only when \(\omega\) is the class of an ample
\(\mathbf Q\)-line bundle, whereas the standing hypothesis here is
\(\omega\in\Amp(X)\) --- the family
\(\mathcal F_{\mathrm{dom}}(\phi)=\{\psi\in C^0(X^{\An})\cap\PSH(\omega)\mid
\psi\ge\phi\}\) is nonempty and, indexed as in Step~6 of the proof of
\Cref{thm:bj-extended-energy-pairing}, produces a decreasing net
\((\phi_j)\) in \(\cH_{\mathrm{dom}}(\omega)\) whose pointwise limit agrees
with \(\phi\) on \(X^{\divis}\).  That limit is \(\omega\)-psh by
\Cref{thm:bj-basic-closure}\,(4), so it equals \(\phi\) on all of
\(X^{\An}\) by comparison at divisorial points
\cite[Theorem~4.22]{BJ22}; thus \(\phi_j\downarrow\phi\) pointwise on
\(X^{\An}\).  Then
\(E_\omega(\phi_j)\to E_\omega(\phi)\) by
\Cref{prop:ch46-energy-usc}\,\cref{itm:ch46-Eusc:decreasing}, and
\(\int\phi_j\,\rd\mu\downarrow\int\phi\,\rd\mu\) in
\([-\infty,\sup\phi_0]\) by monotone convergence for decreasing nets of
upper semicontinuous functions bounded above
(\Cref{cor:appE-usc-monotone}).  If \(\int\phi\,\rd\mu>-\infty\) the two
limits combine and
\(E_\omega(\phi_j)-\int\phi_j\,\rd\mu\to E_\omega(\phi)-\int\phi\,\rd\mu\);
if \(\int\phi\,\rd\mu=-\infty\), then
\(E_\omega(\phi_j)-\int\phi_j\,\rd\mu\to+\infty\) because
\(E_\omega(\phi_j)\) is bounded below by \(E_\omega(\phi)>-\infty\), and
both sides of the asserted equality are \(+\infty\).  In either case the
supremum over \(\cH_{\mathrm{dom}}(\omega)\) is at least
\(E_\omega(\phi)-\int\phi\,\rd\mu\).  This is
\cite[Proposition~9.4]{BJ22}.

\emph{Integrability.}  Let \(\phi\in\cE^1(\omega)\) with \(\sup\phi=0\) and
let \(\phi_t:=\max\{\phi,-t\}\) for \(t>0\), a bounded \(\omega\)-psh
function with \(\sup\phi_t=0\) and \(\phi_t\ge\phi\), so that
\(E_\omega(\phi_t)\ge E_\omega(\phi)\) and hence
\(J(\phi_t)=-E_\omega(\phi_t)\le-E_\omega(\phi)=J(\phi)\).  Applying
\eqref{eq:bj-dual-energy} to \(\phi_t\),
\[
 \int(-\phi_t)\,\rd\mu
 =E_\omega(\phi_t)-\int\phi_t\,\rd\mu-E_\omega(\phi_t)
 \le E^\vee_\omega(\mu)+J(\phi_t)\le E^\vee_\omega(\mu)+J(\phi).
\]
Letting \(t\uparrow\infty\), monotone convergence gives
\(\int(-\phi)\,\rd\mu\le E^\vee_\omega(\mu)+J(\phi)<\infty\); as
\(-\phi\ge0\), the function \(\phi\) is \(\mu\)-integrable.  A general
\(\phi\in\cE^1(\omega)\) differs from a normalized one by a constant.
\end{proof}

\begin{remark}[Why the supremum is not taken over \(\PSH(\omega)\)]
\label{rem:ch47-why-e1-test}
If \(\phi\in\PSH(\omega)\setminus\cE^1(\omega)\) and
\(\int\phi\,\rd\mu=-\infty\), the expression
\(E_\omega(\phi)-\int\phi\,\rd\mu\) is \((-\infty)-(-\infty)\) and is
forbidden by the front matter.  Restricting the test class to
\(\cE^1(\omega)\) removes the ambiguity, and
\Cref{prop:ch47-test-class} shows that nothing is lost: the value of the
supremum is already reached on the smallest of the five classes.  This is
ledger entry \cref{itm:ch47-ledger:integral} of \Cref{sec:ch47-ledger}, and
it is also why \Cref{def:dual-energy-measure} in \Cref{part:bfj} restricts
its test class in the same way and proves the analogous
\Cref{prop:ch28-test-class}.  The proof there is the same monotone
approximation; only the approximating family differs, being model
functions on a varying model rather than members of
\(\cH_{\mathrm{dom}}(\omega)\).
\end{remark}

\subsection{Elementary structure of \texorpdfstring{\(E^\vee_\omega\)}{the dual energy}}
\label{sec:ch47-dual-elementary}

\begin{proposition}
\label{prop:ch47-dual-elementary}
The functional \(E^\vee_\omega:\mathcal M\to[0,+\infty]\) has the following
properties.
\begin{enumerate}
\item\label{itm:ch47-dual:nonneg}
      \(E^\vee_\omega\ge0\), and \(E^\vee_\omega\) is unchanged if
      \(\phi\) is replaced by \(\phi+c\) in
      \eqref{eq:bj-dual-energy}; thus the functional
      \(F_\mu:=E_\omega-\int\cdot\,\rd\mu\) is invariant under translation
      of its argument by real constants.
\item\label{itm:ch47-dual:convex}
      \(E^\vee_\omega\) is convex on the convex set \(\mathcal M\).
\item\label{itm:ch47-dual:lsc}
      \(E^\vee_\omega\) is lower semicontinuous for the weak topology on
      \(\mathcal M\).
\item\label{itm:ch47-dual:compact}
      Every sublevel set \eqref{eq:bj-measure-energy-sublevel} is weakly
      compact, and \(\cM^1(\omega)=\bigcup_{C>0}\cM^1_C\) is a convex,
      \(\sigma\)-compact subset of \(\mathcal M\).
\end{enumerate}
\end{proposition}

\begin{proof}
\cref{itm:ch47-dual:nonneg}.  The constant \(0\) lies in
\(\cE^1(\omega)\) and gives \(E_\omega(0)-\int0\,\rd\mu=0\), so the
supremum is \(\ge0\).  For the translation statement, \(\mu\) has mass one,
so \(E_\omega(\phi+c)-\int(\phi+c)\rd\mu
=E_\omega(\phi)+c-\int\phi\,\rd\mu-c\).  This is the exact analogue of
\eqref{eq:dual-energy-constant-test}, with \(V\) replaced by \(1\).

\cref{itm:ch47-dual:convex}.  For fixed \(\phi\in\cE^1(\omega)\) the map
\(\mu\mapsto E_\omega(\phi)-\int\phi\,\rd\mu\) is affine on \(\mathcal M\),
with values in \([-\infty,+\infty)\).  A supremum of affine functions is
convex.

\cref{itm:ch47-dual:lsc}.  Each \(\phi\in\cE^1(\omega)\) is upper
semicontinuous and bounded above, so by the Portmanteau theorem
\Cref{prop:portmanteau}, in the form
\eqref{eq:appE-portmanteau-lsc} of
\Cref{prop:appE-measure-compactness}, the map
\(\mu\mapsto\int\phi\,\rd\mu\) is upper semicontinuous on
\(\mathcal M\); hence \(\mu\mapsto E_\omega(\phi)-\int\phi\,\rd\mu\) is
lower semicontinuous, and so is the supremum over \(\phi\).

\cref{itm:ch47-dual:compact}.  The space \(\mathcal M\) of Radon
probability measures on the compact Hausdorff space \(X^{\An}\) is weakly
compact by \Cref{prop:appE-measure-compactness}.  A sublevel set of a
lower semicontinuous function is closed, hence compact.  Convexity of
\(\cM^1(\omega)\) is \cref{itm:ch47-dual:convex}, and the countable union is
over \(C\in\mathbf N\).
\end{proof}

\begin{remark}[Weak compactness is not strong compactness]
\label{rem:ch47-weak-not-strong}
\Cref{prop:ch47-dual-elementary}\,\cref{itm:ch47-dual:compact} is a
statement about the \emph{weak} topology on measures, and it is cheap: it
uses only that \(X^{\An}\) is compact and that \(E^\vee_\omega\) is a
supremum of weakly continuous affine functions.  It is emphatically not a
statement about the strong topology of
\Cref{def:bj-measure-quasimetric}, for which \(\cM^1_C\) is not compact ---
the same phenomenon as \Cref{rem:ch28-weak-not-strong} over a discretely
valued base.  The distinction matters in \Cref{sec:ch47-cy}: the direct
method produces a maximizer from weak compactness, but identifying the
limit as a solution of the Monge--Amp\`ere equation is a strong statement
and needs \Cref{thm:bj-energy-differentiability}.
\end{remark}

\subsection{The Fenchel identities}
\label{sec:ch47-fenchel}

\begin{proposition}\label{prop:bj-fenchel-equalities}
For \(\phi\in\mathcal E^1\), one has

\[
 E^\vee(\MA(\phi))
  =E(\phi)-\int\phi\MA(\phi),
\tag{47.14}\label{eq:bj-dual-energy-ma}
\]

and, for every \(\mu\in\mathcal M\),

\[
 E(\phi)\leq E^\vee(\mu)+\int\phi\,\rd \mu.
\tag{47.15}\label{eq:bj-fenchel-inequality}
\]

Equality in \eqref{eq:bj-fenchel-inequality}, with
\(\mu\in\mathcal M^1\), holds if and only if \(\mu=\MA(\phi)\).
\end{proposition}

\begin{proof}
The tangent inequality at \(\phi\) gives, for every \(\psi\in\mathcal E^1\),

\[
 E(\psi)-\int\psi\MA(\phi)
 \leq E(\phi)-\int\phi\MA(\phi),
\]

which proves \eqref{eq:bj-dual-energy-ma}.  Inequality
\eqref{eq:bj-fenchel-inequality} is the definition.  If equality holds,
the constant net \(\phi\) is maximizing for \(\mu\);
\Cref{thm:bj-maximizers-converge} below gives \(\MA(\phi)=\mu\).  The reverse
implication is \eqref{eq:bj-dual-energy-ma}.  This is
\cite[Proposition~9.6 and Corollary~9.12]{BJ22}.

We expand each of the three assertions, since all of them are used
constantly and the first two are complete arguments requiring no import.

\emph{Proof of \eqref{eq:bj-dual-energy-ma}.}  The displayed inequality is
the left-hand tangent inequality of
\eqref{eq:bj-energy-tangent-inequalities}, namely
\(\int(\psi-\phi)\MA(\phi)\le E(\psi)-E(\phi)\), rearranged; it says
precisely that the affine functional
\(\psi\mapsto E(\psi)-\int\psi\,\MA(\phi)\) on \(\cE^1(\omega)\) is
maximized at \(\psi=\phi\).  Its value there is
\(E(\phi)-\int\phi\,\MA(\phi)\), so the supremum
\eqref{eq:bj-dual-energy} for \(\mu=\MA(\phi)\) equals that number.  All
integrals are finite by \Cref{lem:ch47-IJ-elementary}.  In particular
\(\MA_\omega(\phi)\in\cM^1(\omega)\) for every \(\phi\in\cE^1(\omega)\), so
\(\MA_\omega\) does map \(\cE^1(\omega)\) into \(\cM^1(\omega)\); the
content of \Cref{thm:bj-cy-envelope-equivalence} is that it maps
\emph{onto} it.

\emph{The dual energy of a Monge--Amp\`ere measure, in terms of \(J\).}
Rewriting \eqref{eq:bj-dual-energy-ma} with
\Cref{lem:ch47-three-normalizations}, and comparing with
\eqref{eq:bj-I-J-potentials},
\[
 E^\vee_\omega(\MA_\omega(\phi))=J_\phi(0)
 =E(\phi)-E(0)+\int(0-\phi)\MA(\phi),
\]
so by \eqref{eq:bj-I-J-comparison},
\[
 \frac{1}{n+1}\,I(\phi,0)\ \le\ E^\vee_\omega(\MA_\omega(\phi))
 \ \le\ \frac{n}{n+1}\,I(\phi,0).
\]
Thus the dual energy of \(\MA_\omega(\phi)\) is comparable to the
\(I\)-distance from \(\phi\) to the constants --- the first indication that
\(\MA_\omega\) is bi-Lipschitz for the two quasi-metrics, which is
assertion (3) of \Cref{thm:bj-cy-envelope-equivalence}.

\emph{Proof of \eqref{eq:bj-fenchel-inequality} and of the equality case.}
The inequality is \eqref{eq:bj-dual-energy} applied to the single
competitor \(\phi\), together with the convention of ledger entry
\cref{itm:ch47-ledger:integral}; it holds for every
\(\mu\in\mathcal M\), including those with \(E^\vee_\omega(\mu)=+\infty\).
Suppose now \(\mu\in\cM^1(\omega)\) and equality holds.  Then \(\phi\)
attains the supremum \eqref{eq:bj-dual-energy}, so the constant net with
value \(\phi\) is a maximizing net in the sense of
\eqref{eq:bj-maximizing-net}, and \Cref{thm:bj-maximizers-converge} gives
\(\MA_\omega(\phi)=\lim\MA_\omega(\phi)=\mu\).  Conversely if
\(\mu=\MA_\omega(\phi)\), then \eqref{eq:bj-dual-energy-ma} is exactly the
asserted equality.  The forward implication is the only place in this
subsection where \Cref{imp:ch47-imported}\,\cref{itm:ch47-import:diff} is
used, through \Cref{thm:bj-maximizers-converge}; the reader should check
that no circularity results, since the proof of
\Cref{thm:bj-maximizers-converge} uses only
\eqref{eq:bj-fenchel-inequality} and
\Cref{thm:bj-energy-differentiability}.
\end{proof}

\begin{corollary}
\label{cor:ch47-dual-zero}
For \(\mu\in\mathcal M\) one has \(E^\vee_\omega(\mu)=0\) if and only if
\(\mu=\delta_{v_{\triv}}\).  In particular \(E^\vee_\omega\) is a strictly
positive functional on \(\mathcal M\setminus\{\delta_{v_{\triv}}\}\), and
\(\delta_{v_{\triv}}\) is the unique measure whose Monge--Amp\`ere
potential is constant.
\end{corollary}

\begin{proof}
If \(\mu=\delta_{v_{\triv}}=\MA_\omega(0)\)
(\Cref{lem:ch47-three-normalizations}\,\cref{itm:ch47-norm:trivial-measure}),
then \eqref{eq:bj-dual-energy-ma} with \(\phi=0\) gives
\(E^\vee_\omega(\mu)=E_\omega(0)-\int0\,\MA(0)=0\).  Conversely, if
\(E^\vee_\omega(\mu)=0\), then \(\mu\in\cM^1(\omega)\) and equality holds
in \eqref{eq:bj-fenchel-inequality} for \(\phi=0\), because both sides are
\(0\); by \Cref{prop:bj-fenchel-equalities}, \(\mu=\MA_\omega(0)
=\delta_{v_{\triv}}\).
\end{proof}

\subsection{Finite-energy measures and polar sets}
\label{sec:ch47-polar}

Finite-energy measures put no mass on pluripolar sets.  Indeed a pole
function in \(\mathcal E^1\), after bounded truncation and monotone
convergence, would otherwise have integral \(-\infty\), contradicting the
finite Fenchel bound.  In particular they do not charge a proper closed
subvariety of smaller dimension.

The following proposition makes the first sentence quantitative and proves
it.

\begin{proposition}
\label{prop:ch47-no-polar-mass}
Let \(\mu\in\cM^1(\omega)\) and let \(\phi\in\cE^1(\omega)\) with
\(\sup_{X^{\An}}\phi=0\).  Then for every \(t>0\)
\[
 \mu\bigl(\{\phi\le-t\}\bigr)\ \le\ \frac{E^\vee_\omega(\mu)+J(\phi)}{t},
\]
and consequently
\[
 \mu\bigl(\{\phi=-\infty\}\bigr)=0 .
\]
The same holds for any \(\phi\in\cE^1(\omega)\) after normalization.  In
particular, if \(\dim X=1\), then \(\mu\) charges no proper-support point
of \(X^{\An}\).
\end{proposition}

\begin{proof}
By \Cref{prop:ch47-test-class}, \(\int(-\phi)\,\rd\mu\le
E^\vee_\omega(\mu)+J(\phi)=:B<\infty\).  Since \(-\phi\ge0\), Chebyshev's
inequality gives \(t\,\mu(\{-\phi\ge t\})\le\int(-\phi)\rd\mu\le B\), which
is the displayed estimate; the set \(\{\phi\le-t\}\) is Borel because
\(\phi\) is upper semicontinuous.  Letting \(t\to\infty\) and using
\(\{\phi=-\infty\}=\bigcap_t\{\phi\le-t\}\) gives \(\mu(\{\phi=-\infty\})=0\).

For the last assertion let \(\dim X=1\) and let \(x\in X^{\An}\) be a point
with proper support, so that \(x\) is the endpoint at infinity of the ray
through \(\ord_p\) for some closed point \(p\in X(k)\)
(\Cref{ex:ch34-psh-check}).  Put \(d:=\deg\omega>0\), fix
\(a\in(0,d]\), and let \(\phi\) vanish off that ray and satisfy
\(\phi(t\,\ord_p)=-a\log(1+t)\).  The function \(h(t)=-a\log(1+t)\) is
convex, decreasing, and has \(h'(0^+)=-a\ge-d\), so \(\phi\) is
\(\omega\)-psh by the curve criterion recalled in
\Cref{ex:ch34-psh-check}; and by the Dirichlet computation of
\Cref{ex:ch46-e1-curve}, in the form
\(E^{\mathrm u}_\omega(\phi)=\int\phi\,\MA^{\mathrm u}_\omega(0)
-\tfrac12\langle\phi,\phi\rangle_{\mathrm D}
=-\tfrac12\int_0^\infty|h'|^2\,\rd t\) supplied by
\Cref{prop:ch46-dirichlet-vertical}\,\cref{itm:ch46-dir:energy},
\[
 E_\omega(\phi)=-\frac{1}{2d}\int_0^\infty\frac{a^2}{(1+t)^2}\,\rd t
 =-\frac{a^2}{2d}>-\infty .
\]
So \(\phi\in\cE^1(\omega)\), while
\(\phi(x)=\lim_{t\to\infty}\bigl(-a\log(1+t)\bigr)=-\infty\).  Apply the
previous paragraph.
\end{proof}

\begin{remark}
\label{rem:ch47-polar-scope}
\Cref{prop:ch47-no-polar-mass} first gives vanishing on the pole set of
each finite-energy psh function.  The following distinctions matter.
\begin{enumerate}
\item\label{itm:ch47-polar:pluripolar}
      The measure also gives no mass to any pluripolar set.
      Indeed, \cite[Corollary~7.48]{BJ22} places each such set inside
      the pole set of a finite-energy psh function, to which the
      proposition applies.  This additional input does not require
      equality of the original function with its finite-energy witness.
\item\label{itm:ch47-polar:lin}
      Every atom of $\mu$ belongs to $X^{\lin}$: a point outside that
      locus is pluripolar by \Cref{cor:ch49-nonpluripolar}, so it cannot
      have positive mass.  Full mass on $X^{\lin}$ follows under the
      additional hypotheses of \Cref{cor:ch49-mass-in-lin}.
      It cannot be inferred by taking an arbitrary union of the
      individual polar null sets.
\end{enumerate}
For positive mixed measures, \Cref{prop:ch46-no-mass-outside-lin} is the
corresponding finite-energy integrability statement.
These facts make differences of psh constituents meaningful almost
everywhere; they do not require a common full-measure set where
every possible psh test is finite.
\end{remark}

\subsection{Homogeneity of the dual energy}
\label{sec:ch47-dual-scaling}

\begin{proposition}
\label{prop:ch47-dual-scaling}
Let \(t>0\) and let \(t_\star\) act on \(\mathcal M\) by pushforward along
\eqref{eq:valuation-scaling}.  Then for every \(\mu\in\mathcal M\)
\[
 E^\vee_\omega(t_\star\mu)=t\,E^\vee_\omega(\mu),
\tag{47.16}\label{eq:ch47-dual-scaling}
\]
so \(\cM^1(\omega)\) is stable under \(t_\star\) and
\(t_\star\cM^1_C=\cM^1_{tC}\).  Moreover the only \(t_\star\)-invariant
measure of finite dual energy, for any single \(t\ne1\), is
\(\delta_{v_{\triv}}\).
\end{proposition}

\begin{proof}
Since \(\phi\mapsto t\cdot\phi\) is a bijection of \(\cE^1(\omega)\)
(\Cref{prop:ch47-scaling}\,\cref{itm:ch47-scal:e1}), we may substitute
\(\phi=t\cdot\psi\) in \eqref{eq:bj-dual-energy}.  By
\Cref{thm:ch46-energy-homogeneous}, \(E_\omega(t\cdot\psi)=tE_\omega(\psi)\),
and by the change of variables \(v\mapsto tv\),
\[
 \int(t\cdot\psi)\,\rd(t_\star\mu)
 =\int(t\cdot\psi)(tv)\,\rd\mu(v)
 =t\int\psi\,\rd\mu ,
\]
using \((t\cdot\psi)(tv)=t\psi(v)\).  Hence
\[
 E^\vee_\omega(t_\star\mu)
 =\sup_{\psi\in\cE^1(\omega)}
   \Bigl(tE_\omega(\psi)-t\int\psi\,\rd\mu\Bigr)
 =t\,E^\vee_\omega(\mu),
\]
the factor \(t>0\) passing through the supremum by ledger entry~(3) of
\Cref{sec:ch46-ledger}.  Stability and the identity for sublevel sets are
immediate.  If \(t_\star\mu=\mu\) with \(E^\vee_\omega(\mu)<\infty\), then
\eqref{eq:ch47-dual-scaling} gives \(E^\vee_\omega(\mu)=tE^\vee_\omega(\mu)\)
for that \(t\); taking any \(t\ne1\) forces \(E^\vee_\omega(\mu)=0\), so
\(\mu=\delta_{v_{\triv}}\) by \Cref{cor:ch47-dual-zero}.
\end{proof}

\begin{remark}[Comparison with the BFJ measure completion]
\label{rem:ch47-bfj-comparison}
\Cref{def:bj-finite-energy-measure} is the intrinsic BJ22
Legendre class of probability measures on a trivially valued
analytification.  \Cref{chap:bfj-finite-energy-measures} treats a newer
synthetic measure completion
for the nontrivially valued BFJ setting.  The two constructions share a
duality pattern, but neither is being defined by reference to the other.
In particular, potential-side surjectivity below is a separate theorem.

Three differences should be kept in view.  First, the normalization:
\Cref{def:dual-energy-measure} works with measures of mass \(V\) and
\eqref{eq:measure-fenchel-energy} is the unnormalized Fenchel inequality,
whereas \eqref{eq:bj-fenchel-inequality} is normalized; the conversion is
\eqref{eq:ch28-probability-dictionary}.  Second, the logical order:
\Cref{thm:finite-energy-measure-completion} \emph{constructs}
\(\mathcal M^1\) as an abstract completion and then asks which of its
points are Monge--Amp\`ere measures
(\Cref{rem:ch28-completion-not-surjectivity}), whereas
\Cref{def:bj-finite-energy-measure} defines \(\cM^1(\omega)\) by an
inequality and then proves that every one of its points is a
Monge--Amp\`ere measure.  Third, the scaling action:
\Cref{prop:ch47-dual-scaling} has no analogue in
\Cref{chap:bfj-finite-energy-measures}.  The newer synthetic measure-side
theory is \cite{BJ25FiniteEnergy}, and none of it is being attributed
retroactively to \cite{BJ22}.
\end{remark}

\section{Uniform differentiability}
\label{sec:ch47-differentiability}

An arbitrary continuous perturbation of a psh function need not be psh.
For an extended-real function \(h\), write \(E^\downarrow(h)\) for the
decreasing continuous extension of the equilibrium energy, as in BJ22
Definition~8.2.  When the envelope property holds,

\[
                       E^\downarrow(h)=E(P_\omega(h)).
\tag{47.17}\label{eq:bj-equilibrium-energy}
\]

Since \(X\) is unibranch, \eqref{eq:bj-equilibrium-energy} is available
throughout this chapter, and we use the right-hand side as the working
definition.  The functional \(E^\downarrow\) is the trivially valued
counterpart of the functional \(F\) of
\Cref{def:ch27-envelope-energy}, and
\Cref{prop:ch47-F-elementary} below is the counterpart of
\Cref{prop:ch27-F-elementary}; the proofs are the same, the difference
being that here \(P_\omega\) is known to take values in \(\PSH(\omega)\)
whereas there the corresponding statement is a hypothesis.

\subsection{The equilibrium energy}
\label{sec:ch47-equilibrium}

\begin{proposition}
\label{prop:ch47-F-elementary}
Let \(f,g\) be admissible obstacles for \(\omega\) in the sense of
\Cref{def:ch45-admissible-obstacle}, and let \(c\in\mathbf R\),
\(t\in[0,1]\).  Then:
\begin{enumerate}
\item\label{itm:ch47-F:mono}
      if \(f\le g\), then \(E^\downarrow(f)\le E^\downarrow(g)\);
\item\label{itm:ch47-F:transl}
      \(E^\downarrow(f+c)=E^\downarrow(f)+c\);
\item\label{itm:ch47-F:proj}
      \(E^\downarrow(\phi)=E_\omega(\phi)\) for every
      \(\phi\in\PSH(\omega)\);
\item\label{itm:ch47-F:bounds}
      if \(f\) is real valued and bounded, then \(P_\omega(f)\) is bounded
      with \(\inf f\le P_\omega(f)\le\sup f\); in particular
      \(P_\omega(f)\in\cE^1(\omega)\) and
      \(\inf f\le E^\downarrow(f)\le\sup f\);
\item\label{itm:ch47-F:lip}
      if \(f,g\) are real valued and bounded, then
      \(|E^\downarrow(f)-E^\downarrow(g)|\le\lVert f-g\rVert_\infty\);
\item\label{itm:ch47-F:concave}
      \(E^\downarrow((1-t)f+tg)\ge(1-t)E^\downarrow(f)+tE^\downarrow(g)\);
\item\label{itm:ch47-F:e1}
      if \(\phi\in\cE^1(\omega)\) and \(f\in C^0(X^{\An})\), then
      \(\phi+f\) is an admissible obstacle and
      \(P_\omega(\phi+f)\in\cE^1(\omega)\), with
      \(\phi-\lVert f\rVert_\infty\le P_\omega(\phi+f)\le\phi
      +\lVert f\rVert_\infty\).
\end{enumerate}
\end{proposition}

\begin{proof}
\cref{itm:ch47-F:mono} combines
\Cref{prop:bj-final-obstacle-calculus}\,(1) with monotonicity of
\(E_\omega\) on all of \(\PSH(\omega)\)
(\Cref{prop:ch46-energy-usc}\,\cref{itm:ch46-Eusc:props}), which is what is
needed since \(P_\omega(f)\) may have infinite energy for a general
admissible obstacle.

\cref{itm:ch47-F:transl} combines
\Cref{prop:bj-final-obstacle-calculus}\,(2) with
\eqref{eq:bj-energy-basic-identities}.

\cref{itm:ch47-F:proj} is
\Cref{prop:ch45-obstacle-calculus-2}\,\cref{itm:ch45-P:projection}.

\cref{itm:ch47-F:bounds}.  Since \(\omega\) is ample, hence nef, the
constants are \(\omega\)-psh (\Cref{lem:ch34-psh-nonempty}), so
\(\inf f\) is a competitor for \(f\) and \(P_\omega(f)\ge\inf f\); and
\(P_\omega(f)\le f\le\sup f\).  A bounded \(\omega\)-psh function lies in
\(\cE^1(\omega)\) by \Cref{thm:bj-e1-basic}, and monotonicity of
\(E_\omega\) with \(E_\omega(c)=c\) gives the bounds on \(E^\downarrow\).

\cref{itm:ch47-F:lip} is
\Cref{prop:ch45-obstacle-calculus-2}\,\cref{itm:ch45-P:lipschitz} followed
by \Cref{prop:ch46-cocycle}\,\cref{itm:ch46-coc:lipschitz}.

\cref{itm:ch47-F:concave}.  The function
\(u:=(1-t)P_\omega(f)+tP_\omega(g)\) is \(\omega\)-psh by
\Cref{thm:bj-basic-closure}\,(3) and satisfies \(u\le(1-t)f+tg\), so it is
a competitor and \(u\le P_\omega((1-t)f+tg)\).  Hence, by monotonicity and
then concavity of \(E_\omega\)
(\Cref{prop:ch46-energy-usc}\,\cref{itm:ch46-Eusc:props}),
\[
 E^\downarrow((1-t)f+tg)\ \ge\ E_\omega(u)\ \ge\
 (1-t)E_\omega(P_\omega(f))+tE_\omega(P_\omega(g)),
\]
which is the assertion.

\cref{itm:ch47-F:e1}.  The function \(\phi+f\) is upper semicontinuous
(a sum of an upper semicontinuous and a continuous function) and bounded
above by \(\sup\phi+\lVert f\rVert_\infty\); and
\(\phi-\lVert f\rVert_\infty\) is \(\omega\)-psh and \(\le\phi+f\), so the
competitor family is nonempty and \(\phi+f\) is admissible.  The stated
lower bound follows, and the upper bound is
\(P_\omega(\phi+f)\le\phi+f\le\phi+\lVert f\rVert_\infty\).  Finally
\(\cE^1(\omega)\) is upward closed (\Cref{thm:bj-e1-basic}), and
\(\phi-\lVert f\rVert_\infty\in\cE^1(\omega)\), so
\(P_\omega(\phi+f)\in\cE^1(\omega)\).
\end{proof}

\subsection{The exact one-sided estimate}
\label{sec:ch47-one-sided}

The following lemma is elementary, exact --- there is no error term --- and
is used four times below.  It is the half of
\eqref{eq:bj-uniform-differentiability} that requires no import.

\begin{lemma}
\label{lem:ch47-one-sided-tangent}
Let \(\phi\in\cE^1(\omega)\), \(f\in C^0(X^{\An})\), and
\(\varepsilon\in\mathbf R\).  Then

\[
 E^\downarrow(\phi+\varepsilon f)\ \le\
 E_\omega(\phi)+\varepsilon\int f\,\MA_\omega(\phi).
\tag{47.18}\label{eq:ch47-one-sided-tangent}
\]

In particular \(\varepsilon\mapsto E^\downarrow(\phi+\varepsilon f)\) is a
concave function of \(\varepsilon\) whose graph lies below the line through
\((0,E_\omega(\phi))\) of slope \(\int f\,\MA_\omega(\phi)\), and it touches
that line at \(\varepsilon=0\).
\end{lemma}

\begin{proof}
Put \(u:=P_\omega(\phi+\varepsilon f)\).  By
\Cref{prop:ch47-F-elementary}\,\cref{itm:ch47-F:e1} the obstacle is
admissible and \(u\in\cE^1(\omega)\).  Both \(u\) and \(\phi\) lie in
\(\cE^1(\omega)\), so the right-hand tangent inequality of
\eqref{eq:bj-energy-tangent-inequalities} applies to the pair
\((u,\phi)\):
\[
 E_\omega(u)-E_\omega(\phi)\ \le\ \int(u-\phi)\,\MA_\omega(\phi).
\]
Now \(u\le\phi+\varepsilon f\), so \(u-\phi\le\varepsilon f\) pointwise on
\(X^{\lin}\), where \(\phi\) is finite; since \(\MA_\omega(\phi)\) is a
positive measure carrying no mass outside that locus
(\Cref{prop:ch46-no-mass-outside-lin}), the integral is at most
\(\varepsilon\int f\,\MA_\omega(\phi)\).  Combining gives
\eqref{eq:ch47-one-sided-tangent}.  Concavity in \(\varepsilon\) is
\Cref{prop:ch47-F-elementary}\,\cref{itm:ch47-F:concave} applied to the
affine family \(\varepsilon\mapsto\phi+\varepsilon f\), and the value at
\(\varepsilon=0\) is \(E_\omega(\phi)\) by
\cref{itm:ch47-F:proj}.
\end{proof}

\begin{remark}[What the one-sided estimate does and does not give]
\label{rem:ch47-one-sided-scope}
Applying \eqref{eq:ch47-one-sided-tangent} with \(\varepsilon>0\) and with
\(\varepsilon<0\) shows that, writing
\(G(\varepsilon):=E^\downarrow(\phi+\varepsilon f)\),
\[
 G'(0^+)\ \le\ \int f\,\MA_\omega(\phi)\ \le\ G'(0^-),
\]
the one-sided derivatives existing because \(G\) is concave and finite.
Since concavity already forces \(G'(0^+)\le G'(0^-)\), the two inequalities
are consistent with a corner, and \emph{no} differentiability follows.
This is not a defect of the argument: differentiability of the equilibrium
energy is genuinely equivalent to orthogonality, exactly as in
\Cref{part:bfj}, where \Cref{thm:ch27-superdifferential} computes the
superdifferential of \(F\) and
\Cref{prop:ch27-orthogonality-equivalences} lists the equivalent forms of
orthogonality.  The content of
\Cref{thm:bj-energy-differentiability} is that here the corner is absent,
and the input that removes it is
\Cref{imp:ch47-imported}\,\cref{itm:ch47-import:diff}.
\end{remark}

\subsection{Uniform differentiability}
\label{sec:ch47-uniform-diff}

\begin{theorem}\label{thm:bj-energy-differentiability}
For \(f\in\PL(X)_{\mathbf R}\) and \(\phi\in\mathcal E^1\),

\[
 E^\downarrow(\phi+\varepsilon f)
 =E(\phi)+\varepsilon\int f\MA(\phi)+O(\varepsilon^2),
\tag{47.19}\label{eq:bj-uniform-differentiability}
\]

where the implicit constant is independent of \(\phi\).  Consequently, for
all \(f\in C^0(X^{\An})\),

\[
 \left.\frac{\rd}{\rd t}\right|_{t=0}E^\downarrow(\phi+tf)
                =\int f\MA(\phi).
\tag{47.20}\label{eq:bj-first-variation}
\]
\end{theorem}

\begin{proof}[Proof architecture]
First take \(\phi\) continuous and \(f=\psi_1-\psi_2\), with the
\(\psi_i\) PL psh in one auxiliary ample class.  The determinant-norm
interpretation of relative volume bounds the error by
\(V^{-1}(V_{\omega+\varepsilon\omega'}-V)\varepsilon\|f\|_\infty
=O(\varepsilon^2)\).  Approximate a general \(\phi\in\mathcal E^1\) by a
decreasing net and use continuity of mixed measures.  This proves the PL
case uniformly.  Uniform density of PL functions in \(C^0(X^{\An})\),
monotonicity of equilibrium energy, and a squeeze argument prove
\eqref{eq:bj-first-variation}.  This is
\cite[Theorem~8.5 and Corollary~8.6]{BJ22}.

We now separate what is imported from what is proved, and give the second
half in full.

\emph{Step 1: the upper bound, unconditionally.}  For every
\(f\in C^0(X^{\An})\), every \(\phi\in\cE^1(\omega)\) and every
\(\varepsilon\in\mathbf R\),
\[
 E^\downarrow(\phi+\varepsilon f)\le
 E_\omega(\phi)+\varepsilon\int f\,\MA_\omega(\phi),
\]
by \Cref{lem:ch47-one-sided-tangent}.  No error term and no hypothesis on
\(f\) are needed; in particular the \(O(\varepsilon^2)\) of
\eqref{eq:bj-uniform-differentiability} is one-sided.

\emph{Step 2: the lower bound, imported.}  For \(f\in\PL_{\mathbf R}(X)\)
there is a constant \(C=C(f,\omega)\), independent of \(\phi\), such that
for all \(\phi\in\cE^1(\omega)\) and \(\varepsilon\in[0,1]\),
\[
 E^\downarrow(\phi+\varepsilon f)\ \ge\
 E_\omega(\phi)+\varepsilon\int f\,\MA_\omega(\phi)-C\varepsilon^2 .
\]
This is \cite[Theorem~8.5]{BJ22}, quoted as
\Cref{imp:ch47-imported}\,\cref{itm:ch47-import:diff}.  The mechanism is the
one indicated above: writing \(f=\psi_1-\psi_2\) with \(\psi_i\) PL and
psh for one auxiliary ample class \(\omega'\), the function
\(\phi+\varepsilon f\) is psh for the perturbed class
\(\omega+\varepsilon\omega'\), and the loss incurred in projecting back to
the class \(\omega\) is measured by the difference of the two volumes,
which is \(O(\varepsilon)\); multiplied by the amplitude
\(\varepsilon\lVert f\rVert_\infty\) of the perturbation this gives
\(O(\varepsilon^2)\).  Steps 1 and 2 together are
\eqref{eq:bj-uniform-differentiability}.

\emph{Step 3: from PL to continuous test functions.}  Let
\(f\in C^0(X^{\An})\) and \(\eta>0\).  By
\Cref{thm:ch34-pl-density} choose \(g\in\PL_{\mathbf R}(X)\) with
\(\lVert f-g\rVert_\infty\le\eta\).  Then
\(\phi+\varepsilon g-|\varepsilon|\eta\le\phi+\varepsilon f
\le\phi+\varepsilon g+|\varepsilon|\eta\), so by
\Cref{prop:ch47-F-elementary}\,\cref{itm:ch47-F:mono} and
\cref{itm:ch47-F:transl},
\[
 \bigl|E^\downarrow(\phi+\varepsilon f)-E^\downarrow(\phi+\varepsilon g)\bigr|
 \le|\varepsilon|\,\eta,
\]
while \(\bigl|\int f\MA(\phi)-\int g\MA(\phi)\bigr|\le\eta\) because
\(\MA_\omega(\phi)\) is a probability measure.  Combining with Steps 1 and
2 applied to \(g\), for \(\varepsilon\in(0,1]\),
\[
 -C(g)\varepsilon-2\eta
 \ \le\
 \frac{E^\downarrow(\phi+\varepsilon f)-E_\omega(\phi)}{\varepsilon}
 -\int f\,\MA_\omega(\phi)
 \ \le\ 2\eta .
\]
Letting \(\varepsilon\downarrow0\) and then \(\eta\downarrow0\) gives
\[
 \lim_{\varepsilon\downarrow0}
 \frac{E^\downarrow(\phi+\varepsilon f)-E_\omega(\phi)}{\varepsilon}
 =\int f\,\MA_\omega(\phi).
\]
Replacing \(f\) by \(-f\) gives the same limit from the left, since
\(E^\downarrow(\phi+\varepsilon(-f))\) is the same expression with
\(\varepsilon\) negated and \(\int(-f)\MA(\phi)=-\int f\MA(\phi)\).  Hence
the two-sided derivative exists and equals \(\int f\MA_\omega(\phi)\),
which is \eqref{eq:bj-first-variation}.  This is
\cite[Corollary~8.6]{BJ22}.
\end{proof}

\begin{corollary}
\label{cor:ch47-uniformity}
For \(f\in\PL_{\mathbf R}(X)\) there is \(C=C(f,\omega)\) such that for all
\(\phi\in\cE^1(\omega)\) and all \(\varepsilon\in(0,1]\),
\[
 \int f\,\MA_\omega(\phi)\ \le\
 \frac{E^\downarrow(\phi+\varepsilon f)-E_\omega(\phi)}{\varepsilon}
 +C\varepsilon .
\]
The constant does not depend on \(\phi\).
\end{corollary}

\begin{proof}
Rearrange Step~2 of the proof of
\Cref{thm:bj-energy-differentiability}.
\end{proof}

\begin{remark}[Where uniformity is indispensable]
\label{rem:ch47-uniformity-indispensable}
In \Cref{thm:bj-maximizers-converge} one has a net \((\phi_i)\) and a
parameter \(\varepsilon\), and the conclusion requires letting \(i\) run
first and \(\varepsilon\) tend to zero afterwards.  If the constant in
\eqref{eq:bj-uniform-differentiability} depended on \(\phi_i\), the two
limits could not be exchanged and the argument would collapse.  This is the
precise reason for the word ``uniform'' in the title of this section, and
it is also the reason \Cref{imp:ch47-imported}\,\cref{itm:ch47-import:diff}
is stated with the quantifiers in the order ``there is \(C\), for all
\(\phi\)'' rather than ``for all \(\phi\), there is \(C\)''.  The same
remark applies verbatim to \Cref{lem:ch27-uniform-ma} in
\Cref{part:bfj}, whose role in
\Cref{prop:ch27-orthogonality-uniform-limits} is identical.
\end{remark}

\begin{remark}[Comparison with the two differentiability theorems of \Cref{part:bfj}]
\label{rem:ch47-compare-ch27}
\Cref{chap:bfj-differentiability} contains two differentiability
statements, and it is worth saying which of them corresponds to which
statement here.

\Cref{thm:e1-energy-differentiability} differentiates the energy along a
\emph{segment inside} \(\cE^1\), and its trivially valued counterpart is
\Cref{thm:ch47-segment-expansion}: in both cases the function of the
parameter is a polynomial, the proof is multilinearity, and no envelope and
no envelope property are involved.  That is the ``easy'' differentiability,
and the same polarization argument works in the two legal formalisms.  It
does not identify a compactified pair power with a nonexistent absolute
power of closed model forms over the discretely valued base.

\Cref{thm:ch27-superdifferential} computes the superdifferential of the
\emph{equilibrium} energy \(F\), a functional of an obstacle rather than of
a psh function, and shows that \(F\) is differentiable at \(f\) exactly
when orthogonality holds at \(f\).  Its trivially valued counterpart is
\Cref{thm:bj-energy-differentiability}, and the difference in status is the
whole point: over a discretely valued base orthogonality is a hypothesis
(\Cref{def:orthogonality-property}) verified in special cases
(\Cref{prop:ch27-curve-orthogonality},
\Cref{thm:ch27-legendre-orthogonality}), whereas here it is a theorem,
\Cref{thm:bj-obstacle-orthogonality-energy}, because the envelope property
holds.  A reader who wants the shortest description of what
\Cref{part:bj} buys over \Cref{part:bfj} may take this: the same variational
machine, with a hypothesis replaced by a theorem.
\end{remark}

\section{Maximizers and the variational equation}
\label{sec:ch47-orthogonality}

\begin{definition}[Maximizing net]
\label{def:ch47-maximizing-net}
For \(\mu\in\mathcal M^1\), a net \((\phi_i)\) in \(\mathcal E^1\) is
maximizing if

\[
       E(\phi_i)-\int\phi_i\,\rd \mu\longrightarrow E^\vee(\mu).
\tag{47.21}\label{eq:bj-maximizing-net}
\]
\end{definition}

The use of a net avoids an unnecessary separability hypothesis.  Of course
a maximizing sequence can be chosen from the real supremum, but weakly
convergent subobjects of the compact psh space are naturally subnets.

Throughout this section we abbreviate
\(F_\mu(\phi):=E_\omega(\phi)-\int\phi\,\rd\mu\), so that
\(\sup_{\cE^1(\omega)}F_\mu=E^\vee_\omega(\mu)\) by
\eqref{eq:bj-dual-energy}, and
\(\delta_i:=E^\vee_\omega(\mu)-F_\mu(\phi_i)\ge0\) is the \emph{defect} of
the \(i\)-th term of a maximizing net; thus
\eqref{eq:bj-maximizing-net} says \(\delta_i\to0\).  By
\Cref{prop:ch47-dual-elementary}\,\cref{itm:ch47-dual:nonneg} the functional
\(F_\mu\) is invariant under adding a constant to its argument, so a
maximizing net may always be normalized by \(\sup\phi_i=0\) without
changing its defects.

\subsection{Maximizing nets solve the equation in the limit}
\label{sec:ch47-maximizers-converge}

\begin{theorem}\label{thm:bj-maximizers-converge}
For every maximizing net \((\phi_i)\) for
\(\mu\in\mathcal M^1\),

\[
                         \MA(\phi_i)\longrightarrow\mu
\tag{47.22}\label{eq:bj-maximizer-ma-weak}
\]

weakly as Radon probability measures.
\end{theorem}

\begin{proof}
Let
\(\delta_i=E^\vee(\mu)-E(\phi_i)+\int\phi_i\,\rd \mu\to0\).
For \(f\in\PL(X)_{\mathbf R}\) and \(\varepsilon>0\), duality gives

\[
 E^\downarrow(\phi_i+\varepsilon f)
 -\int(\phi_i+\varepsilon f)\,\rd\mu
 \leq E^\vee(\mu).
\]

Use \eqref{eq:bj-uniform-differentiability}, divide by \(\varepsilon\),
first let \(i\) tend along the net and then let \(\varepsilon\downarrow0\).
This yields
\(\limsup_i\int f\MA(\phi_i)\leq\int f\,\rd \mu\).  Replace \(f\) by
\(-f\) for equality.  PL density in \(C^0\) proves weak convergence.  The
uniformity of the quadratic error is exactly what permits the order of
limits.  This is \cite[Theorem~9.11]{BJ22}.

We give the four steps in detail, since the interchange of limits is the
delicate point and since the argument is used again in
\Cref{thm:ch47-existence-admissible} and in
\Cref{prop:bj-fenchel-equalities}.

\emph{Step 1: the competitor.}  Fix \(f\in\PL_{\mathbf R}(X)\) and
\(\varepsilon\in(0,1]\), and set \(u_i:=P_\omega(\phi_i+\varepsilon f)\).
By \Cref{prop:ch47-F-elementary}\,\cref{itm:ch47-F:e1} we have
\(u_i\in\cE^1(\omega)\), so \eqref{eq:bj-fenchel-inequality} applies to it:
\[
 E_\omega(u_i)\ \le\ E^\vee_\omega(\mu)+\int u_i\,\rd\mu .
\]
Since \(u_i\le\phi_i+\varepsilon f\) and \(\mu\) is a positive measure for
which both sides are \(\mu\)-integrable
(\Cref{prop:ch47-test-class} for \(\phi_i\) and \(u_i\), continuity for
\(f\)),
\[
 \int u_i\,\rd\mu\ \le\ \int\phi_i\,\rd\mu+\varepsilon\int f\,\rd\mu .
\]
Combining and using
\(\int\phi_i\,\rd\mu=E_\omega(\phi_i)-E^\vee_\omega(\mu)+\delta_i\),
\[
 E_\omega(u_i)\ \le\ E_\omega(\phi_i)+\delta_i
 +\varepsilon\int f\,\rd\mu .
\tag{47.22a}\label{eq:ch47-maximizer-upper}
\]

\emph{Step 2: the lower bound, with a constant independent of \(i\).}  By
\Cref{cor:ch47-uniformity} there is \(C=C(f,\omega)\), the same for all
\(i\), with
\[
 E_\omega(u_i)=E^\downarrow(\phi_i+\varepsilon f)
 \ \ge\ E_\omega(\phi_i)+\varepsilon\int f\,\MA_\omega(\phi_i)
 -C\varepsilon^2 .
\]

\emph{Step 3: the interchange.}  Comparing
\eqref{eq:ch47-maximizer-upper} with Step 2, cancelling
\(E_\omega(\phi_i)\) --- legitimate because it is a real number, \(\phi_i\)
having finite energy --- and dividing by \(\varepsilon>0\),
\[
 \int f\,\MA_\omega(\phi_i)\ \le\ \int f\,\rd\mu
 +\frac{\delta_i}{\varepsilon}+C\varepsilon .
\]
Now let \(i\) run along the net: \(\delta_i\to0\), so
\[
 \limsup_i\int f\,\MA_\omega(\phi_i)\ \le\ \int f\,\rd\mu+C\varepsilon .
\]
Only now let \(\varepsilon\downarrow0\); the constant \(C\) does not
depend on \(i\), so the bound survives, and
\(\limsup_i\int f\MA_\omega(\phi_i)\le\int f\,\rd\mu\).  Applying the same
to \(-f\in\PL_{\mathbf R}(X)\) gives
\(\liminf_i\int f\MA_\omega(\phi_i)\ge\int f\,\rd\mu\), whence
\(\int f\,\MA_\omega(\phi_i)\to\int f\,\rd\mu\).

\emph{Step 4: from PL to continuous test functions.}  All the measures
involved are probability measures.  Given \(g\in C^0(X^{\An})\) and
\(\eta>0\), choose \(f\in\PL_{\mathbf R}(X)\) with
\(\lVert f-g\rVert_\infty\le\eta\) (\Cref{thm:ch34-pl-density}); then
\(\bigl|\int g\,\MA(\phi_i)-\int g\,\rd\mu\bigr|
\le2\eta+\bigl|\int f\,\MA(\phi_i)-\int f\,\rd\mu\bigr|\),
so \(\limsup_i\bigl|\int g\MA(\phi_i)-\int g\,\rd\mu\bigr|\le2\eta\) for
every \(\eta>0\).  Hence \(\MA_\omega(\phi_i)\to\mu\) weakly.
\end{proof}

\begin{corollary}
\label{cor:ch47-maximizer-solves}
Let \(\mu\in\cM^1(\omega)\) and suppose \(\phi\in\cE^1(\omega)\) attains
the supremum \eqref{eq:bj-dual-energy}, that is,
\(F_\mu(\phi)=E^\vee_\omega(\mu)\).  Then \(\MA_\omega(\phi)=\mu\).
Conversely, if \(\MA_\omega(\phi)=\mu\), then \(\phi\) attains the
supremum.
\end{corollary}

\begin{proof}
For the first assertion take the constant net with value \(\phi\); it is
maximizing, and \Cref{thm:bj-maximizers-converge} gives
\(\MA_\omega(\phi)=\mu\).  The converse is
\eqref{eq:bj-dual-energy-ma}.  This is the trivially valued form of
\Cref{prop:variational-maximizer-solves}; there the corresponding step
needs the orthogonality property as a hypothesis, and the hypothesis is
supplied here by \Cref{thm:bj-obstacle-orthogonality-energy}, itself a
consequence of the envelope theorem.
\end{proof}

\subsection{Orthogonality}
\label{sec:ch47-orthogonality-proper}

\begin{theorem}
\label{thm:ch47-orthogonality-continuous}
Let \(\omega\in\Amp(X)\) and \(h\in C^0(X^{\An})\), and put
\(u:=P_\omega(h)\).  Then \(u\in C^0(X^{\An})\cap\cE^1(\omega)\), the
contact set \(\{u=h\}\) is closed and nonempty, and
\[
 \int_{X^{\An}}(h-u)\,\MA_\omega(u)=0 ,
 \qquad\text{equivalently}\qquad
 \MA_\omega(u)\ \text{is carried by}\ \{u=h\} .
\]
\end{theorem}

\begin{proof}
\emph{Step 1: the envelope is continuous and has finite energy.}  By
\Cref{cor:ch45-continuity-envelopes}, \(u\in C^0(X^{\An})\cap\PSH(\omega)\);
being continuous on a compact space it is bounded, so
\(u\in\cE^1(\omega)\) by \Cref{thm:bj-e1-basic}.  Hence \(h-u\) is a
continuous, nonnegative function.

\emph{Step 2: the envelope is unchanged along the segment towards the
obstacle.}  For \(t\in[0,1]\) put \(h_t:=(1-t)u+th=u+t(h-u)\).  Since
\(h\ge u\) we have \(h_t\ge u\), so \(u\) is a competitor for \(h_t\) and
\(P_\omega(h_t)\ge u\); and \(h_t\le h\), so
\(P_\omega(h_t)\le P_\omega(h)=u\) by
\Cref{prop:bj-final-obstacle-calculus}\,(1).  Therefore
\[
 P_\omega\bigl(u+t(h-u)\bigr)=u,
 \qquad\text{hence}\qquad
 E^\downarrow\bigl(u+t(h-u)\bigr)=E_\omega(u)
 \quad\text{for all }t\in[0,1].
\]

\emph{Step 3: differentiate.}  The function \(f:=h-u\) is continuous, so
\Cref{thm:bj-energy-differentiability}, in the form
\eqref{eq:bj-first-variation} applied at the finite-energy potential
\(u\), gives
\[
 0=\left.\frac{\rd}{\rd t}\right|_{t=0^+}E^\downarrow(u+tf)
 =\int f\,\MA_\omega(u)
 =\int(h-u)\,\MA_\omega(u).
\]
As \(h-u\ge0\) and \(\MA_\omega(u)\ge0\), the integrand vanishes
\(\MA_\omega(u)\)-almost everywhere; and \(\{h=u\}\) is closed because both
functions are continuous, and nonempty because a positive measure cannot be
carried by the empty set.
\end{proof}

\begin{theorem}\label{thm:bj-obstacle-orthogonality-energy}
Assume the envelope property for \(\omega\).  If \(h\) is lower
semicontinuous, bounded below, and \(P_\omega(h)^*\not\equiv+\infty\), then

\[
 P_\omega(h)^*=P_\omega(h)=h
 \quad\text{a.e. for }\MA(P_\omega(h)^*),
\tag{47.23}\label{eq:bj-obstacle-contact}
\]

and hence

\[
 \int(h-P_\omega(h)^*)\MA(P_\omega(h)^*)=0.
\tag{47.24}\label{eq:bj-obstacle-orthogonality}
\]
\end{theorem}

\begin{proof}[Proof architecture]
For continuous \(h\), the envelope of
\((1-t)P(h)+th\) remains \(P(h)\).  Differentiate its energy at \(t=0\)
using \Cref{thm:bj-energy-differentiability}; the nonnegative contact
defect has zero integral.  Approximate a lower semicontinuous obstacle by an
increasing net of continuous functions.  Strong convergence of the
envelopes and the fact that finite-energy measures ignore negligible sets
pass the contact identity to the limit.  This is
\cite[Theorem~13.1]{BJ22}; the envelope property is an explicit hypothesis.

The first half of this architecture is
\Cref{thm:ch47-orthogonality-continuous} and is proved above in full.  We
record precisely what the second half consists of, and what is imported.

\emph{The approximation.}  A lower semicontinuous function \(h\) bounded
below on the compact Hausdorff, hence normal, space \(X^{\An}\) is the
supremum of an upward directed net \((h_j)\) of continuous functions;
each \(h_j\) is an admissible obstacle for \(\omega\), because \(\omega\) is
nef and constants are competitors
(\Cref{lem:ch34-psh-nonempty}).  By
\Cref{thm:ch45-increasing-obstacles}\,\cref{itm:ch45-inc:sup},
\(\sup_jP_\omega(h_j)=\sup\mathcal F_\omega(h)=P_\omega(h)\), the raw
envelope; and by
\cref{itm:ch45-inc:star} its upper semicontinuous regularization
\(P_\omega(h)^*\) is \(\omega\)-psh.  The hypothesis
\(P_\omega(h)^*\not\equiv+\infty\) is exactly the boundedness clause needed
in \Cref{thm:ch45-increasing-obstacles}; note that for a lower
semicontinuous obstacle the raw and regularized envelopes may genuinely
differ, which is why the statement carries the star and why
\Cref{rem:ch45-increasing-asymmetry} exists.

\emph{The limit.}  For each \(j\),
\Cref{thm:ch47-orthogonality-continuous} gives
\(\int(h_j-P_\omega(h_j))\MA_\omega(P_\omega(h_j))=0\).  Passing to the
limit requires two facts that are proved in
\Cref{chap:bj-strong-topologies} and are not available here: that
\(P_\omega(h_j)\to P_\omega(h)^*\) strongly, hence that the measures
\(\MA_\omega(P_\omega(h_j))\) converge, and that the negligible set
\(\{\sup_jP_\omega(h_j)<P_\omega(h)^*\}\) is null for the limiting
finite-energy measure.  The second is
\Cref{thm:bj-small-set-equivalence} together with
\Cref{prop:ch47-no-polar-mass}.  The statement as a whole is
\cite[Theorem~13.1]{BJ22}.
\end{proof}

\begin{remark}[Orthogonality is a theorem here and a hypothesis in \Cref{part:bfj}]
\label{rem:ch47-orthogonality-status}
\Cref{def:orthogonality-property} defines the orthogonality property over a
discretely valued base, and \Cref{prop:ch27-orthogonality-equivalences}
lists its equivalent forms; it is verified there for curves
(\Cref{prop:ch27-curve-orthogonality}) and by a Legendre-transform argument
in a further special case (\Cref{thm:ch27-legendre-orthogonality}), and is
carried as a hypothesis elsewhere, for instance in
\Cref{thm:ch28-existence-admissible}\,\cref{itm:ch28-exist:equation}.  Here
it is \Cref{thm:ch47-orthogonality-continuous}, an unconditional theorem,
and the reason is precisely
\Cref{imp:ch47-envelope}\,\cref{itm:ch47-env:operator}: the envelope
\(P_\omega(h)\) of a continuous obstacle is a continuous psh function, so
that the segment argument of Step~2 stays inside the class where
\Cref{thm:bj-energy-differentiability} applies.  Over a discretely valued
base the corresponding envelope need not be continuous and the segment
argument is unavailable; this is the content of
\Cref{rem:ch27-two-routes}.
\end{remark}

\subsection{Domination and uniqueness}
\label{sec:ch47-domination}

\begin{theorem}
\label{thm:ch47-domination}
Let \(\phi,\psi\in\cE^1(\omega)\) and suppose that \(\phi\le\psi\) holds
\(\MA_\omega(\phi)\)-almost everywhere.  Then \(\phi\le\psi\) everywhere on
\(X^{\An}\).
\end{theorem}

\begin{proof}[Proof architecture and provenance]
This is part of the finite-energy comparison theory of
\cite[\S10]{BJ22}, and together with
\Cref{thm:ch47-uniqueness} it is
\Cref{imp:ch47-imported}\,\cref{itm:ch47-import:uniqueness}; the two
statements are proved there simultaneously.  We record the reduction that
is available here and the point at which it stops.

Let \(v:=\max\{\phi,\psi\}\in\cE^1(\omega)\), which lies in
\(\cE^1(\omega)\) because \(v\ge\phi\) and \(\cE^1(\omega)\) is upward
closed (\Cref{thm:bj-e1-basic}).  By locality
\eqref{eq:bj-locality},
\[
 \mathbf1_{\{\phi>\psi\}}\MA_\omega(v)
 =\mathbf1_{\{\phi>\psi\}}\MA_\omega(\phi)=0 ,
\]
the second equality by hypothesis; so \(\MA_\omega(v)\) is carried by
\(\{\phi\le\psi\}\), where \(v=\psi\).  Writing \(\mu:=\MA_\omega(\phi)\),
the same hypothesis gives \(\int v\,\rd\mu=\int\psi\,\rd\mu\).  What one
would like to conclude is \(\MA_\omega(v)=\MA_\omega(\psi)\), for then
\Cref{thm:ch47-uniqueness} would give \(v=\psi+c\), and evaluating on
\(\{\psi>\phi\}\), where \(v=\psi\), would give \(c=0\) and hence
\(\phi\le\psi\).  The missing step is the behaviour of \(\MA_\omega(v)\) on
the contact set \(\{\phi=\psi\}\), which locality does not control
(\Cref{cor:ch46-contact-mass} bounds it only from below).
\end{proof}

\begin{remark}[Why the classical proof of the domination principle fails here]
\label{rem:ch47-no-nondegenerate-measure}
In the complex setting, and over a discretely valued base after passing to
a model, the domination principle is proved by comparing against a measure
of full support.  One fixes a bounded psh \(\rho\) whose Monge--Amp\`ere
measure dominates a fixed nondegenerate measure, applies the comparison
principle \eqref{eq:bj-comparison-principle} to the pair
\(\bigl((1-\delta)\psi+\delta\rho+\varepsilon,\ \phi\bigr)\), and concludes
that the open set \(\{\phi>(1-\delta)\psi+\delta\rho+\varepsilon\}\) is
null for a measure of full support, hence empty.

Over a trivially valued field this argument collapses at the first step.
By \Cref{lem:ch47-three-normalizations}\,\cref{itm:ch47-norm:trivial-measure}
the reference measure is
\(\MA_\omega(0)=\delta_{v_{\triv}}\), a single atom; and more generally the
Monge--Amp\`ere measure of a \emph{bounded} psh function is, by
\Cref{prop:ch46-pl-measure} and monotone approximation, a limit of finite
atomic measures on \(X^{\divis}\).  No \(\omega\)-psh function has a
Monge--Amp\`ere measure of full support on \(X^{\An}\): the space is a cone
whose generic points are not in the closure of the divisorial locus in any
measure-theoretic sense useful here.  Consequently the conclusion ``a null
open set is empty'' is unavailable, and the domination principle must be
proved by a different route --- in \cite{BJ22} through the strong topology
and the quantitative estimates of its \S10.  This is the single sharpest
structural difference between the two regimes recorded in
\Cref{sec:ch47-part3}, and it is why
\Cref{thm:ch47-domination} is imported.
\end{remark}

\begin{theorem}
\label{thm:ch47-uniqueness}
Let \(\phi,\psi\in\cE^1(\omega)\).  The following are equivalent:
\begin{enumerate}
\item\label{itm:ch47-uniq:ma}
      \(\MA_\omega(\phi)=\MA_\omega(\psi)\);
\item\label{itm:ch47-uniq:I}
      \(I(\phi,\psi)=0\);
\item\label{itm:ch47-uniq:const}
      \(\phi-\psi\) is constant.
\end{enumerate}
In particular, for every \(\mu\in\cM^1(\omega)\) the equation
\(\MA_\omega(\phi)=\mu\) has at most one solution in \(\cE^1(\omega)\)
modulo constants, and the maximizer of \(F_\mu\) is unique modulo
constants.
\end{theorem}

\begin{proof}
\cref{itm:ch47-uniq:const}\(\Rightarrow\)\cref{itm:ch47-uniq:ma} is
\(\MA_\omega(\phi+c)=\MA_\omega(\phi)\).

\cref{itm:ch47-uniq:ma}\(\Rightarrow\)\cref{itm:ch47-uniq:I}.  Put
\(\mu:=\MA_\omega(\phi)=\MA_\omega(\psi)\).  By
\eqref{eq:bj-dual-energy-ma} both \(\phi\) and \(\psi\) attain the
supremum \(E^\vee_\omega(\mu)\).  Apply \Cref{thm:bj-uniform-concavity}
with \(t=\tfrac12\): since \(\mu\mapsto\int\cdot\,\rd\mu\) is affine,
\[
 F_\mu\Bigl(\frac{\phi+\psi}{2}\Bigr)
 \ \ge\ \frac{F_\mu(\phi)+F_\mu(\psi)}{2}
        +\frac{c_n}{4}I(\phi,\psi)
 \ =\ E^\vee_\omega(\mu)+\frac{c_n}{4}I(\phi,\psi).
\]
But \(F_\mu\le E^\vee_\omega(\mu)\) everywhere on \(\cE^1(\omega)\), so
\(I(\phi,\psi)\le0\), hence \(=0\).

\cref{itm:ch47-uniq:I}\(\Rightarrow\)\cref{itm:ch47-uniq:const} is the last
clause of \Cref{prop:bj-energy-I-identities}, that is,
\cite[Corollary~10.4]{BJ22}; for \(n=1\) it is proved in
\Cref{prop:ch47-curve-uniqueness}.  The proof of
\Cref{prop:bj-energy-I-identities} shows unconditionally that
\cref{itm:ch47-uniq:I}\(\Rightarrow\)\cref{itm:ch47-uniq:ma}, so the three
statements form a cycle in which exactly one implication is imported.

For the last sentence: two solutions of \(\MA_\omega(\phi)=\mu\) differ by
a constant by the equivalence just proved; and two maximizers of \(F_\mu\)
have the same Monge--Amp\`ere measure \(\mu\) by
\Cref{cor:ch47-maximizer-solves}, so the same applies.  Alternatively, the
displayed uniform-concavity computation applies verbatim to two maximizers
and gives \(I=0\) directly.
\end{proof}

\begin{corollary}
\label{cor:ch47-strict-monotone}
Let \(\phi,\psi\in\cE^1(\omega)\) with \(\phi\le\psi\) and
\(E_\omega(\phi)=E_\omega(\psi)\).  Then \(\phi=\psi\).
\end{corollary}

\begin{proof}
By \Cref{cor:ch47-monotone-I}, \(I(\phi,\psi)\le(n+1)(E(\psi)-E(\phi))=0\),
so \(\phi-\psi\) is a constant by \Cref{thm:ch47-uniqueness}; that constant
is \(E_\omega(\phi)-E_\omega(\psi)=0\).
\end{proof}

\subsection{A worked orthogonality computation}
\label{sec:ch47-orthogonality-example}

\begin{example}[An obstacle, its envelope, and the contact set]
\label{ex:ch47-orthogonality-p1}
Let \(X=\mathbf P^1_k\), \(\omega=c_1(\cO(1))\), so \(n=1\), \(V=1\), and
fix a closed point \(p\in X(k)\).  Define a PL obstacle by
\[
 \begin{gathered}
 h:=2\max\{\log|\mathfrak m_p|,-1\},\\
 \text{that is,}\qquad
 h(t\,\ord_p)=-2\min\{t,1\},
 \qquad h=0\ \text{off the }p\text{-ray}.
 \end{gathered}
\]
In the notation of \Cref{ex:ch46-p1-energy} this is the function
``\(\phi_c\) with \(c=2\)'', and since the \(\omega\)-psh range there is
\(0\le c\le d=1\), the function \(h\) is \emph{not} \(\omega\)-psh: its
outgoing slope at \(v_{\triv}\) is \(-2<-1\).

\emph{The envelope.}  We claim
\[
 u:=P_\omega(h)
 \quad\text{satisfies}\quad
 u(t\,\ord_p)=\max\{-1-t,\,-2\},\qquad u\equiv-1\ \text{off the }p\text{-ray}.
\]
That the displayed function is \(\omega\)-psh follows from the curve
criterion of \Cref{ex:ch34-psh-check}: it is convex and nonincreasing on
every ray, and the sum of its outgoing slopes at \(v_{\triv}\) is \(-1\),
which is \(\ge-\deg\omega=-1\).  That it is \(\le h\) is checked on the
\(p\)-ray, where \(-1-t\le-2t\) for \(t\le1\) with equality at \(t=1\), and
off it, where \(-1\le0\).  For maximality, let \(\psi\in\PSH(\omega)\) with
\(\psi\le h\), and set \(a:=\psi(v_{\triv})=\sup\psi\)
(\Cref{lem:ch47-three-normalizations}\,\cref{itm:ch47-norm:max}).  On the
\(p\)-ray, \(\psi\) is convex; from \(\psi(t)\le h(t)=-2t\) for
\(t\in[0,1]\) and \(\psi(t)\ge a-t\) --- the outgoing slope at
\(v_{\triv}\) being \(\ge-1\) --- one gets \(a\le-t\) for every
\(t\in[0,1]\), hence \(a\le-1\); and \(\psi(1)\le-2\).  Convexity on
\([0,1]\) now gives
\[
 \psi(t\,\ord_p)\le(1-t)\psi(v_{\triv})+t\,\psi(\ord_p)
 \le-(1-t)-2t=-1-t
 \qquad(0\le t\le1),
\]
while for \(t\ge1\) monotonicity gives \(\psi\le\psi(\ord_p)\le-2\), and
off the \(p\)-ray \(\psi\le a\le-1\).  So \(\psi\le u\), and
\(u=P_\omega(h)\).

\emph{The measure and the contact set.}  The outgoing slopes of \(u\) are
\(-1\) at \(v_{\triv}\) along the \(p\)-ray and \(0\) elsewhere, and the
graph Laplacian has an atom \(+1\) at \(\ord_p\); so by
\eqref{eq:bj-curve-laplacian}
\[
 \MA^{\mathrm u}_\omega(u)=\bigl(1-1\bigr)\delta_{v_{\triv}}
 +1\cdot\delta_{\ord_p}=\delta_{\ord_p}
 =\MA_\omega(u),
\]
a probability measure since \(V=1\).  The contact set is
\(\{u=h\}=\{t\,\ord_p\mid t\ge1\}\), and it does carry the whole measure.
Orthogonality is now visible:
\[
 \int(h-u)\,\MA_\omega(u)=(h-u)(\ord_p)=(-2)-(-2)=0 ,
\]
as \Cref{thm:ch47-orthogonality-continuous} requires.  Note that
\(h-u>0\) on a large set --- for instance \((h-u)(\tfrac12\ord_p)
=-1-(-\tfrac32)=\tfrac12\) --- so the vanishing of the integral is a
genuine cancellation of support, not of the integrand.

\emph{Energies.}  By
\Cref{prop:ch46-dirichlet-vertical}\,\cref{itm:ch46-dir:energy},
\[
 E_\omega(u)=E^{\mathrm u}_\omega(u)
 =u(v_{\triv})-\tfrac12\int_0^\infty|u'|^2\,\rd t
 =-1-\tfrac12=-\tfrac32,
\]
which agrees with
\(E_\omega(u)=\tfrac12\bigl[\int u\,\MA^{\mathrm u}(0)+\int u\,\MA^{\mathrm u}(u)\bigr]
=\tfrac12[(-1)+(-2)]=-\tfrac32\).  Hence
\(E^\downarrow(h)=-\tfrac32\), whereas \(E_\omega\) evaluated on the
obstacle makes no sense, \(h\) not being psh.  Note also that
\(\sup h=0\) while \(\sup u=-1\): the operator \(P_\omega\) does not
preserve suprema, only the Lipschitz bound
\(\lVert P_\omega(h)-P_\omega(h')\rVert_\infty\le\lVert h-h'\rVert_\infty\)
of \Cref{prop:ch45-obstacle-calculus-2}\,\cref{itm:ch45-P:lipschitz}.

\emph{The dual energy of the Dirac mass.}  By
\eqref{eq:bj-dual-energy-ma},
\[
 E^\vee_\omega(\delta_{\ord_p})=E_\omega(u)-\int u\,\MA_\omega(u)
 =-\tfrac32-(-2)=\tfrac12 .
\]
Since \(T_\omega(\ord_p)=1\) --- the largest drop of a normalized
\(\omega\)-psh function at \(\ord_p\), attained by \(\phi_1\) of
\Cref{ex:ch46-p1-energy} --- this is consistent with
\eqref{eq:bj-dirac-energy-T}, whose two bounds for \(n=1\) are both
\(T_\omega(\ord_p)/2=\tfrac12\).  Normalizing so that the solution vanishes
at \(\ord_p\) gives the Green function \(\phi_{\ord_p}=u+2\), with
\(\phi_{\ord_p}(v_{\triv})=1\) and \(E_\omega(\phi_{\ord_p})=\tfrac12\).
\end{example}

\section{The non-Archimedean Calabi--Yau theorem}
\label{sec:ch47-cy}

The variational method has two halves.  The first is soft: a functional
that is upper semicontinuous on a compact set attains its maximum, and by
\Cref{cor:ch47-maximizer-solves} a maximizer solves the equation.  The
obstacle is that \(\phi\mapsto\int\phi\,\rd\mu\) is only \emph{upper}
semicontinuous for a general \(\mu\), and it enters \(F_\mu\) with a minus
sign; so an extra hypothesis on \(\mu\) is needed to run the direct method.
\Cref{sec:ch47-coercivity,sec:ch47-direct} carry this out for an explicit
class of measures, and the proof there is complete.  The second half
removes the hypothesis, at the price of the imported completeness statement
\Cref{imp:ch47-imported}\,\cref{itm:ch47-import:completeness}; this is
\Cref{sec:ch47-general-cy}.

\subsection{Coercivity}
\label{sec:ch47-coercivity}

\begin{definition}[Uniformly integrable and continuously paired measures]
\label{def:ch47-admissible-measure}
\index{admissible measure!BJ existence}
Let \(\mu\in\mathcal M\).  We say that \(\mu\) is
\begin{enumerate}
\item\label{itm:ch47-adm:coercive}
      \emph{uniformly integrable} for \(\omega\) if
      \[
       M(\mu):=\sup\Bigl\{-\int\phi\,\rd\mu\ \Bigm|\
        \phi\in\PSH^{\sup}(\omega)\Bigr\}<+\infty;
      \tag{47.25}\label{eq:ch47-uniform-integrability}
      \]
\item\label{itm:ch47-adm:paired}
      \emph{continuously paired} with \(\omega\) if the map
      \(\phi\mapsto\int\phi\,\rd\mu\) is real valued and continuous on
      \(\PSH^{\sup}(\omega)\) for the topology of pointwise convergence on
      \(X^{\divis}\).
\end{enumerate}
A measure with both properties is \emph{admissible} for \(\omega\).
\end{definition}

This is \Cref{def:ch28-admissible-measure} transplanted, with
\(\PSH_0(X,\omega)\) replaced by \(\PSH^{\sup}(\omega)\) and the mass
normalization changed; the two conditions play exactly the same two roles,
and the discussion following \Cref{def:ch28-admissible-measure} applies
word for word.

\begin{lemma}
\label{lem:ch47-T-integrable}
Let \(\mu\in\mathcal M\).  With \(T_\omega\) as in
\eqref{eq:bj-T-valuation}, the function \(T_\omega\) is lower
semicontinuous on \(X^{\An}\), and
\[
 M(\mu)\ \le\ \int_{X^{\An}}T_\omega\,\rd\mu .
\]
In particular, if \(\int T_\omega\,\rd\mu<\infty\), then \(\mu\) is
uniformly integrable for \(\omega\); this holds whenever \(\mu\) is carried
by a set on which \(T_\omega\) is bounded, in particular by a finite subset
of \(X^{\lin}\).
\end{lemma}

\begin{proof}
Each \(\phi\in\PSH(\omega)\) is upper semicontinuous, so
\(v\mapsto\sup_{X^{\An}}\phi-\phi(v)\) is lower semicontinuous; a supremum
of lower semicontinuous functions is lower semicontinuous, so \(T_\omega\)
is, and in particular Borel.  For \(\phi\in\PSH^{\sup}(\omega)\) one has
\(-\phi(v)=\sup\phi-\phi(v)\le T_\omega(v)\) pointwise, whence
\(-\int\phi\,\rd\mu\le\int T_\omega\,\rd\mu\); take the supremum over
\(\phi\).  The last assertion is \Cref{def:bj-linear-growth} together with
\Cref{prop:bj-linear-growth-potential}: a point of \(X^{\lin}\) has
\(T_\omega<\infty\), and a finite set carries a finite maximum.
\end{proof}

\begin{lemma}
\label{lem:ch47-coercivity}
\index{coercivity!BJ variational functional}
Let \(\mu\in\mathcal M\) be uniformly integrable, with constant
\(M(\mu)\).  Then \(\mu\in\cM^1(\omega)\), with
\[
 E^\vee_\omega(\mu)\ \le\ M(\mu),
\]
the functional \(F_\mu\) is real valued on \(\cE^1(\omega)\) and equals
\(-\infty\) on \(\PSH(\omega)\setminus\cE^1(\omega)\), and for every
\(\phi\in\PSH^{\sup}(\omega)\)
\[
 F_\mu(\phi)\ \le\ E_\omega(\phi)+M(\mu)\ =\ -J(\phi)+M(\mu).
\tag{47.26}\label{eq:ch47-coercive-inequality}
\]
Consequently, for every \(c\in\mathbf R\),
\[
 \{\phi\in\PSH^{\sup}(\omega)\mid F_\mu(\phi)\ge c\}
 \ \subseteq\ \{\phi\in\PSH^{\sup}(\omega)\mid J(\phi)\le M(\mu)-c\},
\]
and the set on the left is empty when \(c>M(\mu)\).
\end{lemma}

\begin{proof}
Let \(\phi\in\PSH^{\sup}(\omega)\), so \(\phi\le0\) and
\(-\int\phi\,\rd\mu\in[0,M(\mu)]\); in particular the integral is finite.
Hence \(F_\mu(\phi)=E_\omega(\phi)-\int\phi\,\rd\mu\) is real if
\(E_\omega(\phi)>-\infty\) and equals \(-\infty\) otherwise, and
\eqref{eq:ch47-coercive-inequality} follows, the last equality being
\(E_\omega(\phi)=-J(\phi)\) for \(\sup\phi=0\)
(\Cref{lem:ch47-three-normalizations}\,\cref{itm:ch47-norm:J}).  Every
\(\phi\in\cE^1(\omega)\) differs from a normalized one by a constant, and
\(F_\mu\) is invariant under that change by
\Cref{prop:ch47-dual-elementary}\,\cref{itm:ch47-dual:nonneg}; so \(F_\mu\)
is real valued on \(\cE^1(\omega)\).  Taking the supremum of
\eqref{eq:ch47-coercive-inequality} over normalized \(\phi\) and using
\(E_\omega(\phi)\le E_\omega(0)=0\)
(\Cref{prop:ch46-energy-usc}\,\cref{itm:ch46-Eusc:identities}) gives
\(E^\vee_\omega(\mu)\le M(\mu)\).  The inclusion of sublevel sets is
\eqref{eq:ch47-coercive-inequality} rearranged, and the emptiness
assertion follows from \(J\ge0\).

The argument is that of \Cref{lem:ch28-coercivity}, with
\eqref{eq:ch28-coercive-inequality} replaced by
\eqref{eq:ch47-coercive-inequality}; the only change is bookkeeping, the
factor \(V\) having been absorbed into the probability normalization.
\end{proof}

\begin{lemma}
\label{lem:ch47-finite-divisorial}
Let \(v_1,\dots,v_r\in X^{\divis}\) and \(c_1,\dots,c_r\ge0\) with
\(\sum_jc_j=1\), and put \(\mu:=\sum_jc_j\delta_{v_j}\).  Then \(\mu\) is
admissible for \(\omega\), with
\(M(\mu)\le\max_jT_\omega(v_j)<\infty\).
\end{lemma}

\begin{proof}
Uniform integrability is \Cref{lem:ch47-T-integrable}, since
\(X^{\divis}\subset X^{\lin}\) by
\cite[Proposition~1.28]{BJ22}, recorded after
\Cref{def:bj-linear-growth}, so each \(T_\omega(v_j)\) is finite.
Continuity of the pairing is immediate: by definition the topology on
\(\PSH(\omega)\) is that of pointwise convergence on \(X^{\divis}\)
(\eqref{eq:normalized-psh-space}, \Cref{thm:ch34-psh-compactness}), so
\(\phi\mapsto\int\phi\,\rd\mu=\sum_jc_j\phi(v_j)\) is a finite positive
combination of coordinate functions, hence continuous; and it is real
valued because a psh function is finite at a divisorial point
(\Cref{def:bj-psh}\,(2)).
\end{proof}

\begin{proposition}
\label{prop:ch47-skeletal-admissible}
Let \(K\subset X^{\lin}\) be a compact subset of \(X^{\An}\) such that
\begin{enumerate}
\item\label{itm:ch47-skel:T}
      \(C_K:=\sup_KT_\omega<\infty\);
\item\label{itm:ch47-skel:dense}
      \(K\cap X^{\divis}\) is dense in \(K\);
\item\label{itm:ch47-skel:equi}
      the family \(\{\phi|_K\mid\phi\in\PSH^{\sup}(\omega)\}\) is
      equicontinuous.
\end{enumerate}
Then every \(\mu\in\mathcal M\) carried by \(K\) is admissible for
\(\omega\), with \(M(\mu)\le C_K\).  For an SNC test configuration
\(\cX\) and \(c>0\), the truncated dual complex
\[
 \Delta_{\cX}^{(c)}:=\{\alpha\in\Delta_{\cX}:
   \alpha_j\ge c\ \text{for every active coordinate }j\}
\]
satisfies \cref{itm:ch47-skel:T}--\cref{itm:ch47-skel:equi}.
\end{proposition}

\begin{proof}
Uniform integrability with the stated constant is
\Cref{lem:ch47-T-integrable} together with \cref{itm:ch47-skel:T}.  For the
pairing, let \((\phi_i)\) be a net in \(\PSH^{\sup}(\omega)\) converging to
\(\phi\) pointwise on \(X^{\divis}\).  By \cref{itm:ch47-skel:equi} the
restrictions \(\phi_i|_K\) form an equicontinuous family on the compact
space \(K\); by \cref{itm:ch47-skel:dense} they converge pointwise on a
dense subset of \(K\).  An equicontinuous net on a compact space that
converges pointwise on a dense subset converges uniformly: given
\(\varepsilon>0\), choose by equicontinuity a finite cover of \(K\) by open
sets \(U_1,\dots,U_m\) on each of which every member of the family
oscillates by at most \(\varepsilon\), pick \(w_l\in U_l\cap X^{\divis}\),
and use pointwise convergence at the finitely many \(w_l\).  Hence
\(\sup_K|\phi_i-\phi|\to0\) and
\(\bigl|\int\phi_i\,\rd\mu-\int\phi\,\rd\mu\bigr|\le\sup_K|\phi_i-\phi|
\to0\).  The values are real because \(K\subset X^{\lin}\) and
\eqref{eq:bj-e1-sublinear-growth} bounds \(|\phi|\) there.

For the last assertion, \(\Delta_{\cX}^{(c)}\) is a compact subset of
\(X^{\lin}\) whose rational points are divisorial and dense, and
\(T_\omega\) is bounded on \(\Delta_{\cX}\) by
\eqref{eq:bj-dinfty-homogeneity}.  For \cref{itm:ch47-skel:equi},
\Cref{prop:ch49-dual-complex}\,\cref{itm:ch49-dc:lipschitz} makes the
inclusion \(\Delta_{\cX}\hookrightarrow(X^{\lin},d_{\infty,\omega})\)
Lipschitz, with constant \(C_{\cX}/c\), on the part of each face at
distance \(\ge c\) from the coordinate hyperplanes, hence on
\(\Delta_{\cX}^{(c)}\); since every \(\omega\)-psh function is
\(1\)-Lipschitz for \(d_{\infty,\omega}\) by \eqref{eq:bj-dinfty}, the
restricted family is uniformly Lipschitz there, and on a set where the two
topologies agree this is equicontinuity for the topology of
\(X^{\An}\).
\end{proof}

\begin{warning}{Why the whole dual complex is not claimed}
\label{warn:ch47-skeletal-boundary}
The truncation by \(c>0\) is not cosmetic.  The Lipschitz constant of
\Cref{prop:ch49-dual-complex}\,\cref{itm:ch49-dc:lipschitz} degenerates as
\(c\downarrow0\), and \Cref{rem:ch49-dual-complex-lipschitz} declines to
assert a global constant: were the inclusion Lipschitz up to the boundary
of every face, compactness of \(\Delta_{\cX}\) would make the weak and the
\(d_{\infty,\omega}\) topologies agree on \(\Delta_{\cX}\), which is not
proved in this book.  In every case computed there --- the toric one in
particular, through the closed formula \eqref{eq:ch49-toric-dinfty} --- the
inclusion \emph{is} Lipschitz up to the boundary, and then
\cref{itm:ch47-skel:T}--\cref{itm:ch47-skel:equi} hold for all of
\(\Delta_{\cX}\).  Statements below that quote
\Cref{prop:ch47-skeletal-admissible} for a measure carried by a dual
complex should be read as quoting it for \(\Delta_{\cX}^{(c)}\), or as
carrying that hypothesis.
\end{warning}

\begin{remark}[The trivially valued analogue of \Cref{lem:ch28-skeleton-admissible} is weaker]
\label{rem:ch47-skeletal-weaker}
Over a discretely valued base, \Cref{lem:ch28-skeleton-admissible} says
that a measure supported on \emph{any} dual complex of an SNC model is
admissible, and since such complexes are cofinal in the model system this
covers a large and natural class.  Over a trivially valued field the dual
complexes of test configurations are still available, but they are
\emph{not} cofinal in \(X^{\An}\) in the relevant sense: the retractions
\(X^{\An}\to\Delta_{\cX}\) do not exhaust \(X^{\lin}\), and a measure of
finite dual energy need not be carried by any one of them.  So
\Cref{prop:ch47-skeletal-admissible} is a genuine restriction, and the
general Calabi--Yau theorem of \Cref{sec:ch47-general-cy} is not a formal
consequence of it.  This is the second of the three structural differences
listed in \Cref{sec:ch47-part3}.
\end{remark}

\subsection{Existence by the direct method}
\label{sec:ch47-direct}

\begin{theorem}
\label{thm:ch47-existence-admissible}
\index{maximizer!existence of, BJ}
Let \(\mu\in\mathcal M\) be admissible for \(\omega\) in the sense of
\Cref{def:ch47-admissible-measure}.  Then:
\begin{enumerate}
\item\label{itm:ch47-exist:max}
      \(F_\mu\) attains its supremum on \(\PSH^{\sup}(\omega)\), at some
      \(\phi\in\cE^1_0(\omega)\) with \(E_\omega(\phi)\ge-M(\mu)\),
      equivalently \(J(\phi)\le M(\mu)\);
\item\label{itm:ch47-exist:value}
      the maximum value is
      \(E^\vee_\omega(\mu)\in[0,M(\mu)]\);
\item\label{itm:ch47-exist:equation}
      \(\MA_\omega(\phi)=\mu\);
\item\label{itm:ch47-exist:unique}
      \(\phi\) is the unique maximizer in \(\PSH^{\sup}(\omega)\), and the
      unique solution of \(\MA_\omega(\cdot)=\mu\) in
      \(\cE^1(\omega)\) with \(\sup\phi=0\).
\end{enumerate}
\end{theorem}

\begin{proof}
\emph{Step 1: the competitor set is compact.}  By
\Cref{cor:ch45-compactness}, quoted as
\Cref{imp:ch47-envelope}\,\cref{itm:ch47-env:compact}, the space
\(\PSH^{\sup}(\omega)\) is compact for the topology of pointwise
convergence on \(X^{\divis}\).  It is nonempty, containing \(0\).  This is
the point at which the envelope theorem enters, and it is the only such
point in this proof; before
\Cref{chap:bj-envelope-consequences} the statement was the conditional
\Cref{warn:ch34-compactness-conditional}.

\emph{Step 2: the functional is upper semicontinuous.}  On
\(\PSH^{\sup}(\omega)\) write \(F_\mu=E_\omega+G\) with
\(G(\phi):=-\int\phi\,\rd\mu\).  The summand \(E_\omega\) is upper
semicontinuous with values in \([-\infty,0]\), by
\Cref{prop:ch46-energy-usc}\,\cref{itm:ch46-Eusc:props} and
\cref{itm:ch46-Eusc:identities}.  The summand \(G\) is continuous with
values in \([0,M(\mu)]\), by
\Cref{def:ch47-admissible-measure}\,\cref{itm:ch47-adm:paired} and
\Cref{lem:ch47-coercivity}.  A sum of an upper semicontinuous function and
a real-valued continuous one is upper semicontinuous, and no cancellation
of infinities occurs because \(G\) is bounded.  Note that upper
semicontinuity of \(G\) --- which is all that
\Cref{prop:portmanteau} gives for a general \(\mu\), since \(-\phi\) is
lower semicontinuous --- points the wrong way, and this is exactly why
\cref{itm:ch47-adm:paired} is hypothesized.

\emph{Step 3: attainment.}  An upper semicontinuous function on a nonempty
compact space attains its supremum (\Cref{cor:usc-attains-max}).  Let
\(\phi\) be a maximizer.  Since \(F_\mu(0)=E_\omega(0)-0=0\), we have
\(F_\mu(\phi)\ge0\), so by \eqref{eq:ch47-coercive-inequality}
\(E_\omega(\phi)\ge-M(\mu)\); in particular \(E_\omega(\phi)>-\infty\) and
\(\phi\in\cE^1_0(\omega)\), with \(J(\phi)=-E_\omega(\phi)\le M(\mu)\).
This proves \cref{itm:ch47-exist:max}.

\emph{Step 4: the value.}  By
\Cref{prop:ch47-dual-elementary}\,\cref{itm:ch47-dual:nonneg} the functional
\(F_\mu\) is invariant under adding constants, so its supremum over
\(\PSH^{\sup}(\omega)\) equals its supremum over
\(\PSH(\omega)\); and on \(\PSH(\omega)\setminus\cE^1(\omega)\) it is
\(-\infty\) by \Cref{lem:ch47-coercivity}.  Hence the maximum equals
\(\sup_{\cE^1(\omega)}F_\mu=E^\vee_\omega(\mu)\), which lies in
\([0,M(\mu)]\) by
\Cref{prop:ch47-dual-elementary}\,\cref{itm:ch47-dual:nonneg} and
\Cref{lem:ch47-coercivity}.  This is \cref{itm:ch47-exist:value}.

\emph{Step 5: the equation.}  By Step~4 the maximizer \(\phi\) attains the
supremum \eqref{eq:bj-dual-energy}, so
\Cref{cor:ch47-maximizer-solves} gives \(\MA_\omega(\phi)=\mu\).  This is
\cref{itm:ch47-exist:equation}, and it is the only step that uses
\Cref{imp:ch47-imported}\,\cref{itm:ch47-import:diff}.

\emph{Step 6: uniqueness.}  Two maximizers have the same Monge--Amp\`ere
measure by Step~5, hence differ by a constant by
\Cref{thm:ch47-uniqueness}; the normalization \(\sup=0\) fixes the
constant.  This is \cref{itm:ch47-exist:unique}.
\end{proof}

\begin{remark}[The compactness inputs, attributed]
\label{rem:ch47-compactness-attributed}
Three compactness statements occur in this chapter and they should not be
confused.

\emph{Compactness of \(\PSH^{\sup}(\omega)\)}
(\Cref{cor:ch45-compactness}) is used in Step~1 above and nowhere else in
this proof.  It is a consequence of the envelope theorem
\Cref{thm:bj-final-envelope}, hence of
\Crefrange{chap:bj-envelope-statement}{chap:bj-envelope-consequences}, and
before those chapters it was conjectural.

\emph{Compactness of \(\mathcal M\)} (\Cref{prop:appE-measure-compactness})
is used in
\Cref{prop:ch47-dual-elementary}\,\cref{itm:ch47-dual:compact} to make the
sublevel sets of \(E^\vee_\omega\) weakly compact.  It is elementary, being
Banach--Alaoglu for the compact Hausdorff space \(X^{\An}\).

\emph{Compactness of the energy superlevel sets}
\(\{\phi\in\PSH^{\sup}(\omega)\mid J(\phi)\le C\}\) follows from the first
together with upper semicontinuity of \(E_\omega\); it is what
\eqref{eq:ch47-coercive-inequality} converts the coercivity hypothesis
into, and it is the trivially valued form of
\Cref{cor:ch26-net-superlevel}.  All three are stated for nets; no
sequential extraction is performed anywhere, in accordance with
\Cref{warn:ch45-nets} and \Cref{rem:ch28-net-form}.
\end{remark}

\begin{corollary}
\label{cor:ch47-divisorial-green}
Let \(v\in X^{\divis}\).  Then there is a unique
\(\phi_v\in\cE^1(\omega)\) with
\[
 \MA_\omega(\phi_v)=\delta_v,\qquad \phi_v(v)=0 .
\]
It satisfies \(0\le\sup\phi_v=\phi_v(v_{\triv})\le T_\omega(v)\) and
\(E^\vee_\omega(\delta_v)=E_\omega(\phi_v)\in
[T_\omega(v)/(n+1),\,nT_\omega(v)/(n+1)]\).
\end{corollary}

\begin{proof}
Existence and uniqueness modulo constants are
\Cref{thm:ch47-existence-admissible} applied to the admissible measure
\(\delta_v\) (\Cref{lem:ch47-finite-divisorial}); normalizing by
\(\phi_v(v)=0\) --- legitimate because \(\phi_v(v)\) is finite, \(v\) being
divisorial --- fixes the constant.  For the bounds,
\(\sup\phi_v-\phi_v(v)\le T_\omega(v)\) by
\eqref{eq:bj-T-valuation}, and \(\ge0\) trivially.  Finally, with
\(\psi:=\phi_v-\sup\phi_v\in\cE^1_0(\omega)\),
\eqref{eq:bj-dual-energy-ma} gives
\(E^\vee_\omega(\delta_v)=E_\omega(\phi_v)-\int\phi_v\,\rd\delta_v
=E_\omega(\phi_v)-\phi_v(v)=E_\omega(\phi_v)\), and the two-sided estimate
is \eqref{eq:bj-dirac-energy-T}.
\end{proof}

\subsection{Maximizing nets are Cauchy}
\label{sec:ch47-cauchy}

The hypothesis of \Cref{thm:ch47-existence-admissible} is used only to make
the linear term of \(F_\mu\) continuous.  For an arbitrary
\(\mu\in\cM^1(\omega)\) this fails, and
\Cref{rem:ch28-coercivity-shape} explains why no cheap repair is possible:
the Fenchel inequality gives back exactly the term one is trying to
control.  The route that does work replaces compactness by completeness,
and the following proposition is what makes it available.  It is proved
here in full and has an explicit constant.

\begin{proposition}
\label{prop:ch47-maximizing-cauchy}
Let \(\mu\in\cM^1(\omega)\), and let \(\phi,\psi\in\cE^1(\omega)\) have
defects \(\delta_\phi:=E^\vee_\omega(\mu)-F_\mu(\phi)\) and
\(\delta_\psi:=E^\vee_\omega(\mu)-F_\mu(\psi)\).  Then
\[
 I(\phi,\psi)\ \le\ \frac{2}{c_n}\bigl(\delta_\phi+\delta_\psi\bigr)
 =\frac{2(n+1)(n+2)}{n}\bigl(\delta_\phi+\delta_\psi\bigr).
\]
Consequently every maximizing net for \(\mu\) is Cauchy for \(I\), and any
two maximizing nets for \(\mu\) have the same limit behaviour: if
\(\phi_i\) and \(\psi_j\) are maximizing, then
\(I(\phi_i,\psi_j)\to0\) along the product net.
\end{proposition}

\begin{proof}
Put \(\chi:=\tfrac12(\phi+\psi)\in\cE^1(\omega)\)
(\Cref{thm:bj-e1-basic}).  The linear part of \(F_\mu\) is affine, and
\(\int\chi\,\rd\mu=\tfrac12(\int\phi\,\rd\mu+\int\psi\,\rd\mu)\), all three
integrals being finite by \Cref{prop:ch47-test-class}.  Hence by
\Cref{thm:bj-uniform-concavity} with \(t=\tfrac12\),
\[
 F_\mu(\chi)
 =E_\omega(\chi)-\int\chi\,\rd\mu
 \ \ge\ \frac{F_\mu(\phi)+F_\mu(\psi)}{2}+\frac{c_n}{4}I(\phi,\psi).
\]
Since \(F_\mu(\chi)\le E^\vee_\omega(\mu)\) and
\(\tfrac12(F_\mu(\phi)+F_\mu(\psi))
=E^\vee_\omega(\mu)-\tfrac12(\delta_\phi+\delta_\psi)\), we get
\[
 E^\vee_\omega(\mu)\ \ge\ E^\vee_\omega(\mu)
 -\tfrac12(\delta_\phi+\delta_\psi)+\tfrac{c_n}{4}I(\phi,\psi),
\]
whence \(I(\phi,\psi)\le\frac{2}{c_n}(\delta_\phi+\delta_\psi)\).  The
value of \(c_n\) is \Cref{cor:ch47-concavity-constants}.  If
\((\phi_i)\) is maximizing then \(\delta_{\phi_i}\to0\), so
\(I(\phi_i,\phi_{i'})\to0\) along the product net, which is the Cauchy
condition for the quasi-metric \(I\) on \(\cE^1(\omega)/\mathbf R\); the
last assertion is the same computation applied to the two nets.
\end{proof}

\begin{remark}[This is a genuine strict-concavity statement]
\label{rem:ch47-cauchy-is-strict-concavity}
\Cref{prop:ch47-maximizing-cauchy} is the reason
\Cref{thm:bj-uniform-concavity} was proved with a quantitative constant
rather than as bare concavity.  Bare concavity of \(F_\mu\) gives that the
set of maximizers is convex; it does not give that almost-maximizers are
close to each other, and without that one cannot pass from a maximizing net
to a limit.  The situation is the familiar one for uniformly convex
functionals on Banach spaces, and the analogy is exact: \(I\) plays the
role of the squared distance, \eqref{eq:bj-uniform-concavity} is uniform
concavity, and \Cref{prop:ch47-maximizing-cauchy} is the standard
consequence that minimizing sequences of uniformly convex functionals are
Cauchy.  What is \emph{not} standard is that the ambient space is not a
Banach space and that \(I\) is only a quasi-metric; this is why
completeness has to be imported rather than quoted from functional
analysis.
\end{remark}

\subsection{The general theorem}
\label{sec:ch47-general-cy}

\begin{theorem}\label{thm:bj-cy-envelope-equivalence}
For an integral projective \(X\) and \(\omega\in\Amp(X)\), the following
are equivalent:

\begin{enumerate}
\item \(\omega\) has the envelope property;
\item \(\MA:\mathcal E^1(\omega)\to\mathcal M^1(\omega)\) is onto;
\item \(\MA\) induces a bi-Lipschitz isomorphism of the natural
      quasi-metric spaces \(\mathcal E^1(\omega)/\mathbf R\) and
      \(\mathcal M^1(\omega)\);
\item \(\mathcal E^1(\omega)/\mathbf R\), equivalently
      \(\mathcal E^1(\omega)\), is complete for its strong uniformity.
\end{enumerate}
\par
\end{theorem}

\begin{proof}[Proof architecture and dependence]
Assume the envelope property.  Normalize a maximizing sequence by
\(\sup\phi_i=0\).  Compactness of normalized psh functions gives a weakly
convergent subnet.  Energy bounds put the limit in \(\mathcal E^1\), and
upper semicontinuity shows that it attains the dual supremum.
\Cref{prop:bj-fenchel-equalities} then solves the equation.  The energy
estimates identify the potential and measure quasi-metrics.  Conversely,
completeness makes every bounded increasing psh net converge; the envelope
criterion of \Cref{chap:bj-envelope-history} yields the envelope property.
This is exactly
\cite[Theorem~12.8]{BJ22}.

We now give the architecture as a chain of implications, marking at each
step what is proved here and what is imported.

\emph{(1)\(\Rightarrow\)(4).}  Assume the envelope property, so that
\Cref{imp:ch47-envelope} is in force.  Let \((\phi_i)\) be
\(I\)-Cauchy in \(\cE^1(\omega)\), normalized by \(\sup\phi_i=0\).  By
\Cref{cor:ch47-J-vanishes} and the quasi-triangle inequality
\eqref{eq:bj-I-quasitriangle}, the numbers \(I(\phi_i,0)\), hence the
energies \(J(\phi_i)=-E_\omega(\phi_i)\), are bounded, say by \(C\).  The
set \(\{J\le C\}\) is a closed subset of the compact space
\(\PSH^{\sup}(\omega)\) (\Cref{cor:ch45-compactness} and
\Cref{prop:ch46-energy-usc}), so some subnet converges pointwise on
\(X^{\divis}\) to a limit \(\phi\) with \(E_\omega(\phi)\ge-C\), that is,
\(\phi\in\cE^1_0(\omega)\).  Upgrading this to
\(I(\phi_i,\phi)\to0\) is where
\Cref{imp:ch47-imported}\,\cref{itm:ch47-import:completeness} is used: one
replaces the net by the increasing net of envelopes
\(P_\omega\bigl(\inf_{i\ge i_0}\phi_i\bigr)\), which is legitimate because
an infimum of upper semicontinuous functions is upper semicontinuous and
\Cref{imp:ch47-envelope}\,\cref{itm:ch47-env:operator} applies, and then
invokes continuity of the energy along increasing nets,
\cite[Theorem~12.5]{BJ22}.  The full argument is
\cite[Theorem~12.4]{BJ22}, and it is carried out in
\Cref{chap:bj-strong-topologies}, where the strong topology is the subject.

\emph{(4)\(\Rightarrow\)(2).}  Let \(\mu\in\cM^1(\omega)\) and let
\((\phi_i)\) be a maximizing net for \(F_\mu\), normalized by
\(\sup\phi_i=0\); such a net exists because
\(E^\vee_\omega(\mu)=\sup_{\cE^1}F_\mu\) is finite and \(F_\mu\) is
translation invariant.  By \Cref{prop:ch47-maximizing-cauchy} the net is
\(I\)-Cauchy, so by (4) it converges strongly to some
\(\phi\in\cE^1(\omega)\).  Strong convergence implies convergence of
energies and of the pairings against \(\mu\) --- this is the content of
\Cref{thm:bj-strong-measures}, and again belongs to
\Cref{chap:bj-strong-topologies} --- so \(F_\mu(\phi)=E^\vee_\omega(\mu)\)
and \(\phi\) is a maximizer.  \Cref{cor:ch47-maximizer-solves} gives
\(\MA_\omega(\phi)=\mu\).  Only the strong-continuity step is imported; the
Cauchy property, which is the substantive point, is
\Cref{prop:ch47-maximizing-cauchy} and is proved here.

\emph{(2)\(\Rightarrow\)(3).}  Injectivity modulo constants is
\Cref{thm:ch47-uniqueness}.  For the bi-Lipschitz statement, the measure
quasi-metric \(I^\vee\) of \Cref{def:bj-measure-quasimetric} evaluated at
two Monge--Amp\`ere measures is, by
\eqref{eq:bj-dual-energy-ma}, the infimum over \(\chi\) of
\(J_\phi(\chi)+J_\psi(\chi)\), and
\eqref{eq:bj-I-J-comparison} bounds this above and below by dimensional
multiples of \(I(\phi,\psi)\); the details are
\Cref{prop:bj-ma-strong-embedding}.

\emph{(3)\(\Rightarrow\)(4).}  \(\cM^1(\omega)\) is complete for
\(I^\vee\) by \Cref{thm:bj-strong-measures}, a statement that does not
require the envelope property; a bi-Lipschitz image of a complete
quasi-metric space is complete.

\emph{(4)\(\Rightarrow\)(1).}  Completeness makes every increasing net in
\(\cE^1(\omega)\) that is bounded in energy converge, hence makes
regularized suprema psh; by the criterion of
\Cref{chap:bj-envelope-history}, recorded there as
\eqref{eq:continuity-of-envelopes} and its equivalents, this is the
envelope property for \(\omega\).

The whole equivalence is \cite[Theorem~12.8]{BJ22}.  The reader should
notice which single implication carries the analytic weight: it is
(1)\(\Rightarrow\)(4), and it is the one whose proof is deferred to
\Cref{chap:bj-strong-topologies}.  Everything else in the cycle is either
proved above or is formal.
\end{proof}

\begin{corollary}\label{cor:bj-cy-unibranch}
Assume \(X\) is unibranch.  Then for every \(\mu\in\mathcal M^1(\omega)\)
there exists \(\phi\in\mathcal E^1(\omega)\), unique modulo constants, such
that

\[
                              \MA_\omega(\phi)=\mu.
\tag{47.27}\label{eq:bj-cy-equation}
\]
\end{corollary}

\begin{proof}
The theorem proved in
\Crefrange{chap:bj-envelope-statement}{chap:bj-envelope-consequences}
gives the envelope property for every
ample class on a unibranch projective variety.  Apply
\Cref{thm:bj-cy-envelope-equivalence}; uniqueness is
\Cref{prop:bj-energy-I-identities}.

In detail: \Cref{thm:bj-final-envelope} applies to \(X\) integral
projective unibranch and every real class, in particular to
\(\omega\in\Amp(X)\), so hypothesis (1) of
\Cref{thm:bj-cy-envelope-equivalence} holds and (2) gives existence.
Uniqueness modulo constants is \Cref{thm:ch47-uniqueness}, whose proof
reduces it to the vanishing clause of
\Cref{prop:bj-energy-I-identities}.  For the subclass of admissible
measures of \Cref{def:ch47-admissible-measure} --- which by
\Cref{lem:ch47-finite-divisorial} and
\Cref{prop:ch47-skeletal-admissible} already contains every finitely
supported divisorial measure and every measure carried by a dual complex
--- the statement is \Cref{thm:ch47-existence-admissible} and its proof is
complete, with no import beyond
\Cref{imp:ch47-imported}\,\cref{itm:ch47-import:diff}
and~\cref{itm:ch47-import:uniqueness}.
\end{proof}

\begin{corollary}
\label{cor:ch47-scaling-cy}
Assume \(X\) unibranch and let \(t>0\).  If
\(\phi\in\cE^1(\omega)\) solves \(\MA_\omega(\phi)=\mu\), then
\(t\cdot\phi\in\cE^1(\omega)\) solves
\(\MA_\omega(t\cdot\phi)=t_\star\mu\).  Hence the bijection
\[
 \MA_\omega:\ \cE^1(\omega)/\mathbf R\ \xrightarrow{\ \sim\ }\ \cM^1(\omega)
\]
of \Cref{cor:bj-cy-unibranch} intertwines the scaling action on potentials
with the pushforward action on measures, and multiplies the two
quasi-metrics \(I\) and \(I^\vee\) by \(t\).
\end{corollary}

\begin{proof}
The first assertion is \eqref{eq:bj-homogeneous-ma-scaling}.  The two
actions are free away from the fixed points \([0]\) and
\(\delta_{v_{\triv}}\) (\Cref{cor:ch47-J-vanishes},
\Cref{prop:ch47-dual-scaling}), and the scaling of the quasi-metrics is
\Cref{prop:ch47-scaling}\,\cref{itm:ch47-scal:JI} together with
\eqref{eq:ch47-dual-scaling}.
\end{proof}

Historically, \cite[Theorem~12.9]{BJ22} asserted the corollary when \(X\) is smooth
and either \(\operatorname{char}k=0\) or \(\dim X\leq2\), using their then
available envelope theorem.  \Cref{cor:bj-cy-unibranch} is a
post-BJ22 consequence and is labeled as such.

\begin{remark}[What changed, and what did not]
\label{rem:ch47-what-changed}
The variational apparatus of \cite[\S\S8--12]{BJ22} was complete when it
was written; what was missing was the envelope property outside the two
cases just quoted.  Every statement in this chapter whose proof invokes
\Cref{imp:ch47-envelope} was therefore, in 2022, a theorem conditional on
\(\operatorname{EP}_X(\omega)\), and
\Crefrange{chap:bj-envelope-statement}{chap:bj-envelope-consequences}
discharge that condition for unibranch \(X\).  Nothing else changed: no
constant improved, no hypothesis on \(\mu\) was removed, and the imported
statements of \Cref{imp:ch47-imported} are exactly the ones that were
unconditional in \cite{BJ22} to begin with.  A reader comparing this
chapter with the source should therefore expect the same theorems with
smaller hypothesis lists, and \Cref{sec:ch47-hypothesis-use} tabulates the
difference.
\end{remark}

\begin{example}[Dirac data and Green functions]
\label{ex:ch47-dirac-green}
For \(v\in X^{\lin}\), the measure \(\delta_v\) has finite energy
(\Cref{chap:bj-linear-growth}).  On a unibranch \(X\) there is
therefore a unique normalized solution \(\phi_v\in\mathcal E^1\) with
\(\MA(\phi_v)=\delta_v\) and \(\phi_v(v)=0\).  It is continuous and is the
Green function of \(v\): continuity follows from
\cite[Theorem~12.12 and Proposition~12.14(i)]{BJ22}, since
\(\operatorname{supp}\delta_v=\{v\}\subset\{T\leq T(v)\}\).
This does not make \(\phi_v\) PL or a model function, nor does it imply
invariance on any fixed skeleton.  A proper-support point is not in
\(X^{\lin}\), so this statement neither asserts nor suggests that its
Dirac mass has finite energy.

Three refinements are worth recording.  First, when \(v\in X^{\divis}\) the
existence and uniqueness are \Cref{cor:ch47-divisorial-green} and require
no import beyond \Cref{imp:ch47-imported}\,\cref{itm:ch47-import:diff}
and~\cref{itm:ch47-import:uniqueness}; for a general \(v\in X^{\lin}\) they
require \Cref{cor:bj-cy-unibranch}.  Second, the normalization
\(\phi_v(v)=0\) is available exactly because \(v\in X^{\lin}\), where
finite-energy potentials are finite by
\eqref{eq:bj-e1-sublinear-growth}; at a point outside \(X^{\lin}\) the
value could be \(-\infty\) and the normalization would be meaningless.
Third, \(\phi_v\) is in general \emph{not} a \(\mathbf Q\)-PL function even
for divisorial \(v\); \cite[Theorem~A]{BJGreen24} constructs a divisorial
point on a projective variety whose Green function is not \(\mathbf Q\)-PL,
and the warning at the end of \Cref{chap:bj-linear-growth} records the
consequence.  An explicit Green function is computed in
\Cref{ex:ch47-orthogonality-p1} and, in the toric case, in
\eqref{eq:bj-toric-green}.
\end{example}

\section{Worked examples}
\label{sec:ch47-examples}

The four examples below are arranged so that each isolates one assertion of
the chapter and carries it to explicit numbers: a finite-energy potential
with unbounded singularities, a finite-energy measure whose potential is
written down, a toric Calabi--Yau equation solved from the polytope, and a
psh function of infinite energy.  Together they show that all the
inclusions
\[
 \{\text{bounded}\}\ \subsetneq\ \cE^1(\omega)\ \subsetneq\ \PSH(\omega),
 \qquad
 \MA_\omega\bigl(\cE^1(\omega)\bigr)=\cM^1(\omega)\ \subsetneq\
 \mathcal M
\]
are strict, and that the hypotheses of
\Cref{cor:bj-cy-unibranch} are not vacuous.

\subsection{A finite-energy potential with unbounded singularities}
\label{sec:ch47-ex-unbounded}

\begin{example}[A one-parameter family on \(\mathbf P^1\), with all its invariants]
\label{ex:ch47-p1-unbounded}
Let \(X=\mathbf P^1_k\), let \(\omega=c_1(\cO(d))\) with \(d>0\), so
\(n=1\) and \(V=d\), and fix \(p\in X(k)\).  For \(\alpha\in(0,d]\) define
\(\phi_\alpha\) to vanish off the \(p\)-ray and to satisfy
\[
 \phi_\alpha(t\,\ord_p)=h_\alpha(t):=-\alpha\log(1+t),
 \qquad t\ge0 .
\]

\emph{Pshness.}  The function \(h_\alpha\) is convex and decreasing with
\(h_\alpha(0)=0\) and \(h_\alpha'(0^+)=-\alpha\ge-d\), so \(\phi_\alpha\)
is \(\omega\)-psh by the curve criterion of \Cref{ex:ch34-psh-check}.  It
is unbounded below: \(h_\alpha(t)\to-\infty\), so \(\phi_\alpha\) takes the
value \(-\infty\) at the proper-support endpoint of the \(p\)-ray.  Thus
\(\phi_\alpha\) is not bounded, and \(\{\phi_\alpha=-\infty\}\) is a
nonempty one-point set.

\emph{Measure.}  By \eqref{eq:bj-curve-laplacian} the graph Laplacian
contributes \(h_\alpha'(0^+)=-\alpha\) at \(v_{\triv}\) and the density
\(h_\alpha''(t)=\alpha(1+t)^{-2}\) along the ray, so
\[
 \MA^{\mathrm u}_\omega(\phi_\alpha)
 =(d-\alpha)\,\delta_{v_{\triv}}+\frac{\alpha}{(1+t)^2}\,\rd t,
 \qquad
 \MA_\omega(\phi_\alpha)
 =\frac{d-\alpha}{d}\,\delta_{v_{\triv}}
 +\frac{\alpha}{d}\cdot\frac{\rd t}{(1+t)^2},
\]
of total mass \((d-\alpha)+\alpha\int_0^\infty(1+t)^{-2}\rd t
=(d-\alpha)+\alpha=d=V\), as
\Cref{prop:ch46-pl-measure}\,\cref{itm:ch46-pl-ma:mass} requires; it is
positive exactly for \(0\le\alpha\le d\), matching the psh range.  Note
that for \(\alpha=d\) the atom at \(v_{\triv}\) disappears and the whole
mass sits on the open ray: a diffuse Monge--Amp\`ere measure, which no PL
function can produce.

\emph{Energy.}  By
\Cref{prop:ch46-dirichlet-vertical}\,\cref{itm:ch46-dir:energy},
\[
 E^{\mathrm u}_\omega(\phi_\alpha)
 =\int\phi_\alpha\,\MA^{\mathrm u}_\omega(0)
  -\tfrac12\langle\phi_\alpha,\phi_\alpha\rangle_{\mathrm D}
 =h_\alpha(0)-\tfrac12\int_0^\infty|h_\alpha'|^2\,\rd t
 =-\frac{\alpha^2}{2},
\]
since \(\int_0^\infty\alpha^2(1+t)^{-2}\rd t=\alpha^2\); hence
\[
 E_\omega(\phi_\alpha)=-\frac{\alpha^2}{2d}\ >\ -\infty,
 \qquad\text{so}\qquad
 \phi_\alpha\in\cE^1(\omega)\setminus\{\text{bounded}\}.
\]
The same number arises from
\eqref{eq:ch46-pure-pairing-as-integrals}:
\(2V E_\omega=\int\phi_\alpha\MA^{\mathrm u}(0)
+\int\phi_\alpha\MA^{\mathrm u}(\phi_\alpha)
=0+\alpha\int_0^\infty\frac{-\alpha\log(1+t)}{(1+t)^2}\rd t=-\alpha^2\),
using \(\int_0^\infty\log(1+t)(1+t)^{-2}\rd t=1\), which is
\(\int_1^\infty s^{-2}\log s\,\rd s=1\) after \(s=1+t\).

\emph{The functionals.}  Since \(\sup\phi_\alpha=\phi_\alpha(v_{\triv})=0\),
\[
 J(\phi_\alpha)=\frac{\alpha^2}{2d},
 \qquad
 I(\phi_\alpha,0)=\phi_\alpha(v_{\triv})
   -\int\phi_\alpha\,\MA_\omega(\phi_\alpha)
 =0+\frac{\alpha^2}{d}
 =\frac{\alpha^2}{d},
\]
so \(J=\tfrac12I\), which is \eqref{eq:bj-I-J-comparison} in the case
\(n=1\), where both constants equal \(\tfrac12\).  Finally
\eqref{eq:bj-dual-energy-ma} gives
\[
 E^\vee_\omega\bigl(\MA_\omega(\phi_\alpha)\bigr)
 =E_\omega(\phi_\alpha)-\int\phi_\alpha\,\MA_\omega(\phi_\alpha)
 =-\frac{\alpha^2}{2d}+\frac{\alpha^2}{d}
 =\frac{\alpha^2}{2d}\ \in(0,\infty),
\]
so the measure displayed above lies in \(\cM^1(\omega)\).

\emph{Comparison.}  This is the trivially valued counterpart of
\Cref{ex:ch26-unbounded-finite-energy}, which exhibits an unbounded
finite-energy potential over a discretely valued base; there the potential
is unbounded along a sequence of models and here it is unbounded along a
single ray, but the mechanism --- a square-integrable slope with divergent
primitive --- is identical.  The instance \(d=1\), \(\alpha=1\) is
\Cref{ex:ch46-e1-curve}, where \(E_\omega=-\tfrac12\).
\end{example}

\subsection{A finite-energy measure whose potential is explicit}
\label{sec:ch47-ex-measure}

\begin{example}[Solving \(\MA_\omega(\phi)=\mu\) on \(\mathbf P^1\) by hand]
\label{ex:ch47-explicit-potential}
Keep \(X=\mathbf P^1\), \(\omega=c_1(\cO(d))\), and fix
\(\alpha\in(0,d]\).  Take as \emph{data} the probability measure
\[
 \mu_\alpha:=\frac{d-\alpha}{d}\,\delta_{v_{\triv}}
 +\frac{\alpha}{d}\cdot\frac{\rd t}{(1+t)^2}
 \quad\text{on }\{t\,\ord_p\}\cup\{v_{\triv}\},
\]
where \(\rd t\) is the arclength parameter of the \(p\)-ray.  We solve
\(\MA_\omega(\phi)=\mu_\alpha\) directly, and check every assertion of
\Cref{sec:ch47-cy} on it.

\emph{Solving.}  Look for \(\phi\) vanishing off the \(p\)-ray, with
\(\phi(t\,\ord_p)=h(t)\) for a convex decreasing \(h\) with \(h(0)=0\).  By
\eqref{eq:bj-curve-laplacian} the equation
\(\MA^{\mathrm u}_\omega(\phi)=d\mu_\alpha\) reads
\[
 h''(t)=\frac{\alpha}{(1+t)^2}\ \ (t>0),
 \qquad
 d+h'(0^+)=d-\alpha .
\]
The second condition gives \(h'(0^+)=-\alpha\); integrating the first,
\(h'(t)=-\alpha/(1+t)\), and integrating again with \(h(0)=0\),
\[
 h(t)=-\alpha\log(1+t).
\]
So \(\phi=\phi_\alpha\) of \Cref{ex:ch47-p1-unbounded}, and the two
prescriptions --- one datum, one potential --- match.  Uniqueness modulo
constants is \Cref{thm:ch47-uniqueness}; here it is visible, since the two
integrations introduced exactly one constant, fixed by \(h(0)=0\).

\emph{The numbers.}  Collecting from \Cref{ex:ch47-p1-unbounded},
\[
 E_\omega(\phi_\alpha)=-\frac{\alpha^2}{2d},
 \quad
 \int\phi_\alpha\,\rd\mu_\alpha=-\frac{\alpha^2}{d},
 \quad
 E^\vee_\omega(\mu_\alpha)=\frac{\alpha^2}{2d},
 \quad
 J(\phi_\alpha)=\frac{\alpha^2}{2d}.
\]
The Fenchel inequality \eqref{eq:bj-fenchel-inequality} reads
\(-\alpha^2/(2d)\le\alpha^2/(2d)+(-\alpha^2/d)\), an equality, as
\Cref{prop:bj-fenchel-equalities} requires for
\(\mu=\MA_\omega(\phi)\).  The coercivity hypothesis of
\Cref{lem:ch47-coercivity}, on the other hand, \emph{fails} here.  For
\(s>0\) let \(\psi_s:=d\max\{\log|\mathfrak m_p|,-s\}\), so that
\(\psi_s(t\,\ord_p)=-d\min\{t,s\}\); this is \(\omega\)-psh with
\(\sup\psi_s=0\), the outgoing slope at \(v_{\triv}\) being \(-d\).  Then
\[
 \begin{aligned}
 -\int\psi_s\,\rd\mu_\alpha
 &=\frac{\alpha}{d}
   \int_0^\infty\frac{d\min\{t,s\}}{(1+t)^2}\,\rd t\\
 &=\alpha\int_0^\infty
   \frac{\min\{t,s\}}{(1+t)^2}\,\rd t\\
 &\xrightarrow[\ s\to\infty\ ]{}
   \alpha\int_0^\infty\frac{t\,\rd t}{(1+t)^2}
   =+\infty .
 \end{aligned}
\]
the last integral diverging logarithmically.  Hence
\(M(\mu_\alpha)=+\infty\): this \(\mu_\alpha\) is \emph{not} uniformly
integrable, and
\Cref{thm:ch47-existence-admissible} does not apply to it: the general
\Cref{cor:bj-cy-unibranch} is needed, or, as here, an explicit solution.
This is a useful warning: \(\cM^1(\omega)\) is strictly larger than the
class of admissible measures of \Cref{def:ch47-admissible-measure}.

\emph{Truncating to reach the admissible class.}  Replacing
\(\rd t/(1+t)^2\) by its restriction to \([0,R]\), renormalized, gives a
measure carried by the compact set \(\{t\,\ord_p\mid t\le R\}\), on which
\(T_\omega\le dR\); such a measure is admissible by
\Cref{lem:ch47-T-integrable} and
\Cref{prop:ch47-skeletal-admissible}, and the corresponding potential is
obtained by the same two integrations, being affine beyond \(t=R\).  As
\(R\to\infty\) these potentials \emph{decrease} to \(\phi_\alpha\) and
their energies decrease to \(-\alpha^2/(2d)\).  The direction is worth
checking, because renormalization pushes the opposite way to deletion:
removing the tail mass \(\alpha/(1+R)\) and rescaling multiplies
\(\mu^{\mathrm u}([0,t])\) by \((1-\varepsilon)^{-1}>1\), where
\(\varepsilon=\alpha/\bigl(d(1+R)\bigr)\), so that
\(h_R'(t)=-d+\mu^{\mathrm u}_R([0,t])\ge h_\alpha'(t)\) for every \(t\)
while \(h_R(0)=h_\alpha(0)=0\), whence \(\phi_R\ge\phi_\alpha\); and
\(\varepsilon\) decreases in \(R\), so the net decreases.  With
\(d=\alpha=1\) one gets \(h_R'(t)=(t/R-1)/(1+t)\) on \([0,R]\) and \(0\)
beyond, hence \(E_\omega(\phi_R)=-\tfrac12(3-4\log2)\approx-0.114\) at
\(R=1\) and \(\approx-0.463\) at \(R=100\), decreasing to
\(-\tfrac12=-\alpha^2/(2d)\).  \Cref{cor:ch47-monotone-I} applies with
\(\phi=\phi_\alpha\le\psi=\phi_R\) and gives
\(I(\phi_\alpha,\phi_R)\le(n+1)\bigl(E(\phi_R)-E(\phi_\alpha)\bigr)\to0\),
so the convergence is strong.  The example
therefore also illustrates the passage from
\Cref{thm:ch47-existence-admissible} to
\Cref{cor:bj-cy-unibranch}: exhaustion by admissible data, plus
completeness.
\end{example}

\subsection{A toric Calabi--Yau equation solved from the polytope}
\label{sec:ch47-ex-toric}

\begin{example}[Two toric solutions, read off the polytope]
\label{ex:ch47-toric-cy}
We use the toric dictionary of
\Cref{thm:bj-toric-psh-convex,thm:bj-toric-ma-energy}, with the sign
convention \(P_{\mathrm{BJ}}=-P\) of
\eqref{eq:bj-toric-sign-conversion} and
\(V_P=\operatorname{vol}(P_{\mathrm{BJ}})=(\omega^n)/n!\).  Recall that for
a torus-invariant \(\phi\) the equation \(\MA_\omega(\phi)=\mu\) with
\(\mu\) carried by \(N_{\mathbf R}\) becomes, by
\eqref{eq:bj-toric-ma}, the real Monge--Amp\`ere equation
\[
 \operatorname{MA}_{\mathbf R}(f_\phi)=V_P\,\mu
 \qquad\text{on }N_{\mathbf R},
 \qquad f_\phi=(\phi+f_P)|_{N_{\mathbf R}},
\]
subject to \(f_\phi\) convex with \(f_\phi\le f_P+O(1)\), which by
\Cref{thm:bj-toric-psh-convex} is the same as
\(\partial f_\phi(N_{\mathbf R})\subseteq P_{\mathrm{BJ}}\).

\emph{First solution: an atomic datum on \(\mathbf P^2\).}  Take
\(X=\mathbf P^2\), \(\omega=c_1(\cO(1))\), \(n=2\), \(V=1\),
\(V_P=\tfrac12\), and the data of \Cref{ex:ch46-toric-p2}.  For
\(c\in[0,1]\) let
\[
 \mu_c:=c^2\,\delta_{\ord_0}+(1-c^2)\,\delta_{v_{\triv}},
\]
a probability measure carried by two torus-invariant points, namely
\(v_{(1,1)}=\ord_0\) and \(v_{(0,0)}=v_{\triv}\).  By
\Cref{lem:ch47-finite-divisorial} it is admissible, so
\Cref{thm:ch47-existence-admissible} already guarantees a unique normalized
solution; and \Cref{ex:ch46-toric-p2} exhibits it:
\[
 \phi_c=\max\{c\log|\mathfrak m_0|,-c\},
 \qquad
 \MA_\omega(\phi_c)=\mu_c,
 \qquad
 E_\omega(\phi_c)=-\frac{c^3}{3}.
\]
On the polytope side this is the statement that the convex function
\(f_{\phi_c}=f_P+\phi_c\), whose conjugate is
\(f_{\phi_c}^\vee(m)=(c+m_1+m_2)_+\), has Alexandrov measure
\(\tfrac{c^2}{2}\delta_{(1,1)}+\tfrac{1-c^2}{2}\delta_{(0,0)}
=V_P\,\mu_c\); the two atoms are the images of the two subgradient cells,
of volumes \(c^2/2\) and \((1-c^2)/2\).  The dual energy is
\[
 E^\vee_\omega(\mu_c)
 =E_\omega(\phi_c)-\int\phi_c\,\rd\mu_c
 =-\frac{c^3}{3}-c^2\cdot(-c)
 =\frac{2c^3}{3},
\]
and the two comparison constants of \eqref{eq:bj-I-J-comparison} are both
attained here: with \(\psi=0\) one has
\(I(\phi_c,0)=c^3\), \(J_0(\phi_c)=c^3/3=I/(n+1)\) and
\(J_{\phi_c}(0)=2c^3/3=nI/(n+1)\).  This is the sharpness configuration
described in the proof of \Cref{prop:bj-energy-I-identities}.

\emph{Second solution: a diffuse datum on \(\mathbf P^1\times\mathbf P^1\).}
Take \(X=\mathbf P^1\times\mathbf P^1\) and
\(\omega=c_1(\cO(1,1))\), so \(n=2\), \(V=(\omega^2)=2\),
\(P=[0,1]^2\), \(P_{\mathrm{BJ}}=[-1,0]^2\), \(V_P=1=V/2!\), and
\[
 f_P(u)=\max\{0,-u_1\}+\max\{0,-u_2\}.
\]
Let \(\nu\) be the probability measure \((1+u)^{-2}\rd u\) on
\([0,\infty)\subset\mathbf R\) and take as datum the product
\[
 \mu:=\nu\otimes\nu\quad\text{on }N_{\mathbf R}=\mathbf R^2 .
\]
Set \(g(u):=-u\) for \(u\le0\) and \(g(u):=-\log(1+u)\) for \(u\ge0\); then
\(g\) is convex with \(g'(-\infty)=-1\), \(g'(+\infty)=0\), so
\(g'(\mathbf R)=(-1,0)\), and \(g''=\nu\) as a measure.  Put
\[
 f(u):=g(u_1)+g(u_2),
 \qquad
 \phi:=f-f_P\ \text{ on }N_{\mathbf R}.
\]
Then \(\partial f(N_{\mathbf R})=(-1,0)^2\subset P_{\mathrm{BJ}}\), so
\(f\le f_P+O(1)\) --- indeed \(f\le f_P\) with equality on the negative
quadrant --- and \(\phi\) extends uniquely to a torus-invariant
\(\omega\)-psh function by \eqref{eq:bj-toric-psh-inverse}.  Its real
Monge--Amp\`ere measure is the product of the one-dimensional ones,
\[
 \operatorname{MA}_{\mathbf R}(f)=g''(u_1)\,g''(u_2)\,\rd u_1\rd u_2
 =\nu\otimes\nu=V_P\,\mu ,
\]
so \(\MA_\omega(\phi)=\mu\) by \eqref{eq:bj-toric-ma}.  Explicitly,
\[
 \phi(v_u)=-\log(1+u_1^+)-\log(1+u_2^+),
 \qquad u^+:=\max\{u,0\},
\]
which is unbounded below; the solution of a Calabi--Yau equation with a
perfectly smooth datum is therefore not bounded, in sharp contrast with the
complex case.

\emph{The numbers.}  By \eqref{eq:bj-toric-energy} and the product
structure,
\[
 \begin{aligned}
  E_\omega(\phi)
   &=-\frac1{V_P}\int_{P_{\mathrm{BJ}}}f^\vee\,\rd m\\
   &=-\int_{[-1,0]^2}
     \bigl(g^\vee(m_1)+g^\vee(m_2)\bigr)\rd m\\
   &=-2\int_{-1}^0g^\vee(m)\,\rd m=-1,
 \end{aligned}
\]
because the one-dimensional computation of
\Cref{ex:ch47-p1-unbounded} with \(d=\alpha=1\) gives
\(\int_{-1}^0g^\vee=\tfrac12\).  Independently, by
\eqref{eq:ch46-pure-pairing-as-integrals} with \(n=2\), \(V=2\),
\[
 6E_\omega(\phi)=\int\phi\,\omega^2+\int\phi\,(\omega+\ddc\phi)\wedge\omega
 +\int\phi\,(\omega+\ddc\phi)^2
 =0+(-2)+(-4)=-6 ,
\]
where the three measures are, in the toric picture,
\(2(\delta_0\otimes\delta_0)\),
\(\nu\otimes\delta_0+\delta_0\otimes\nu\) and \(2(\nu\otimes\nu)\), and
\(\int\log(1+u)\,\rd\nu=1\).  Finally
\[
 E^\vee_\omega(\mu)=E_\omega(\phi)-\int\phi\,\MA_\omega(\phi)
 =-1-(-2)=1 ,
\]
and \(I(\phi,0)=\phi(v_{\triv})-\int\phi\,\MA_\omega(\phi)=0+2=2\), so
\(J_0(\phi)=1\) and \(J_\phi(0)=1\), both strictly inside the interval
\([I/3,2I/3]=[2/3,4/3]\) of \eqref{eq:bj-I-J-comparison}, as they must be
for a datum that is neither of the two extreme configurations.

\emph{A check against \Cref{part:bfj}.}  The same product construction over
a discretely valued base appears in \Cref{ex:ch28-chain-dual-energy}, and
the numbers there differ from these by the factor \(V\) of
\Cref{warn:ch46-normalization} and by the reflection
\eqref{eq:bj-toric-sign-conversion}.  A reader who obtains
\(E_\omega(\phi)=-2\) here has used the unnormalized energy
\(E^{\mathrm u}_\omega=VE_\omega\); one who obtains \(+1\) has dropped the
minus sign in \eqref{eq:bj-toric-energy}, which comes from the reflection.
\end{example}

\subsection{A psh function of infinite energy}
\label{sec:ch47-ex-infinite}

\begin{example}[\(\cE^1(\omega)\) is a proper subset of \(\PSH(\omega)\), and \(\cM^1(\omega)\) of \(\mathcal M\)]
\label{ex:ch47-infinite-energy}
Let \(X=\mathbf P^1\), \(\omega=c_1(\cO(1))\), so \(d=V=1\), and let
\(\psi\) vanish off the \(p\)-ray with
\[
 \psi(t\,\ord_p)=h(t):=1-\sqrt{1+t},\qquad t\ge0 .
\]

\emph{Pshness.}  \(h\) is decreasing, \(h(0)=0\),
\(h'(t)=-\tfrac12(1+t)^{-1/2}\) so \(h'(0^+)=-\tfrac12\ge-1\), and
\(h''(t)=\tfrac14(1+t)^{-3/2}>0\), so \(h\) is convex and \(\psi\) is
\(\omega\)-psh by \Cref{ex:ch34-psh-check}.  Note that \(\psi\) is
\emph{less} singular than the function \(\log|\mathfrak m_p|\) of
\Cref{ex:ch46-e1-curve}, whose slope is \(-1\) throughout: here the slope
tends to \(0\).

\emph{Infinite energy.}  By the Dirichlet formula of
\Cref{prop:ch46-dirichlet-vertical}\,\cref{itm:ch46-dir:energy},
\[
 E_\omega(\psi)=h(0)-\tfrac12\int_0^\infty|h'|^2\,\rd t
 =-\frac18\int_0^\infty\frac{\rd t}{1+t}=-\infty .
\]
So \(\psi\in\PSH(\omega)\setminus\cE^1(\omega)\), and
\(\cE^1(\omega)\subsetneq\PSH(\omega)\).  The mechanism is that
\(h'\in L^2\) fails only logarithmically; the family
\(h_\beta(t)=1-(1+t)^{\beta}\) with \(\beta\in(0,1)\) has
\(\int|h_\beta'|^2=\beta^2\int(1+t)^{2\beta-2}\rd t<\infty\) exactly for
\(\beta<\tfrac12\), so \(\beta=\tfrac12\) is the borderline and
\(\cE^1(\omega)\) has no ``uniform slope'' description.

\emph{The associated measure has infinite dual energy.}  Consider the
probability measure
\[
 \mu:=\tfrac12\,\delta_{v_{\triv}}
 +\tfrac14\,(1+t)^{-3/2}\,\rd t
\]
on \(\{v_{\triv}\}\cup\{t\ord_p\}\), which is what
\eqref{eq:bj-curve-laplacian} produces from \(\psi\): the atom is
\(1+h'(0^+)=\tfrac12\), the density is \(h''\), and the total mass is
\(\tfrac12+\tfrac12=1\).  We show \(E^\vee_\omega(\mu)=+\infty\).  Let
\(\chi_j:=\max\{\psi,-j\}\), a bounded \(\omega\)-psh function with
\(\sup\chi_j=0\); on the ray it equals \(h\) up to
\(T_j:=(1+j)^2-1\) and \(-j\) afterwards.  Then
\[
 E_\omega(\chi_j)=-\tfrac12\int_0^{T_j}|h'|^2\,\rd t
 =-\tfrac18\log(1+T_j)=-\tfrac14\log(1+j),
\]
while, substituting \(s=\sqrt{1+t}\),
\[
 \begin{aligned}
 \int\chi_j\,\rd\mu
 &=\tfrac14\int_0^{T_j}\bigl(1-\sqrt{1+t}\bigr)(1+t)^{-3/2}\rd t\\
 &\quad{}-\tfrac{j}{4}\int_{T_j}^\infty(1+t)^{-3/2}\rd t\\
 &=\tfrac12\Bigl(1-\tfrac1{1+j}-\log(1+j)\Bigr)
   -\tfrac{j}{2(1+j)},
 \end{aligned}
\]
which simplifies to \(-\tfrac12\log(1+j)\).  Hence
\[
 F_\mu(\chi_j)=E_\omega(\chi_j)-\int\chi_j\,\rd\mu
 =-\tfrac14\log(1+j)+\tfrac12\log(1+j)
 =\tfrac14\log(1+j)\ \longrightarrow\ +\infty ,
\]
so \(E^\vee_\omega(\mu)=+\infty\) and \(\mu\notin\cM^1(\omega)\).  Thus
\(\cM^1(\omega)\subsetneq\mathcal M\), and the hypothesis
\(\mu\in\cM^1(\omega)\) in \Cref{cor:bj-cy-unibranch} is a genuine
restriction: there is no \(\phi\in\cE^1(\omega)\) with
\(\MA_\omega(\phi)=\mu\), by \eqref{eq:bj-dual-energy-ma}.

\emph{What this example is not.}  It does not say that \(\mu\) fails to be
a Monge--Amp\`ere measure in any weaker sense; it is one, of the
infinite-energy potential \(\psi\), by the graph-Laplacian formula
\eqref{eq:bj-curve-laplacian}, which on a curve is available for every psh
function.  What fails is finiteness of the dual energy, and with it every
statement of \Cref{sec:ch47-cy}.  Nor does the example produce a pluripolar
set of positive \(\mu\)-measure: \(\mu\) gives no mass to the single point
\(\{\psi=-\infty\}\), which is the endpoint at infinity of the ray, since
\(\int_0^\infty(1+t)^{-3/2}\rd t<\infty\).  The failure detected by
\(E^\vee_\omega\) is quantitative, not set-theoretic --- which is exactly
the difference between \Cref{prop:ch47-no-polar-mass} and membership in
\(\cM^1(\omega)\).

\emph{Comparison.}  This is the trivially valued counterpart of
\Cref{ex:ch26-infinite-energy} and of
\Cref{ex:ch28-infinite-dual-energy}, which make the two points ---
infinite energy on the potential side, infinite dual energy on the measure
side --- over a discretely valued base.  The comparison is exact, and the
proofs are the same computation with the ray replaced by a decreasing
sequence of models.
\end{example}

\section{Scope and hypotheses}
\label{sec:ch47-scope}

\subsection{Where each hypothesis is used}
\label{sec:ch47-hypothesis-use}

\begin{center}
\begin{tabular}{p{.30\textwidth}p{.62\textwidth}}
\toprule
hypothesis & where it is used\\
\midrule
\(X\) integral projective &
throughout, through \Cref{chap:bj-energy}; integrality is used again in the
vanishing clause of \Cref{prop:bj-energy-I-identities}, where reducibility
would weaken ``constant'' to ``locally constant''\\
\(k\) algebraically closed and trivially valued &
in \Cref{lem:ch47-three-normalizations}, hence in every use of the identity
\(J=J_0\); and in every statement about the scaling action
(\Cref{prop:ch47-scaling}, \Cref{prop:ch47-dual-scaling},
\Cref{cor:ch47-scaling-cy})\\
\(X\) unibranch &
only through \Cref{imp:ch47-envelope}: the operator \(P_\omega\)
(\Cref{sec:ch47-differentiability},
\Cref{sec:ch47-orthogonality}) and the compactness of
\(\PSH^{\sup}(\omega)\) (Step~1 of
\Cref{thm:ch47-existence-admissible}, and (1)\(\Rightarrow\)(4) of
\Cref{thm:bj-cy-envelope-equivalence})\\
\(\omega\) ample &
so that \(V>0\) and the probability normalization is legitimate; so that
the constants are \(\omega\)-psh, used in
\Cref{lem:ch47-mass-forced},
\Cref{prop:ch47-F-elementary}\,\cref{itm:ch47-F:bounds} and
\Cref{cor:ch47-dual-zero}; and for the continuity of envelopes
\Cref{cor:ch45-continuity-envelopes}, on which
\Cref{thm:ch47-orthogonality-continuous} rests
(\Cref{warn:ch45-continuity-ample})\\
finite energy of both arguments &
in every identity involving a difference: \eqref{eq:bj-I-J-potentials},
\eqref{eq:ch47-bernstein-derivative}, \eqref{eq:bj-I-J-comparison};
see \Cref{warn:ch46-no-subtraction}\\
\(\mu\) a probability measure &
forced, \Cref{lem:ch47-mass-forced}\\
\(\mu\) uniformly integrable &
\Cref{lem:ch47-coercivity} and Step~3 of
\Cref{thm:ch47-existence-admissible}\\
\(\mu\) continuously paired &
Step~2 of \Cref{thm:ch47-existence-admissible}, and nowhere else; it is not
needed for \Cref{cor:bj-cy-unibranch}\\
\(h\) continuous &
\Cref{thm:ch47-orthogonality-continuous}; for lower semicontinuous \(h\)
see \Cref{thm:bj-obstacle-orthogonality-energy} and the star in
\eqref{eq:bj-obstacle-contact}\\
\bottomrule
\end{tabular}
\end{center}

\subsection{What is not claimed}
\label{sec:ch47-not-claimed}

\begin{enumerate}
\item No regularity of the solution \(\phi\) of
      \eqref{eq:bj-cy-equation} is asserted beyond membership in
      \(\cE^1(\omega)\).  Continuity holds when \(\mu\) is carried by a
      \(T_\omega\)-ball, by \cite[Theorem~12.12 and
      Proposition~12.14(i)]{BJ22} quoted in
      \Cref{ex:ch47-dirac-green}; it fails in general, and no PL or model
      regularity is available even for \(\mu=\delta_v\) with \(v\)
      divisorial (\cite[Theorem~A]{BJGreen24}).
\item Nothing is asserted about zero-volume or merely pseudoeffective
      classes.  \(E_\omega\), \(\MA_\omega\), \(\cE^1(\omega)\),
      \(E^\vee_\omega\) and \(\cM^1(\omega)\) all divide by \(V\) and are
      used only when \(V>0\).  For a big class the correct substitute
      measures energy relative to \(V_\theta=P_\theta(0)\) of
      \Cref{cor:ch45-extremal}; this is not developed here.
\item The strong topologies are not constructed here.  The quasi-metric
      \(\mathbb I\) is introduced and its axioms verified
      (\Cref{prop:ch47-quasimetric}), but the identification of the induced
      topology, its completeness, and the corresponding statements on the
      measure side are \Cref{chap:bj-strong-topologies}
      (\Cref{rem:ch47-strong-deferred}).
\item Capacity estimates and quasicontinuity are developed in
      \Cref{chap:bj-strong-topologies}.  Here the nonpluripolarity
      deduction in \Cref{rem:ch47-polar-scope} uses the finite-energy
      witness theorem as an explicit import.
\item The domination principle \Cref{thm:ch47-domination} and the
      uniqueness statement underlying
      \Cref{thm:ch47-uniqueness}\,\cref{itm:ch47-uniq:I}\(\Rightarrow\)\cref{itm:ch47-uniq:const}
      are imported, not proved, except for \(n=1\)
      (\Cref{prop:ch47-curve-uniqueness}).  The reason is structural and is
      explained in \Cref{rem:ch47-no-nondegenerate-measure}.
\item No comparison is made with the Chambert--Loir--Ducros
      Monge--Amp\`ere operator; that is
      \Cref{chap:compare-ma}.  Nor with the transcendental theories of
      \Cref{part:complex}; the finite-energy classes there are defined by
      different means and are compared only after their interfaces have
      been fixed.
\end{enumerate}

\section{Chapter dictionary}
\label{sec:ch47-dictionary}

\begingroup
\begin{longtable}{p{.24\textwidth}p{.36\textwidth}p{.30\textwidth}}
\toprule
object & content & where\\
\midrule
\endfirsthead
\toprule
object & content & where\\
\midrule
\endhead
segment expansion &
\(E_\omega((1-t)\phi+t\psi)\) is a Bernstein polynomial of degree
\(\le n+1\) with slopes \(\lambda_k\) and second differences
\(-\langle f,f\rangle_{\mathrm D,k}\) &
\Cref{thm:ch47-segment-expansion}\\
\(J_\psi\), \(I\), \(J\) &
\(I=J_\psi(\phi)+J_\phi(\psi)\);
\(J(\phi)=\sup\phi-E(\phi)=J_0(\phi)\) &
\eqref{eq:bj-I-J-potentials}, \eqref{eq:ch47-J-functional},
\Cref{lem:ch47-three-normalizations}\\
comparison &
\(\frac1{n+1}I\le J_\psi(\phi)\le\frac n{n+1}I\), both constants sharp &
\Cref{prop:bj-energy-I-identities},
\Cref{ex:ch47-toric-cy}\\
uniform concavity &
gap above the chord \(\ge c_nt(1-t)I\) with
\(c_n=\frac{n}{(n+1)(n+2)}\); \(\le t(1-t)I\) &
\Cref{thm:bj-uniform-concavity},
\Cref{cor:ch47-concavity-constants}\\
quasi-metric &
\(\mathbb I=I+|\sup\phi-\sup\psi|\); symmetric, separating,
quasi-triangle &
\eqref{eq:bj-I-normalized}, \Cref{prop:ch47-quasimetric}\\
cone structure &
\(I(t\cdot\phi,t\cdot\psi)=tI(\phi,\psi)\),
\(E^\vee(t_\star\mu)=tE^\vee(\mu)\) &
\Cref{prop:ch47-scaling},
\Cref{prop:ch47-dual-scaling}\\
dual energy &
\(E^\vee(\mu)=\sup_{\cE^1}(E-\int\cdot\,\rd\mu)\); convex, weakly lsc,
weakly compact sublevels &
\Cref{def:bj-finite-energy-measure},
\Cref{prop:ch47-dual-elementary}\\
test class &
\(\cH_{\mathrm{dom}}\), PL, continuous, bounded, \(\cE^1\) all give the
same supremum &
\Cref{prop:ch47-test-class}\\
Fenchel &
\(E^\vee(\MA(\phi))=E(\phi)-\int\phi\MA(\phi)=J_\phi(0)\); equality
characterizes solutions &
\Cref{prop:bj-fenchel-equalities},
\Cref{cor:ch47-dual-zero}\\
polar sets &
\(\mu(\{\phi\le-t\})\le(E^\vee(\mu)+J(\phi))/t\) for
\(\phi\in\cE^1_0\) &
\Cref{prop:ch47-no-polar-mass}\\
equilibrium energy &
\(E^\downarrow(h)=E(P_\omega(h))\); monotone, concave, \(1\)-Lipschitz &
\eqref{eq:bj-equilibrium-energy},
\Cref{prop:ch47-F-elementary}\\
one-sided tangent &
\(E^\downarrow(\phi+\varepsilon f)\le E(\phi)+\varepsilon\int f\MA(\phi)\),
exactly &
\Cref{lem:ch47-one-sided-tangent}\\
differentiability &
matching lower bound with a constant independent of \(\phi\); derivative
\(\int f\MA(\phi)\) &
\Cref{thm:bj-energy-differentiability},
\Cref{cor:ch47-uniformity}\\
orthogonality &
\(\int(h-P_\omega(h))\MA(P_\omega(h))=0\) &
\Cref{thm:ch47-orthogonality-continuous},
\Cref{thm:bj-obstacle-orthogonality-energy},
\Cref{ex:ch47-orthogonality-p1}\\
maximizers &
maximizing nets have \(\MA\to\mu\); a maximizer solves the equation &
\Cref{thm:bj-maximizers-converge},
\Cref{cor:ch47-maximizer-solves}\\
uniqueness &
\(\MA(\phi)=\MA(\psi)\Leftrightarrow I=0\Leftrightarrow\phi-\psi\)
constant &
\Cref{thm:ch47-uniqueness},
\Cref{prop:ch47-curve-uniqueness}\\
coercivity &
\(M(\mu)\le\int T_\omega\rd\mu\);
\(F_\mu\le-J+M(\mu)\) on \(\PSH^{\sup}\) &
\Cref{lem:ch47-T-integrable},
\Cref{lem:ch47-coercivity}\\
existence, direct method &
admissible \(\mu\): maximizer exists, is unique, and solves &
\Cref{thm:ch47-existence-admissible},
\Cref{lem:ch47-finite-divisorial},
\Cref{prop:ch47-skeletal-admissible}\\
Cauchy maximizers &
\(I(\phi,\psi)\le\frac{2}{c_n}(\delta_\phi+\delta_\psi)\) &
\Cref{prop:ch47-maximizing-cauchy}\\
Calabi--Yau &
envelope property \(\Leftrightarrow\) \(\MA\) onto \(\cM^1\)
\(\Leftrightarrow\) bi-Lipschitz \(\Leftrightarrow\) completeness &
\Cref{thm:bj-cy-envelope-equivalence},
\Cref{cor:bj-cy-unibranch}\\
Green functions &
\(\MA(\phi_v)=\delta_v\), \(\phi_v(v)=0\), for \(v\in X^{\lin}\) &
\Cref{cor:ch47-divisorial-green},
\Cref{ex:ch47-dirac-green}\\
\bottomrule
\end{longtable}
\endgroup

Four separations should be carried forward.

\emph{Potential side versus measure side.}  \(\cE^1(\omega)\) is defined by
a pointwise condition on psh functions and \(\cM^1(\omega)\) by an
inequality on measures; the Calabi--Yau theorem says they correspond, and
that correspondence is the content of the chapter, not its definition.  The
different construction in \Cref{chap:bfj-finite-energy-measures}, where the
measure class is an abstract completion, must not be conflated with this
one (\Cref{rem:ch47-bfj-comparison}).

\emph{Weak versus strong.}  Every compactness statement used here is weak
(\Cref{rem:ch47-compactness-attributed}); every convergence statement about
solutions is strong, and the strong theory is
\Cref{chap:bj-strong-topologies}
(\Cref{rem:ch47-strong-deferred}).

\emph{Proved versus imported.}  Four statements are imported and are listed
once, in \Cref{imp:ch47-imported}.  Two of them --- the quasi-triangle
inequality and uniqueness --- are proved here for \(n=1\); one --- the
uniform quadratic lower bound --- is the analytic heart and is not; one ---
completeness --- is deferred to the next chapter and is not used for
admissible data.

\emph{Admissible versus finite dual energy.}  The direct method needs
uniform integrability and continuous pairing;
\(\cM^1(\omega)\) needs neither, and
\Cref{ex:ch47-explicit-potential} exhibits a measure in \(\cM^1(\omega)\)
with \(M(\mu)=+\infty\).  Confusing the two classes is the trivially
valued analogue of confusing \Cref{def:ch28-admissible-measure} with
\(\mathcal M^1_V\) in \Cref{part:bfj}.

\section{Source and provenance}
\label{sec:ch47-source}

\subsection{Source note}
\label{sec:ch47-source-note}

The potential functionals are \cite[Definitions~7.24--7.25 and
Proposition~7.27]{BJ22}.  Uniform concavity is \cite[Theorem~7.28]{BJ22}.
Measure energy, Legendre duality, and maximizing nets are treated in
\cite[\S\S9.1--9.3]{BJ22}, especially
\cite[Proposition~9.6 and Theorem~9.11]{BJ22}.  Differentiability is
\cite[Theorem~8.5 and Corollary~8.6]{BJ22}.  The envelope/Calabi--Yau
equivalence is \cite[Theorem~12.8]{BJ22}.
The newer synthetic measure-side theory cited in
\Cref{chap:bfj-finite-energy-measures} is
\cite{BJ25FiniteEnergy}; it should not be retroactively attributed to BJ22.

Three further attributions belong here.  The comparison
\eqref{eq:bj-I-J-comparison} and the quasi-triangle inequality
\eqref{eq:bj-I-quasitriangle} are \cite[\S7.6 and Theorem~3.31]{BJ22}; the
first is proved in full above, the second is not.  Uniqueness for the
Monge--Amp\`ere equation on \(\cE^1\) is \cite[Corollary~10.4]{BJ22}, and
the domination principle belongs to the same section \cite[\S10]{BJ22}.
Orthogonality for lower semicontinuous obstacles is
\cite[Theorem~13.1]{BJ22}; the continuous case proved here is the first
step of its proof.  Continuity of Green functions and the structure of
solutions with support in a \(T_\omega\)-ball are
\cite[Theorem~12.12 and Proposition~12.14]{BJ22}, and the failure of
\(\mathbf Q\)-PL regularity is \cite[Theorem~A]{BJGreen24}.

\subsection{What is supplied here and is not in the source}
\label{sec:ch47-additions}

The following material is not in \cite{BJ22} in this form and is supplied
here.

\Cref{lem:ch47-three-normalizations} and
\Cref{warn:ch47-J-identity}: the identification
\(\MA_\omega(0)=\delta_{v_{\triv}}\), the maximum principle at
\(v_{\triv}\), and the resulting identity
\(J=\sup\phi-E_\omega=J_0\).  All three ingredients are in
\Cref{chap:bj-fs-pl-psh,chap:bj-energy}; the packaging, and the warning
that the identity fails over a discretely valued base, are supplied here
because they are used on almost every page below.

\Cref{thm:ch47-segment-expansion}: the Bernstein expansion of the energy
along a segment, with the identification of the second differences as
Dirichlet forms.  This is the trivially valued counterpart of
\Cref{prop:ch27-energy-polynomial} and is written out because it makes
\Cref{prop:bj-energy-I-identities} and
\Cref{thm:bj-uniform-concavity} into computations rather than citations.

\Cref{cor:ch47-concavity-constants}: the explicit constant
\(c_n=n/((n+1)(n+2))\) in \eqref{eq:bj-uniform-concavity}, together with
the reverse estimate \eqref{eq:ch47-concavity-upper} and the exact value
\(\tfrac12t(1-t)I\) for \(n=1\).  \cite[Theorem~7.28]{BJ22} asserts the
existence of \(c_n\) without a value; the Green's-function proof given
above produces one, and \Cref{rem:ch47-concavity-shape} explains why
polynomiality is indispensable.

\Cref{prop:ch47-curve-I-norm} and \Cref{prop:ch47-curve-uniqueness}: the
complete treatment of the case \(n=1\), where \(I^{1/2}\) is a genuine
seminorm and uniqueness is elementary.  These make two of the four imported
statements unnecessary in dimension one, and
\Cref{rem:ch47-curve-not-general} identifies exactly what fails in higher
dimension.

\Cref{prop:ch47-test-class}, \Cref{lem:ch47-mass-forced},
\Cref{prop:ch47-dual-elementary},
\Cref{prop:ch47-no-polar-mass}: the elementary structure theory of
\(E^\vee_\omega\), with the Chebyshev estimate.  The statements are
implicit in \cite[\S9]{BJ22}; the proofs are written out because they are
the exact analogues of
\Cref{prop:ch28-test-class,lem:ch28-mass-forced,prop:ch28-dual-order}
and the comparison between the two regimes is one of the stated purposes of
this chapter.

\Cref{prop:ch47-scaling}, \Cref{prop:ch47-dual-scaling},
\Cref{cor:ch47-scaling-cy}: the interaction with the scaling action.  That
\(E_\omega\) is homogeneous of degree one is
\Cref{thm:ch46-energy-homogeneous}; that \(I\), \(J\) and
\(E^\vee_\omega\) inherit the homogeneity, that
\(\delta_{v_{\triv}}\) is the only invariant measure of finite dual
energy, and that the Calabi--Yau correspondence is equivariant, are
supplied here.

\Cref{lem:ch47-one-sided-tangent} and
\Cref{prop:ch47-F-elementary}: the exact one-sided tangent inequality for
the equilibrium energy and the elementary calculus of \(E^\downarrow\).
Isolating the half of \eqref{eq:bj-uniform-differentiability} that is free
is what makes \Cref{thm:ch47-orthogonality-continuous} a three-line
argument and what confines the import to a single inequality.

\Cref{def:ch47-admissible-measure}, \Cref{lem:ch47-T-integrable},
\Cref{lem:ch47-coercivity}, \Cref{lem:ch47-finite-divisorial},
\Cref{prop:ch47-skeletal-admissible},
\Cref{thm:ch47-existence-admissible},
\Cref{cor:ch47-divisorial-green}: the direct-method existence theorem for
admissible data.  \cite{BJ22} proves surjectivity of \(\MA_\omega\) in one
step, through completeness; the version given here is weaker but
self-contained, and it identifies precisely which compactness statement
does the work.  It is modelled on
\Cref{def:ch28-admissible-measure,lem:ch28-coercivity,thm:ch28-existence-admissible}.

\Cref{prop:ch47-maximizing-cauchy}: that almost-maximizers are close, with
an explicit constant.  This is the step that converts uniform concavity
into a Cauchy property and hence into existence and uniqueness; it is
implicit in \cite[\S12]{BJ22} and is isolated here because it is the only
place where \eqref{eq:bj-uniform-concavity} is used, and because it makes
the logical role of completeness transparent.

\Cref{ex:ch47-p1-numbers},
\Cref{ex:ch47-orthogonality-p1},
\Cref{ex:ch47-p1-unbounded},
\Cref{ex:ch47-explicit-potential},
\Cref{ex:ch47-toric-cy},
\Cref{ex:ch47-infinite-energy}: all new.  They exist because the two claims
this chapter makes --- that the correspondence
\(\cE^1\leftrightarrow\cM^1\) is a bijection, and that the hypotheses are
not vacuous --- are exactly the claims a plausible reading can appear to
establish without doing so.  A reader who has solved a Calabi--Yau equation
by two integrations on \(\mathbf P^1\), read one off a polytope on
\(\mathbf P^2\) and on \(\mathbf P^1\times\mathbf P^1\), and watched a
perfectly ordinary probability measure fail to have finite dual energy,
will not mistake \(\cM^1(\omega)\) for \(\mathcal M\).

\subsection{Historical note}
\label{sec:ch47-history}

The variational solution of a Monge--Amp\`ere equation has a short and
well-documented history, and the non-Archimedean chapter of it is the third
of three.

The first is complex.  Aubin and Yau solved the complex Monge--Amp\`ere
equation on a compact Kähler manifold by continuity methods
\cite{Aub84,Y78}, answering Calabi's question \cite{Cal57}; the variational
reformulation, in which the solution is a maximizer of
\(E-\int\cdot\,\rd\mu\) over a finite-energy class, is due to
Berman--Boucksom--Guedj--Zeriahi \cite{BBGZ13}, building on the energy
classes of Guedj--Zeriahi \cite{GZ07} and on the pluripotential theory
surveyed in \cite{GZ17}.  Every structural feature of the present chapter
--- the Legendre duality, the differentiability of the equilibrium energy,
the orthogonality relation, the role of uniform concavity --- appears there
first, and the reader who knows \cite{BBGZ13} will recognize
\Cref{sec:ch47-dual,sec:ch47-differentiability,sec:ch47-orthogonality} as
its non-Archimedean transcription.

The second is the model-theoretic non-Archimedean one.  Boucksom, Favre and
Jonsson transported the variational method to a discretely valued base
\cite{BFJ15,BFJ16}, where the same four ingredients appear but where
orthogonality had to be assumed, or proved in special cases;
\Cref{chap:bfj-differentiability,chap:bfj-finite-energy-measures} give that
theory, and \Cref{rem:ch47-orthogonality-status} records the exact point of
divergence.  It is worth stressing that the difficulty there is not
technical: over a discretely valued base the envelope of a continuous
function need not be continuous, and the segment argument of
\Cref{thm:ch47-orthogonality-continuous} is simply unavailable.

The third is the trivially valued theory of Boucksom and Jonsson
\cite{BJ22}, of which this chapter is an account.  Their \S\S7--13 build
the energy, the duality, the differentiability, the orthogonality and the
Calabi--Yau theorem, and prove them under the hypothesis
\(\operatorname{EP}_X(\omega)\); the hypothesis was verified there for
smooth \(X\) in characteristic zero and for surfaces, whence the shape of
\cite[Theorem~12.9]{BJ22}.  The envelope theorem of
\Crefrange{chap:bj-envelope-statement}{chap:bj-envelope-consequences}
removes the hypothesis for unibranch \(X\), and
\Cref{cor:bj-cy-unibranch} is the resulting unconditional statement.  It is
a corollary of an envelope theorem, and it should be cited as such: the
variational argument is BJ's, the hypothesis-removal is not.

Two threads run out of this chapter.  On the analytic side, the strong
topologies and the capacity theory of
\Cref{chap:bj-strong-topologies} complete the picture and, in particular,
supply the completeness statement imported here.  On the algebraic side,
the identification of \(\cM^1(\omega)\) with a completion of a space of
filtrations and norms, \Cref{chap:bj-stability}, turns the Calabi--Yau
theorem into a statement about K-stability; the metric \(d_1\) of
\Cref{thm:stability-norm-completion} is the norm-theoretic incarnation of
the quasi-metric \(\mathbb I\) of \eqref{eq:bj-I-normalized}, and
\cite{BJKS1} is where the two are compared.  The reader interested in the
Yau--Tian--Donaldson correspondence should read
\Cref{chap:bj-stability} next; the reader interested in the comparison with
complex geometry should read \Cref{part:complex}.
